\documentclass[11pt,reqno]{amsart}
\usepackage[margin=1.1in]{geometry}
\usepackage{amsmath,amssymb,amsthm,mathtools}
\usepackage{enumitem}
\usepackage{booktabs,longtable}
\usepackage{xcolor}
\usepackage[colorlinks=true,linkcolor=blue!50!black,citecolor=green!40!black,urlcolor=blue!50!black]{hyperref}

% ---------------------------------------------------------------------------
% Theorem environments
% ---------------------------------------------------------------------------
\newtheorem{theorem}{Theorem}[section]
\newtheorem{proposition}[theorem]{Proposition}
\newtheorem{lemma}[theorem]{Lemma}
\newtheorem{corollary}[theorem]{Corollary}
\newtheorem{fact}[theorem]{Fact}
\theoremstyle{definition}
\newtheorem{definition}[theorem]{Definition}
\newtheorem{example}[theorem]{Example}
\theoremstyle{remark}
\newtheorem{remark}[theorem]{Remark}

% ---------------------------------------------------------------------------
% Macros
% ---------------------------------------------------------------------------
\newcommand{\R}{\mathbb{R}}
\newcommand{\N}{\mathbb{N}}
\newcommand{\Z}{\mathbb{Z}}
\newcommand{\C}{\mathbb{C}}
\newcommand{\Q}{\mathbb{Q}}
\newcommand{\D}{\mathbb{D}}
\newcommand{\T}{\mathbb{T}}
\newcommand{\Hp}{\mathbb{H}}
\newcommand{\Sph}{\mathbb{S}^2}
\newcommand{\cl}{\operatorname{Cl}}
\newcommand{\inter}{\operatorname{int}}
\newcommand{\diam}{\operatorname{diam}}
\newcommand{\dist}{\operatorname{dist}}
\newcommand{\acc}{\operatorname{acc}}
\newcommand{\CC}{\operatorname{CC}}
\newcommand{\Trap}{\operatorname{Trap}}
\newcommand{\Tri}{\operatorname{Tri}}
\newcommand{\expi}{\operatorname{expi}}
\newcommand{\abs}[1]{\lvert #1\rvert}
\newcommand{\norm}[1]{\lVert #1\rVert}
\newcommand{\eps}{\varepsilon}
\newcommand{\Fs}{F_\sigma}
\newcommand{\Gd}{G_\delta}
\newcommand{\Fsd}{F_{\sigma\delta}}
\newcommand{\todo}[1]{\textcolor{red!70!black}{[#1]}}

% ---------------------------------------------------------------------------
% Blueprint annotations (leanblueprint-compatible). They are invisible in the
% PDF; scripts/blueprint.py reads them to build blueprint/map.md.
%   \lean{Decl1, Decl2}  Lean declaration(s) formalizing this item
%   \uses{label1, ...}   items this statement / proof depends on
%   \source{text}        where an imported Fact comes from
% ---------------------------------------------------------------------------
\newcommand{\lean}[1]{}
\newcommand{\uses}[1]{}
\newcommand{\leanok}{}
\newcommand{\source}[1]{}

\numberwithin{equation}{section}

\title{The mathematical works of T.\,J.~Kaczynski:\\ a unified account and formalization blueprint}
\date{October 2026}

\begin{document}

\begin{abstract}
We give a single, self-contained account, in traditional mathematical language, of every result in
the following seven sources: the doctoral dissertation \emph{Boundary functions} (1967) \cite{Diss};
the papers \cite{JMM} (described only by its review MR0176080), \cite{MMJ}, \cite{TAMSa},
\cite{TAMSb}, \cite{PAMS}; and the note \cite{MathMag} on a problem of Alan Sutcliffe. The one
result known to us only through a review (Theorem~1 of \cite{JMM}, on homeomorphisms of the disk) is
given a complete proof by the plane-topological methods of the other papers. Proofs that the
originals leave to the reader (the ``readily verifiable facts'' of \cite{TAMSb}, the existence of
special functions, the explicit Sutcliffe solutions, the extension theorems for Baire functions) are
written out. Classical results from the literature that the originals cite are isolated as numbered
\emph{Facts}. The document is organised as a blueprint for a Lean~4 formalization: every numbered
item names the Lean declarations that formalize it and the items it depends on.
\end{abstract}

\maketitle
\tableofcontents

% ===========================================================================
\section{Introduction}\label{sec:intro}
% ===========================================================================

\subsection{The sources}
Seven documents are covered.
\begin{itemize}[leftmargin=2em]
\item \cite{MathMag} (1968): an elementary number-theoretic note settling a conjecture of Sutcliffe
  on integers that are multiplied by an integer when their digits are reversed. Section~\ref{sec:arith}.
\item \cite{JMM} (1965), known to us only through its review. Its theorems concern boundary functions of
  homeomorphisms, continuous, Baire and measurable functions in a disk. Theorems~2--6 of \cite{JMM}
  reappear, with proofs, in the dissertation \cite{Diss} (which says so explicitly); Theorem~1 (on
  homeomorphisms) does not, and is proved here in Section~\ref{sec:bf} by the crosscut methods of
  \cite{MMJ,Diss}.
\item \cite{MMJ} (1966): a short proof of McMillan's theorem that the set of curvilinear convergence of
  a continuous function is of type $\Fsd$, and a generalization to cluster sets. Section~\ref{sec:cc}.
\item \cite{Diss} (1967): the dissertation. Its Introduction (Baire functions) is in
  Section~\ref{sec:baire}; Chapter~I in Sections~\ref{sec:cc}, \ref{sec:bf}, \ref{sec:construct};
  Chapter~II in Section~\ref{sec:discont}.
\item \cite{TAMSb} (1969): every $\Fsd$ set is the set of curvilinear convergence of a continuous
  function with a prescribed boundary function (McMillan's problem); appendix on Baire functions into
  the sphere. Sections~\ref{sec:construct} and~\ref{sec:baire}.
\item \cite{TAMSa} (1969): boundary functions of bounded harmonic functions. Section~\ref{sec:harmonic}.
\item \cite{PAMS} (1969): a bounded continuous function in a cube whose set of curvilinear
  convergence is not Borel. Section~\ref{sec:cube}.
\end{itemize}
A concordance between the numbered results of the sources and the items of this document is given in
Appendix~\ref{app:concordance}.

\subsection{How the document is organised}
Several results occur in more than one source in slightly different settings (the disk $\D$, the upper
half-plane $\Hp$, the strip-like region $Q$). We prove each result once, in the most convenient model,
and transfer it to the others by an explicit homeomorphism (Lemma~\ref{lem:transfer}). When two
sources prove the same statement by different arguments, we give the argument that is shorter or more
general, and say so in the concordance.

\emph{Facts.} Results from the literature that the sources use without proof are stated with their
provenance and proved here (Tietze's arc theorem, the crosscut theorem, the Poisson integral, Banach's
theorem, an analytic non-Borel set, Bagemihl's theorem for continuous functions). Two deep theorems that
the sources rely on, the Jordan curve theorem with accessibility and Kierst's theorem on asymptotic
values, are avoided: we replace the arguments that use them by direct ones (Lemmas~\ref{lem:ha-access}
and~\ref{lem:cube-lemma-elementary}), and record the original arguments in remarks.

\emph{Departures from the sources.} We record here the places where our proof differs in substance from
the source.
\begin{enumerate}[label=(\roman*)]
\item The step (B)(III) of \cite{TAMSb} cites a theorem of Bagemihl and McMillan; we use instead the
  approximation theorem of the dissertation (Theorem~\ref{thm:baire1-char}), as \cite{TAMSb} itself
  suggests.
\item The proof of (F)(I) in \cite{TAMSb} cites an invariance-of-domain corollary from \cite{Newman}; we
  write down the inverse maps explicitly (Lemma~\ref{lem:fold}).
\item The extension theorems for Baire-one functions (cited from \cite{Hausdorff} in \cite{Diss} and
  used in Appendix~(b),(c) of \cite{TAMSb}) are proved (Theorem~\ref{thm:extension}).
\item The ``Figure~5'' argument of \cite{Diss} and the monotonicity of first hitting points are derived
  from a single imported crosscut theorem (Theorem~\ref{fact:crosscut}).
\item Sutcliffe's Theorems~2 and~3, used in \cite{MathMag}, are proved with explicit solutions
  (Propositions~\ref{prop:sut-two-composite}, \ref{prop:sut-explicit}).
\item In the existence of special functions (statement (XX) of \cite{TAMSb}) we give a non-recursive
  construction: all disjointness constraints turn out to be independent of the other member's size.
\end{enumerate}

\subsection{Conventions}\label{sec:conventions}
$\N=\{1,2,3,\dots\}$. ``Countable'' includes ``finite''. We identify $\R^2$ with $\C$ and write points as
$z=x+iy=\langle x,y\rangle$. We use
\begin{align*}
\D&=\{z:\abs z<1\}, & C&=\{z:\abs z=1\}, & \Hp&=\{z:\operatorname{Im}z>0\},\\
Q&=\{\langle x,y\rangle:-1\le x\le1,\ 0<y\le1\}, & Q^\circ&=(-1,1)\times(0,1), &
X&=(-1,1)\times\{0\}.
\end{align*}
The real axis is identified with $\R$; in Sections~\ref{sec:construct} we identify $X$ with $(-1,1)$.
For $n\in\N$, $\Hp_n=\{0<y<1/n\}$, $\overline{\Hp}_n=\{0\le y\le 1/n\}$ and $X_n=\{y=1/n\}$.
$\Sph=\{v\in\R^3:\abs v=1\}$ is the Riemann sphere with the Euclidean metric.
$S(r,p)$ is the open ball of radius $r$ about $p$; $\cl E$ or $\overline E$ is closure, $E^*$ or
$\inter E$ interior (for subsets of $C$ or of $\R$, relative to $C$ or $\R$).
For an interval or arc $I$, $I^*$ is its relative interior; an interval or arc is
\emph{nondegenerate} if it has more than one point.

% ===========================================================================
\section{Sutcliffe's problem}\label{sec:arith}
% ===========================================================================
Source: \cite{MathMag}, which uses Theorems~2 and~3 of \cite{Sut}; we prove these too.

\begin{definition}[digit solutions]\label{def:sut-solution}
\lean{TJK.Sutcliffe.IsSolution, TJK.Sutcliffe.HasSolution}
For an integer $n\ge2$ and nonnegative integers $a_h,\dots,a_0$ write
$(a_h,\dots,a_1,a_0)_n=a_hn^h+\dots+a_1n+a_0$. An \emph{$(h+1)$-digit solution for $n$} is a
collection of integers $n,a_h,\dots,a_0,k$ with
\[
  n\ge2,\quad k>1,\quad a_h>0,\quad a_0>0,\quad 0\le a_i\le n-1\ (0\le i\le h),\qquad
  k\,(a_h,\dots,a_1,a_0)_n=(a_0,a_1,\dots,a_h)_n .
\]
We say \emph{there is an $(h+1)$-digit solution for $n$} if such $a_i,k$ exist. Thus a 2-digit
solution is $(a,c,k)$ with $k(an+c)=cn+a$, and a 3-digit solution is $(a,b,c,k)$ with
$k(an^2+bn+c)=cn^2+bn+a$, where $1\le a,c\le n-1$, $0\le b\le n-1$, $k\ge2$.
\end{definition}

\begin{lemma}[basic inequalities]\label{lem:sut-basic}
\lean{TJK.Sutcliffe.two_digit_basic, TJK.Sutcliffe.three_digit_basic}
\uses{def:sut-solution}
For a 2-digit solution $(a,c,k)$ for $n$, and for a 3-digit solution $(a,b,c,k)$ for $n$, we have
$k<n$ and $a<c$.
\end{lemma}
\begin{proof}
In the 3-digit case $an^2+bn+c\ge n^2$ because $a\ge1$, and $cn^2+bn+a\le (n-1)(n^2+n+1)=n^3-1<n^3$;
so $kn^2\le k(an^2+bn+c)<n^3$ and $k<n$. Since $k\ge2$,
$cn^2+bn+a\ge 2(an^2+bn+c)>an^2+bn+c$, i.e.\ $(c-a)(n^2-1)>0$, so $c>a$. The 2-digit case is the
same with $an+c\ge n$, $cn+a\le n^2-1$ and $(c-a)(n-1)>0$.
\end{proof}

\begin{proposition}[Sutcliffe, Theorem 2, first half]\label{prop:sut-two-composite}
\lean{TJK.Sutcliffe.not_prime_of_two_digit}
\uses{def:sut-solution, lem:sut-basic}
If there is a 2-digit solution for $n$, then $n+1$ is not prime.
\end{proposition}
\begin{proof}
Let $k(an+c)=cn+a$. Modulo $p=n+1$ we have $n\equiv-1$, so $k(c-a)\equiv a-c$, i.e.\
$p\mid(k+1)(c-a)$. By Lemma~\ref{lem:sut-basic}, $0<c-a<n<p$ and $3\le k+1\le n<p$. If $p$ were
prime it would divide one of the two factors, which is impossible.
\end{proof}

\begin{proposition}[explicit solutions]\label{prop:sut-explicit}
\lean{TJK.Sutcliffe.explicit_two_digit, TJK.Sutcliffe.explicit_three_digit}
\uses{def:sut-solution}
Let $n\ge2$ and suppose $n+1=de$ with integers $d\ge3$ and $e\ge2$. Put $a=e-1$, $c=n-e$, $k=d-1$.
Then $(a,c,k)$ is a 2-digit solution for $n$ and $(a,n-1,c,k)$ is a 3-digit solution for $n$.
\end{proposition}
\begin{proof}
Digits: $a=e-1\ge1$; $e\le(n+1)/3\le n-1$ gives $c=n-e\ge1$; clearly $a,c\le n-1$; and $k=d-1\ge2$.
Using $n+1=de$:
\begin{align*}
k(an+c)&=(d-1)\bigl((e-1)n+n-e\bigr)=(d-1)e(n-1),\\
cn+a&=(n-e)n+e-1=(n-1)(n+1-e)=(n-1)(de-e)=(d-1)e(n-1),\\
k\bigl(an^2+(n-1)n+c\bigr)&=(d-1)\bigl((e-1)n^2+n^2-n+n-e\bigr)=(d-1)e(n^2-1),\\
cn^2+(n-1)n+a&=(n-e)n^2+n^2-n+e-1=(n^2-1)(n+1-e)=(d-1)e(n^2-1).\qedhere
\end{align*}
\end{proof}

\begin{lemma}[factorization]\label{lem:sut-factor}
\lean{TJK.Sutcliffe.exists_factor}
If $n\ge4$ and $n+1$ is not prime, then $n+1=de$ with $d\ge3$ and $e\ge2$.
\end{lemma}
\begin{proof}
Write $n+1=uv$ with $2\le u\le v$. Then $v\ge3$, since $v=2$ would force $n+1=4$. Take $d=v$, $e=u$.
\end{proof}

\begin{theorem}[Sutcliffe, Theorem 2]\label{thm:sut-sutcliffe2}
\lean{TJK.Sutcliffe.two_digit_iff}
\uses{prop:sut-two-composite, prop:sut-explicit, lem:sut-factor}
Let $n\ge4$. There is a 2-digit solution for $n$ if and only if $n+1$ is not prime.
\end{theorem}
\begin{proof}
One direction is Proposition~\ref{prop:sut-two-composite}; the other is Lemma~\ref{lem:sut-factor} with
Proposition~\ref{prop:sut-explicit}.
\end{proof}

\begin{theorem}[Sutcliffe, Theorem 3]\label{thm:sut-sutcliffe3}
\lean{TJK.Sutcliffe.three_digit_of_two_digit}
\uses{lem:sut-basic, lem:sut-factor, prop:sut-explicit, prop:sut-two-composite}
If there is a 2-digit solution for $n$, then there is a 3-digit solution for $n$.
\end{theorem}
\begin{proof}
By Lemma~\ref{lem:sut-basic}, $1<k<n$, so $n\ge3$. For $n=3$ we would have $k=2$ and $2(3a+c)=3c+a$,
i.e.\ $c=5a\ge5>n-1$; impossible. So $n\ge4$, $n+1$ is not prime by
Proposition~\ref{prop:sut-two-composite}, and Lemma~\ref{lem:sut-factor} with
Proposition~\ref{prop:sut-explicit} gives a 3-digit solution. (In $\N$-valued arithmetic the
expressions $n(a-1)$, $a-b+c$ below are to be read in $\Z$.)
\end{proof}

\begin{lemma}[{\cite[Lemma 1]{MathMag}}]\label{lem:sut-lemma1}
\lean{TJK.Sutcliffe.kaczynski_lemma1}
There are no natural numbers $a,c,n$ with
$n\ge2$, $a<c\le n-1$, $a+c\ge n+1$, and $n(a-1)+c\mid n(c-1)+a$.
\end{lemma}
\begin{proof}
Suppose there are, and write
\begin{equation}\label{eq:sut1}
k\bigl(n(a-1)+c\bigr)=n(c-1)+a,\qquad k\ge1 .
\end{equation}
(Here $a\ge2$ because $a\ge n+1-c\ge2$, so both sides are positive and $k\ge1$.) Modulo $n-1$ we have
$n\equiv1$, so $k(a-n+c)\equiv c-n+a$, i.e.\ $n-1\mid(k-1)(a-n+c)$. Since $a+c\ge n+1$, $a-n+c>0$. If
$k=1$ then $n(a-1)+c=n(c-1)+a$, so $c-a=n(c-a)$ and $n=1$, contrary to $n\ge2$. Hence $k>1$,
$(k-1)(a-n+c)>0$, and so $(k-1)(a+c-n)\ge n-1$. Since $c-n\le-1$, $(k-1)(a-1)\ge(k-1)(a+c-n)\ge n-1$.
Therefore
\[
k\bigl(n(a-1)+c\bigr)=nk(a-1)+kc\ge n(k-1)(a-1)+kc\ge n(n-1)+kc\ge nc+kc>n(c-1)+a,
\]
using $c\le n-1$ and $kc>0\ge a-n$. This contradicts \eqref{eq:sut1}.
\end{proof}

\begin{lemma}[{\cite[Lemma 2]{MathMag}}]\label{lem:sut-lemma2}
\lean{TJK.Sutcliffe.kaczynski_lemma2}
\uses{lem:sut-basic}
If $n=2$ or $n=3$, there is no 3-digit solution for $n$.
\end{lemma}
\begin{proof}
By Lemma~\ref{lem:sut-basic}, $1<k<n$, which is impossible for $n=2$. For $n=3$ we get $k=2$ and
$2(9a+3b+c)=9c+3b+a$, i.e.\ $0=17a+3b-7c\ge17+3b-14>0$, a contradiction.
\end{proof}

\begin{theorem}[Kaczynski {\cite{MathMag}}; Sutcliffe's conjecture]\label{thm:sut-main}
\lean{TJK.Sutcliffe.two_digit_of_three_digit}
\uses{lem:sut-lemma1, lem:sut-lemma2, lem:sut-basic, thm:sut-sutcliffe2}
Let $n\ge2$. If there is a 3-digit solution for $n$, then there is a 2-digit solution for $n$.
\end{theorem}
\begin{proof}
By Lemma~\ref{lem:sut-lemma2} we may assume $n>3$, and by Theorem~\ref{thm:sut-sutcliffe2} we may
assume that $p=n+1$ is prime. Let $(a,b,c,k)$ be a 3-digit solution:
\begin{equation}\label{eq:sut3}
k(an^2+bn+c)=cn^2+bn+a .
\end{equation}
Modulo $p$, $n\equiv-1$, so $k(a-b+c)\equiv c-b+a$ and $p\mid(k-1)(a-b+c)$. If $p\nmid a-b+c$ then
$p\mid k-1$; since $k-1>0$ this gives $k-1\ge p$, impossible because $k<n$
(Lemma~\ref{lem:sut-basic}). Hence $p\mid a-b+c$. Now
$-p<-n<3-n\le a-b+c\le 2n-2<2p$, so $a-b+c=ep$ with $e\in\{0,1\}$. Substituting $b=a+c-ep$,
\[
an^2+bn+c=a(n^2+n)+c(n+1)-epn=p\bigl((a-e)n+c\bigr),\qquad cn^2+bn+a=p\bigl((c-e)n+a\bigr),
\]
so \eqref{eq:sut3} becomes $k\bigl((a-e)n+c\bigr)=(c-e)n+a$. If $e=1$ then
$k\bigl(n(a-1)+c\bigr)=n(c-1)+a$; moreover $0\le b=a+c-p$ gives $a+c\ge n+1$, and $a<c\le n-1$ by
Lemma~\ref{lem:sut-basic}. This is impossible by Lemma~\ref{lem:sut-lemma1}. Hence $e=0$ and
$k(an+c)=cn+a$ with $1\le a,c\le n-1$ and $k>1$: a 2-digit solution.
\end{proof}

\begin{corollary}[{\cite{MathMag}}]\label{cor:sut-main}
\lean{TJK.Sutcliffe.three_digit_iff}
\uses{thm:sut-main, thm:sut-sutcliffe2, thm:sut-sutcliffe3}
If $n>3$, there is a 3-digit solution for $n$ if and only if $n+1$ is not prime.
\end{corollary}
\begin{proof}
If $n+1$ is not prime, Theorem~\ref{thm:sut-sutcliffe2} gives a 2-digit solution and
Theorem~\ref{thm:sut-sutcliffe3} a 3-digit one. Conversely a 3-digit solution gives a 2-digit solution
(Theorem~\ref{thm:sut-main}), and then $n+1$ is not prime (Theorem~\ref{thm:sut-sutcliffe2}).
\end{proof}

\begin{remark}
\cite{MathMag} notes that a similar result was obtained independently by Prasert Na Nagara.
\end{remark}

% ===========================================================================
\section{Borel classes and Baire functions}\label{sec:baire}
% ===========================================================================
Sources: \cite[Introduction, \S3]{Diss}, \cite[Appendix]{TAMSb}. Throughout, $M$ is a metric space.

\subsection{Classes of sets}

\begin{definition}[$\Fs$, $\Gd$, $\Fsd$, ambiguous sets]\label{def:borel-classes}
\lean{TJK.IsFσ, TJK.IsFσδ, TJK.IsAmbiguous}
A subset of a topological space is $\Fs$ if it is a countable union of closed sets, $\Gd$ if it is a
countable intersection of open sets, $\Fsd$ if it is a countable intersection of $\Fs$ sets, and
\emph{ambiguous} if it is both $\Fs$ and $\Gd$. For $L\subseteq M$ we write $\Fs(L)$, $\Gd(L)$,
$\Fsd(L)$ for the corresponding classes of the subspace $L$.
\end{definition}

\begin{lemma}[closure properties]\label{lem:borel-basic}
\lean{TJK.IsFσ.iUnion, TJK.IsFσ.inter, TJK.IsFσδ.iInter, TJK.IsFσδ.union, TJK.IsOpen.isFσ,
  TJK.Set.Countable.isFσ, TJK.isFσ_subtype_iff, TJK.isGδ_subtype_iff, TJK.isFσδ_subtype_iff,
  TJK.IsFσδ.of_subtype}
\uses{def:borel-classes}
\begin{enumerate}[label=(\alph*)]
\item In $M$, open sets are $\Fs$, closed sets are $\Gd$, and countable sets are $\Fs$.
\item $\Fs$ is closed under countable unions and finite intersections; $\Gd$ under countable
  intersections and finite unions; $\Fsd$ under countable intersections and finite unions. The
  complement of an $\Fs$ set is $\Gd$ and conversely.
\item For $L\subseteq M$: $E\in\Fs(L)$ iff $E=P\cap L$ for some $\Fs$ set $P$ of $M$; likewise for
  $\Gd$ and $\Fsd$.
\item If $L$ is $\Fsd$ (resp.\ $\Gd$) in $M$ and $E\in\Fsd(L)$ (resp.\ $\Gd(L)$), then $E$ is $\Fsd$
  (resp.\ $\Gd$) in $M$. Every $\Fs$ set is $\Fsd$, and every $\Gd$ set is $\Fsd$.
\item Homeomorphisms preserve all these classes.
\end{enumerate}
\end{lemma}
\begin{proof}
(a) $U=\bigcup_k\{x:\dist(x,M\setminus U)\ge1/k\}$ for $U\ne M$ open; complements give the rest, and
points are closed. (b) $(\bigcup_nA_n)\cap(\bigcup_mB_m)=\bigcup_{n,m}A_n\cap B_m$, and
$(\bigcap_nP_n)\cup(\bigcap_mR_m)=\bigcap_{n,m}(P_n\cup R_m)$; the rest is De Morgan. (c) Closed sets
of $L$ are traces of closed sets of $M$, and traces commute with countable unions and intersections.
(d) $E=P\cap L$ with $P$ in the class, and the class is closed under finite intersections; a $\Gd$ set
$\bigcap_nU_n$ is $\Fsd$ by (a). (e) is clear.
\end{proof}

\subsection{Classes of functions}

\begin{definition}[function classes]\label{def:fn-classes}
\lean{TJK.IsFσMeasurable, TJK.IsBaireOne, TJK.IsHonoraryBaireTwo, TJK.BaireClass}
\uses{def:borel-classes}
Let $L$ be a metric space and $Y$ a metric space.
\begin{enumerate}[label=(\alph*)]
\item $f:L\to Y$ is \emph{of class $\Fs$} if $f^{-1}(U)\in\Fs(L)$ for every open $U\subseteq Y$.
\item $f$ is \emph{of Baire class $\le1$}, written $f\in\mathcal B_1(L,Y)$, if there are continuous
  $f_n:L\to Y$ with $f_n\to f$ pointwise on $L$.
\item $f$ is \emph{of honorary Baire class $\le2$} if there are a countable $N\subseteq L$ and
  $\psi\in\mathcal B_1(L,Y)$ with $f=\psi$ on $L\setminus N$.
\item (Baire classes of all orders.) $f$ is of Baire class $0$ if it is continuous; for an ordinal
  $\xi\ge1$, $f$ is of Baire class $\xi$ if it is the pointwise limit of a sequence $f_n$, each of
  Baire class $\eta_n$ for some $\eta_n<\xi$. (So class $\xi$ contains all classes $\eta<\xi$ when
  $\xi\ge1$.)
\end{enumerate}
\end{definition}

\begin{lemma}[reduction; {\cite[Lemma 1]{Diss}}]\label{lem:reduction}
\lean{TJK.reduction}
\uses{lem:borel-basic}
Let $E,F$ be $\Fs$ subsets of $M$. There are disjoint $\Fs$ sets $A\subseteq E$, $B\subseteq F$ with
$A\cup B=E\cup F$. In particular $E\setminus F\subseteq A$ and $F\setminus E\subseteq B$.
\end{lemma}
\begin{proof}
Write $E=\bigcup_nE_n$, $F=\bigcup_nF_n$ with $E_n,F_n$ closed, and put
\[
A=\bigcup_n\Bigl(E_n\setminus\bigcup_{j<n}F_j\Bigr),\qquad
B=\bigcup_n\Bigl(F_n\setminus\bigcup_{j\le n}E_j\Bigr).
\]
Each term is closed minus closed, i.e.\ closed $\cap$ open, hence $\Fs$ (Lemma~\ref{lem:borel-basic}).
If $x\in(E_n\setminus\bigcup_{j<n}F_j)\cap(F_m\setminus\bigcup_{j\le m}E_j)$ then $m<n$ is impossible
because $x\in F_m$, and $m\ge n$ is impossible because $x\in E_n$; so $A\cap B=\emptyset$. If
$x\in E\cup F$, let $n$ be least with $x\in E_n\cup F_n$. If $x\in E_n$ then $x\in A$; otherwise
$x\in F_n\setminus\bigcup_{j\le n}E_j\subseteq B$.
\end{proof}

\begin{corollary}[finite reduction]\label{cor:reduction-finite}
\lean{TJK.reduction_finite}
\uses{lem:reduction}
If $P_1,\dots,P_q$ are $\Fs$ subsets of $M$, there are pairwise disjoint $\Fs$ sets $R_i\subseteq P_i$
with $\bigcup_iR_i=\bigcup_iP_i$. If $\bigcup_iP_i=M$, each $R_i$ is ambiguous.
\end{corollary}
\begin{proof}
Induction on $q$. Reduce $P_1$ against $P'=P_2\cup\dots\cup P_q$ (Lemma~\ref{lem:reduction}), getting
disjoint $\Fs$ sets $A\subseteq P_1$, $B\subseteq P'$ with $A\cup B=P_1\cup P'$; apply the induction
hypothesis to $P_2,\dots,P_q$, getting $R_2',\dots,R_q'$, and put $R_1=A$, $R_i=R_i'\cap B$. If the $R_i$
cover $M$, the complement of each is the finite union of the others, hence $\Fs$.
\end{proof}

\begin{theorem}[Sierpi\'nski {\cite{Sierpinski}}]\label{thm:sierpinski}
\lean{TJK.sierpinski_ambiguous}
\uses{lem:reduction, lem:borel-basic}
If $E\subseteq F\subseteq M$ with $E$ a $\Gd$ set and $F$ an $\Fs$ set, there is an ambiguous set
$K$ with $E\subseteq K\subseteq F$.
\end{theorem}
\begin{proof}
$P=M\setminus E$ and $F$ are $\Fs$ with $P\cup F=M$. Reduce them (Lemma~\ref{lem:reduction}):
$A\subseteq P$, $B\subseteq F$ disjoint $\Fs$ with $A\cup B=M$. Put $K=B$. Then $K\subseteq F$,
$K=M\setminus A\supseteq M\setminus P=E$, and $K$ is $\Fs$ with $\Fs$ complement $A$.
\end{proof}

\begin{lemma}\label{lem:baire1-fsigma}
\lean{TJK.IsBaireOne.isFσMeasurable}
\uses{def:fn-classes, lem:borel-basic}
Every $f\in\mathcal B_1(L,Y)$ is of class $\Fs$.
\end{lemma}
\begin{proof}
Let $f_n\to f$ with $f_n$ continuous, and $U\subsetneq Y$ open. Put
$V_k=\{y:\dist(y,Y\setminus U)\ge1/k\}$, a closed subset of $U$. Then
$f^{-1}(U)=\bigcup_{k,N}\bigcap_{n\ge N}f_n^{-1}(V_k)$: if $\dist(f(x),Y\setminus U)>2/k$ then
eventually $f_n(x)\in V_k$; conversely if $f_n(x)\in V_k$ for $n\ge N$ then $f(x)\in V_k\subseteq U$.
\end{proof}

\begin{lemma}[{Banach's Hilfssatz 3}]\label{lem:finite-valued}
\lean{TJK.finiteValued_ambiguous}
\uses{def:fn-classes, lem:borel-basic}
If $f:L\to Y$ is of class $\Fs$ and takes only the values $a_1,\dots,a_q$ (distinct), then the sets
$E_j=f^{-1}(a_j)$ are ambiguous in $L$ and partition $L$.
\end{lemma}
\begin{proof}
With $r$ less than half the least distance between the $a_j$, $E_j=f^{-1}(S(r,a_j))$ is $\Fs$, and its
complement $\bigcup_{i\ne j}E_i$ is $\Fs$ too.
\end{proof}

\begin{lemma}\label{lem:finite-valued-ops}
\lean{TJK.FiniteValued.add, TJK.FiniteValued.sub}
\uses{lem:finite-valued}
If $f,g:L\to\R$ are finitely valued and of class $\Fs$, so are $f+g$ and $f-g$.
\end{lemma}
\begin{proof}
By Lemma~\ref{lem:finite-valued} the level sets of $f$ and $g$ are ambiguous; every preimage under $f\pm g$
is a finite union of sets $f^{-1}(a)\cap g^{-1}(b)$, which are $\Fs$.
\end{proof}

\begin{lemma}[{\cite[Lemma 2]{Diss}}]\label{lem:approx-char}
\lean{TJK.approx_indicator}
\uses{lem:reduction, lem:borel-basic}
Let $L\subseteq M$ and let $E$ be ambiguous in $L$. There are continuous $f_n:M\to[0,1]$ with
$f_n\to\chi_E$ pointwise on $L$.
\end{lemma}
\begin{proof}
Write $E=E_1\cap L$ and $L\setminus E=F_1\cap L$ with $E_1,F_1$ $\Fs$ in $M$. Reduce: disjoint $\Fs$
sets $A\subseteq E_1$, $B\subseteq F_1$ with $A\cup B=E_1\cup F_1\supseteq L$. Then $E=A\cap L$ and
$L\setminus E=B\cap L$ (a point of $E$ is not in $F_1$, hence not in $B$, hence in $A$). Write
$A=\bigcup A_n$, $B=\bigcup B_n$ with $A_n\subseteq A_{n+1}$, $B_n\subseteq B_{n+1}$ closed, and let
$f_n$ be the Urysohn function $\dist(x,B_n)/(\dist(x,A_n)+\dist(x,B_n))$ (if $A_n=\emptyset$ take
$f_n=0$; if $B_n=\emptyset$ take $f_n=1$), which is $1$ on $A_n$ and $0$ on $B_n$. For $x\in L$, eventually $x\in A_n$ or eventually
$x\in B_n$.
\end{proof}

\begin{lemma}[{\cite[Lemma 3]{Diss}}]\label{lem:approx-finite}
\lean{TJK.approx_finiteValued}
\uses{lem:finite-valued, lem:approx-char}
If $f:L\to\R$ is of class $\Fs$ and finitely valued, and $L\subseteq M$, there are continuous
$f_n:M\to\R$ with $f_n\to f$ pointwise on $L$.
\end{lemma}
\begin{proof}
$f=\sum_ja_j\chi_{E_j}$ with $E_j$ ambiguous (Lemma~\ref{lem:finite-valued}); take
$f_n=\sum_ja_jf^{(j)}_n$ with $f^{(j)}_n\to\chi_{E_j}$ from Lemma~\ref{lem:approx-char}.
\end{proof}

\begin{lemma}[{\cite[Lemma 4]{Diss}}, Banach's Hilfssatz 4]\label{lem:uniform-approx}
\lean{TJK.uniform_approx_finiteValued}
\uses{cor:reduction-finite}
If $f:L\to\R$ is bounded and of class $\Fs$, then for every $k\in\N$ there is a finitely valued
$g:L\to\R$ of class $\Fs$ with $\abs{g-f}<1/k$ and $\sup\abs g\le\sup\abs f+1/k$.
\end{lemma}
\begin{proof}
Cover $[-\sup\abs f,\sup\abs f]$ by the finitely many intervals $U_j=((j-1)/k,(j+1)/k)$ that meet it.
The sets $f^{-1}(U_j)$ are $\Fs$ in $L$ and cover $L$; Corollary~\ref{cor:reduction-finite} gives a
partition of $L$ into ambiguous $R_j\subseteq f^{-1}(U_j)$. Put $g=j/k$ on $R_j$. Then $\abs{g-f}<1/k$,
and $g^{-1}(W)$ is a finite union of $R_j$'s for every $W$.
\end{proof}

\begin{lemma}[{\cite[Lemma 5]{Diss}}]\label{lem:approx-real}
\lean{TJK.approx_real}
\uses{lem:uniform-approx, lem:approx-finite, lem:finite-valued-ops}
Let $L\subseteq M$ and let $f:L\to\R$ be of class $\Fs$. There are continuous $f_j:M\to\R$ with
$f_j\to f$ pointwise on $L$.
\end{lemma}
\begin{proof}
\emph{Bounded $f$.} Choose by Lemma~\ref{lem:uniform-approx} finitely valued $g_n$ of class $\Fs$ with
$\norm{g_n-f}_\infty\le2^{-n}$; put $h_1=g_1$, $h_n=g_n-g_{n-1}$. Each $h_n$ is finitely valued of class
$\Fs$, and $\norm{h_n}\le 2^{-n+2}$ for $n\ge2$. By Lemma~\ref{lem:approx-finite} choose continuous
$h_{n,j}\to h_n$ on $L$, and let $k_{n,j}$ be $h_{n,j}$ clamped to $[-\norm{h_n},\norm{h_n}]$; still
$k_{n,j}\to h_n$ on $L$. Then $f_j=\sum_nk_{n,j}$ converges uniformly, so it is continuous. For
$x\in L$ and $m\in\N$,
\[
\abs{f_j(x)-f(x)}\le\sum_{n>m}\norm{h_n}+\sum_{n\le m}\abs{k_{n,j}(x)-h_n(x)}+\abs{g_m(x)-f(x)},
\]
which is small for $m$ large and then $j$ large.

\emph{General $f$.} $\arctan\circ f$ is bounded of class $\Fs$; take continuous $G_j\to\arctan\circ f$
on $L$, clamp $G_j$ into $[-\pi/2+1/j,\pi/2-1/j]$ to get $H_j$ (still convergent on $L$, since the target
values lie in the open interval), and put $f_j=\tan\circ H_j$.
\end{proof}

\begin{theorem}[{\cite[Lemma 6]{Diss}}]\label{thm:baire1-char}
\lean{TJK.baireOne_tfae}
\uses{lem:approx-real, lem:baire1-fsigma}
Let $L\subseteq M$ and $f:L\to\R^m$. The following are equivalent:
(i) $f\in\mathcal B_1(L,\R^m)$; (ii) $f$ is of class $\Fs$; (iii) there are continuous
$f_n:M\to\R^m$ with $f_n\to f$ pointwise on $L$.
\end{theorem}
\begin{proof}
(i)$\Rightarrow$(ii) is Lemma~\ref{lem:baire1-fsigma}; (iii)$\Rightarrow$(i) by restriction. (ii)$\Rightarrow$(iii):
each coordinate $f_i=\pi_i\circ f$ is of class $\Fs$; apply Lemma~\ref{lem:approx-real} to each. (Conversely
the coordinates being of class $\Fs$ implies $f$ is: an open set of $\R^m$ is a countable union of open
boxes, and $f^{-1}(\prod_iI_i)=\bigcap_if_i^{-1}(I_i)$.)
\end{proof}

\begin{lemma}[uniform limits]\label{lem:uniform-limit-fsigma}
\lean{TJK.IsFσMeasurable.of_tendstoUniformly}
\uses{def:fn-classes, lem:borel-basic}
If $S_n:L\to Y$ are of class $\Fs$ and $S_n\to F$ uniformly, then $F$ is of class $\Fs$.
\end{lemma}
\begin{proof}
Let $U\subsetneq Y$ be open and $U_k=\{y:\dist(y,Y\setminus U)>2/k\}$; for each $k$ choose $n_k$ with
$\dist(S_{n_k}(x),F(x))<1/k$ for all $x$. Then $F^{-1}(U)=\bigcup_kS_{n_k}^{-1}(U_k)$.
\end{proof}

\begin{theorem}[extension of Baire-one functions]\label{thm:extension}
\lean{TJK.IsBaireOne.extend_of_isGδ, TJK.IsBaireOne.extend_of_countable_compl}
\uses{lem:finite-valued, lem:finite-valued-ops, cor:reduction-finite, lem:uniform-approx, lem:uniform-limit-fsigma,
  thm:baire1-char, lem:baire1-fsigma, lem:borel-basic, lem:approx-real}
Let $f:L\to\R^m$ be of Baire class $\le1$, $L\subseteq M$.
\begin{enumerate}[label=(\alph*)]
\item If $L$ is $\Gd$ in $M$ and $f$ is bounded, $f$ extends to a bounded $F\in\mathcal B_1(M,\R^m)$.
\item If $M\setminus L$ is countable, $f$ extends to some $F\in\mathcal B_1(M,\R^m)$.
\end{enumerate}
\end{theorem}
\begin{proof}
It suffices to treat $m=1$ (extend coordinatewise).

(a) \emph{Finitely valued $f$ of class $\Fs$ on a $\Gd$ set $L$.} By Lemma~\ref{lem:finite-valued}, $f=\sum_ja_j\chi_{E_j}$ with
$E_1,\dots,E_q$ an ambiguous partition of $L$; write $E_j=P_j\cap L$ with $P_j$ $\Fs$ in $M$. The sets
$P_1,\dots,P_q,M\setminus L$ are $\Fs$ and cover $M$; let $R_1,\dots,R_q,R_0$ be the disjoint $\Fs$
cover given by Corollary~\ref{cor:reduction-finite}. Put $F=a_j$ on $R_j$ and $F=0$ on $R_0$. For
$x\in L\cap E_j$, $x\notin R_0$, and $x\in R_i$ forces $x\in P_i\cap L=E_i$, so $i=j$ and $F(x)=f(x)$.
$F$ is finitely valued of class $\Fs$ and $\sup\abs F\le\max\abs{a_j}$.

\emph{Bounded $f$.} $f$ is of class $\Fs$ (Lemma~\ref{lem:baire1-fsigma}). With $g_n$, $h_n$ as in the proof
of Lemma~\ref{lem:approx-real}, extend each
$h_n$ to a finitely valued $H_n$ of class $\Fs$ with $\sup\abs{H_n}\le\sup\abs{h_n}$. The partial sums
$S_N=\sum_{n\le N}H_n$ are finitely valued of class $\Fs$ and converge uniformly to some $F$, which is
bounded, of class $\Fs$ (Lemma~\ref{lem:uniform-limit-fsigma}), hence in $\mathcal B_1(M,\R)$
(Theorem~\ref{thm:baire1-char} with $L=M$); and $F|_L=\lim g_N=f$.

(b) Let $\tau(v)=v/(1+\abs v)$, a homeomorphism of $\R^m$ onto the open unit ball $B$. Since $M\setminus L$
is countable, $L$ is $\Gd$; by (a) extend $\tau\circ f$ to a bounded $G\in\mathcal B_1(M,\R^m)$. The
set $Z=G^{-1}(\R^m\setminus B)$ lies in $M\setminus L$, so it is countable. Put $G'=G$ off $Z$ and
$G'=0$ on $Z$. For open $U\subseteq B$, $G'^{-1}(U)=G^{-1}(U)\cup(Z\text{ or }\emptyset)$ is $\Fs$. So
$G':M\to B$ is of class $\Fs$, hence so is $\tau^{-1}\circ G':M\to\R^m$, which is therefore in
$\mathcal B_1$ (Theorem~\ref{thm:baire1-char}) and extends $f$.
\end{proof}

\subsection{Functions into the sphere}
We use the folding device of \cite[(F)(I)]{TAMSb}. Write points of $\Sph$ as $\langle x,y,z\rangle$ and
let $T(z)=4z^3-3z$.

\begin{definition}[the fold]\label{def:fold}
\lean{TJK.Fold.fold, TJK.Fold.U, TJK.Fold.V, TJK.Fold.ZU, TJK.Fold.ZV}
Let $s=2^{-1/2}$ and
\begin{gather*}
U=\{-s<z\le1\},\quad V=\{-1\le z<s\},\quad Z_U=\{s<z\le1\},\quad Z_V=\{-1\le z<-s\},\\
\Phi\langle x,y,z\rangle=\begin{cases}\langle x(4z^2-1),\,y(4z^2-1),\,T(z)\rangle,&\abs z\ge s,\\
\langle x,y,-z\rangle,&\abs z\le s,\end{cases}
\end{gather*}
all subsets of $\Sph$. Let $p_1=\langle1,0,0\rangle$ and let $\sigma:\Sph\setminus\{p_1\}\to\R^2$ be
stereographic projection from $p_1$.
\end{definition}

\begin{lemma}[properties of the fold]\label{lem:fold}
\lean{TJK.Fold.continuous_fold, TJK.Fold.fold_mapsTo, TJK.Fold.homeomorph_ZU, TJK.Fold.homeomorph_ZV,
  TJK.Fold.fibre_card_le_three, TJK.Fold.isOpen_U, TJK.Fold.isOpen_V, TJK.Fold.U_union_V,
  TJK.Fold.bounded_chart}
\uses{def:fold}
(i) $\Phi$ is a continuous map $\Sph\to\Sph$. (ii) $\Phi$ restricts to homeomorphisms $Z_U\to U$ and
$Z_V\to V$. (iii) Each fibre of $\Phi$ has at most three points. (iv) $U,V$ are open and
$U\cup V=\Sph$. (v) $Z_U\cup Z_V\subseteq\Sph\setminus\{p_1\}$ and $\sigma(Z_U\cup Z_V)$ is bounded.
\end{lemma}
\begin{proof}
On $\Sph$, $(x^2+y^2)(4z^2-1)^2+T(z)^2=(1-z^2)(16z^4-8z^2+1)+z^2(16z^4-24z^2+9)=1$, so $\Phi$ maps into
$\Sph$. At $z=\pm s$ we have $4z^2-1=1$ and $T(\pm s)=\mp s$, so the two formulas agree and $\Phi$ is
continuous by the pasting lemma. $T'(z)=12z^2-3>0$ for $\abs z>1/2$, so $T$ is a strictly increasing
bijection $[s,1]\to[-s,1]$ and $[-1,-s]\to[-1,s]$. Given $\langle x',y',z'\rangle\in U$, put
$z=T^{-1}(z')\in(s,1]$, $c=4z^2-1\in(1,3]$, $x=x'/c$, $y=y'/c$; the identity above shows
$\langle x,y,z\rangle\in\Sph$, and this is a continuous inverse of $\Phi|_{Z_U}$. Likewise for $Z_V$.
$\Phi$ is injective on each of $Z_U$, $\{\abs z\le s\}$, $Z_V$, which cover $\Sph$; this gives (iii). (iv)
is clear. (v): $\cl(Z_U\cup Z_V)=\{\abs z\ge s\}$ is compact and misses $p_1$.
\end{proof}

\begin{lemma}[fold factorisation]\label{lem:fold-factor}
\lean{TJK.Fold.exists_factor}
\uses{lem:fold, thm:sierpinski, lem:borel-basic}
Let $f:M\to\Sph$ be of class $\Fs$. There is $\psi:M\to Z_U\cup Z_V$ with $\Phi\circ\psi=f$ such that
$\sigma\circ\psi:M\to\R^2$ is bounded and of class $\Fs$.
\end{lemma}
\begin{proof}
$f^{-1}(U)$ is $\Fs$ and $M\setminus f^{-1}(V)$ is $\Gd$, and $M\setminus f^{-1}(V)\subseteq f^{-1}(U)$ since $U\cup V=\Sph$.
Theorem~\ref{thm:sierpinski} gives an ambiguous $K$ with $M\setminus f^{-1}(V)\subseteq K\subseteq f^{-1}(U)$; put
$L=M\setminus K$, so $f(K)\subseteq U$, $f(L)\subseteq V$. Let $\Phi_U=\Phi|_{Z_U}:Z_U\to U$ and $\Phi_V=\Phi|_{Z_V}:Z_V\to V$
(homeomorphisms by Lemma~\ref{lem:fold}), and define $\psi=\Phi_U^{-1}\circ f$ on $K$, $\psi=\Phi_V^{-1}\circ f$ on $L$. For
$G\subseteq\Sph\setminus\{p_1\}$ open, $\psi^{-1}(G)=[K\cap f^{-1}(\Phi_U(Z_U\cap G))]\cup[L\cap f^{-1}(\Phi_V(Z_V\cap G))]$, and
$\Phi_U(Z_U\cap G)$ is open in $U$, hence in $\Sph$; so $\psi^{-1}(G)$ is $\Fs$. Boundedness is Lemma~\ref{lem:fold}(v).
\end{proof}

\begin{theorem}[{\cite[Appendix, Theorem (a)]{TAMSb}}; cf.\ {\cite[Theorem 1]{Diss}}]\label{thm:sphere-baire1}
\lean{TJK.sphere_baireOne_iff}
\uses{lem:fold-factor, thm:baire1-char, lem:baire1-fsigma}
A function $f:M\to\Sph$ is in $\mathcal B_1(M,\Sph)$ if and only if it is of class $\Fs$.
\end{theorem}
\begin{proof}
One direction is Lemma~\ref{lem:baire1-fsigma}. Conversely, take $\psi$ from Lemma~\ref{lem:fold-factor};
by Theorem~\ref{thm:baire1-char} there are continuous
$\psi_n:M\to\R^2$ with $\psi_n\to\sigma\circ\psi$. Then $\Phi\circ\sigma^{-1}\circ\psi_n$ are
continuous and converge to $\Phi\circ\psi=f$.
\end{proof}

\begin{theorem}[{\cite[Appendix, Theorem (c)]{TAMSb}}]\label{thm:sphere-extension}
\lean{TJK.sphere_baireOne_extend}
\uses{lem:fold-factor, thm:extension, thm:baire1-char, lem:baire1-fsigma}
If $E$ is a $\Gd$ subset of $M$ and $f\in\mathcal B_1(E,\Sph)$, then $f$ extends to some
$F\in\mathcal B_1(M,\Sph)$.
\end{theorem}
\begin{proof}
By Lemmas~\ref{lem:baire1-fsigma} and~\ref{lem:fold-factor}, applied in the metric space $E$, write
$f=\Phi\circ\psi$ with $\sigma\circ\psi:E\to\R^2$ bounded of class $\Fs$, hence in $\mathcal B_1$. By
Theorem~\ref{thm:extension}(a) it extends to a bounded $G\in\mathcal B_1(M,\R^2)$; let continuous
$G_n\to G$ on $M$. Then $\Phi\circ\sigma^{-1}\circ G_n$ are continuous and converge on $M$ to
$F=\Phi\circ\sigma^{-1}\circ G$, and $F|_E=\Phi\circ\psi=f$.
\end{proof}

\begin{theorem}[honorary Baire class 2; {\cite[Theorem 2]{Diss}}, {\cite[Appendix, Theorem (b)]{TAMSb}}]\label{thm:honorary-char}
\lean{TJK.isHonoraryBaireTwo_iff_euclidean, TJK.isHonoraryBaireTwo_iff_sphere,
  TJK.isHonoraryBaireTwo_iff_closed_preimage}
\uses{thm:baire1-char, thm:sphere-baire1, thm:extension, thm:sphere-extension, lem:borel-basic, lem:baire1-fsigma}
Let $Y$ be $\R^m$ or $\Sph$ and $f:M\to Y$. The following are equivalent.
\begin{enumerate}[label=(\arabic*)]
\item $f$ is of honorary Baire class $\le2$.
\item There is a countable $N\subseteq M$ such that $f|_{M\setminus N}$ is of class $\Fs$.
\item There is a countable $N\subseteq M$ such that $f^{-1}(F)\setminus N$ is $\Gd$ in $M$ for every
  closed $F\subseteq Y$.
\end{enumerate}
\end{theorem}
\begin{proof}
(1)$\Rightarrow$(2): restrict the Baire-one function and use Lemma~\ref{lem:baire1-fsigma}.
(2)$\Rightarrow$(1): $f|_{M\setminus N}\in\mathcal B_1$ by Theorem~\ref{thm:baire1-char} or
\ref{thm:sphere-baire1}, and $M\setminus N$ is $\Gd$; extend it to $M$ by
Theorem~\ref{thm:extension}(b) or \ref{thm:sphere-extension}.
(2)$\Leftrightarrow$(3): $f^{-1}(F)\setminus N=(M\setminus N)\setminus(f|_{M\setminus N})^{-1}(Y\setminus F)$;
a subset of the $\Gd$ set $M\setminus N$ is $\Gd$ in $M$ iff it is $\Gd$ in $M\setminus N$, and complements
in $M\setminus N$ exchange $\Fs(M\setminus N)$ and $\Gd(M\setminus N)$.
\end{proof}

\begin{lemma}[invariance of honorary Baire class 2]\label{lem:honorary-invariance}
\lean{TJK.IsHonoraryBaireTwo.restrict, TJK.IsHonoraryBaireTwo.comp_homeomorph}
\uses{def:fn-classes}
If $f:M\to Y$ is of honorary Baire class $\le2$, so are $f|_{M'}$ for $M'\subseteq M$ and $f\circ h$ for a
homeomorphism $h:M''\to M$.
\end{lemma}
\begin{proof}
Restrict (resp.\ compose with $h$) the Baire-one function, its continuous approximants and the countable set.
\end{proof}

\subsection{Points of continuity}

\begin{theorem}[Baire, Kuratowski {\cite[p.~301]{Kuratowski}}]\label{thm:kuratowski}
\lean{TJK.IsBaireOne.meagre_discontinuity}
\uses{def:fn-classes}
If $A$ and $Y$ are metric spaces and $f\in\mathcal B_1(A,Y)$, then the set of points of $A$ at which
$f$ is discontinuous is of the first category in $A$.
\end{theorem}
\begin{proof}
Let $f_n\to f$ be continuous. For $\eps>0$ and $N\in\N$ the set
$F^\eps_N=\{x:\dist(f_n(x),f_m(x))\le\eps\ \forall n,m\ge N\}$ is closed and $\bigcup_NF^\eps_N=A$. Put
$G^\eps=\bigcup_N\inter F^\eps_N$. Then $A\setminus G^\eps\subseteq\bigcup_N(F^\eps_N\setminus\inter
F^\eps_N)$, a countable union of closed nowhere dense sets. If $x\in\inter F^\eps_N$, then
$\dist(f,f_N)\le\eps$ on a neighbourhood of $x$, and $f_N$ is continuous, so
$\dist(f(x'),f(x))\le3\eps$ for $x'$ near $x$ and the oscillation of $f$ at $x$ is at most $6\eps$. Hence the discontinuity set is contained in $\bigcup_k(A\setminus G^{1/k})$.
\end{proof}

\begin{corollary}\label{cor:baire-continuity-point}
\lean{TJK.IsBaireOne.exists_continuousAt}
\uses{thm:kuratowski}
If $P$ is a nonempty complete metric space, $Y$ a metric space and $f\in\mathcal B_1(P,Y)$, then $f$ is continuous at some
point of $P$.
\end{corollary}
\begin{proof}
By the Baire category theorem $P$ is not of the first category in itself.
\end{proof}

\begin{lemma}\label{lem:meager-subspace}
\lean{TJK.isMeagre_of_isMeagre_subtype}
A subset of $A\subseteq M$ that is of the first category in $A$ is of the first category in $M$.
\end{lemma}
\begin{proof}
If $N\subseteq A$ is nowhere dense in $A$ and an open $O\ne\emptyset$ of $M$ lies in $\cl_M N$, then
$O$ meets $A$ and $O\cap A\subseteq\cl_A N$, a contradiction.
\end{proof}

\subsection{Higher Borel and Baire classes}
These are needed only in Section~\ref{sec:discont}.

\begin{definition}[{\cite[\S3]{Diss}}]\label{def:additive-classes}
\lean{TJK.AddClass, TJK.MulClass}
\uses{def:borel-classes}
$P^1(M)$ is the family of open subsets of $M$ and $Q^1(M)$ that of closed subsets. For an ordinal
$\xi>1$, $P^\xi(M)$ consists of the countable unions of members of $\bigcup_{\eta<\xi}Q^\eta(M)$, and
$Q^\xi(M)$ of the countable intersections of members of $\bigcup_{\eta<\xi}P^\eta(M)$. A function
$f:M\to Y$ is \emph{of class $P^\xi$} if $f^{-1}(U)\in P^\xi(M)$ for every open $U\subseteq Y$.
\end{definition}

\begin{lemma}\label{lem:additive-basic}
\lean{TJK.mulClass_iff_compl, TJK.AddClass.mono, TJK.AddClass.iUnion, TJK.AddClass.inter,
  TJK.MulClass.iInter, TJK.AddClass.preimage, TJK.AddClass.subtype_iff, TJK.AddClass.measurableSet}
\uses{def:additive-classes, lem:borel-basic}
(a) $E\in Q^\xi$ iff $M\setminus E\in P^\xi$. (b) $P^\xi$ and $Q^\xi$ are closed under finite unions and
intersections, $P^\xi$ under countable unions, $Q^\xi$ under countable intersections. (c) If $\eta<\xi$
then $P^\eta\cup Q^\eta\subseteq P^\xi\cap Q^\xi$. (d) $P^2(M)$ is the class of $\Fs$ sets and
$Q^2(M)$ that of $\Gd$ sets. (e) Continuous preimages preserve each class. (f) For $L\subseteq M$,
$P^\xi(L)=\{D\cap L:D\in P^\xi(M)\}$. (g) Every member of every class is a Borel set.
\end{lemma}
\begin{proof}
Transfinite induction on $\xi$. (c) for $\eta=1$ uses that in a metric space open sets are $\Fs$ and closed
sets are $\Gd$ (Lemma~\ref{lem:borel-basic}(a)); for $\eta>1$ it is immediate from the definitions. For (b), $(\bigcup A_n)\cap(\bigcup B_m)=\bigcup A_n\cap B_m$, and
$A_n\in Q^{\eta_1}$, $B_m\in Q^{\eta_2}$ both lie in $Q^{\max(\eta_1,\eta_2)}$ by (c). The rest is
routine.
\end{proof}

\begin{lemma}\label{lem:baire-precomp}
\lean{TJK.BaireClass.comp_continuous, TJK.BaireClass.mono}
\uses{def:fn-classes}
If $f:M\to Y$ is of Baire class $\xi$ and $g:K\to M$ is continuous, then $f\circ g$ is of Baire class
$\xi$. In particular restrictions to subspaces keep the class. If $1\le\xi\le\xi'$, Baire class $\xi$
is contained in Baire class $\xi'$.
\end{lemma}
\begin{proof}
Induction on $\xi$; for the last claim use constant sequences.
\end{proof}

\begin{lemma}\label{lem:baire-imp-class}
\lean{TJK.BaireClass.addClass_preimage}
\uses{def:additive-classes, lem:additive-basic, lem:baire-precomp, lem:baire1-fsigma}
If $f:M\to Y$ is of Baire class $\xi$ ($Y$ any metric space), then $f$ is of class $P^{\xi+1}$. In
particular it is Borel measurable.
\end{lemma}
\begin{proof}
For $\xi=0$ preimages of open sets are open. For $\xi\ge1$ let $f_n\to f$ with $f_n$ of class
$\eta_n<\xi$. With $V_k$ as in Lemma~\ref{lem:baire1-fsigma},
$f^{-1}(U)=\bigcup_{k,N}\bigcap_{n\ge N}f_n^{-1}(V_k)$; by induction $f_n^{-1}(V_k)\in
Q^{\eta_n+1}\subseteq Q^\xi$, so the intersections lie in $Q^\xi$ and the union in $P^{\xi+1}$.
\end{proof}

\begin{fact}[Banach {\cite{Banach}}]\label{fact:banach}
\lean{TJK.banach_baireClass_of_addClass}
\uses{def:additive-classes, def:fn-classes}
\source{Banach 1931, quoted in Diss \S3}
Let $Y$ be a separable, arcwise connected metric space and $\xi\ge2$. If $f:M\to Y$ is of class
$P^{\xi+1}$, then $f$ is of Baire class $\xi$.
\end{fact}

\begin{fact}[Lebesgue--Hausdorff--Banach]\label{fact:borel-baire}
\lean{TJK.borel_measurable_exists_baireClass}
\uses{def:fn-classes}
\source{Hausdorff, Set Theory p.~294 (for Y = R, as quoted in Diss \S3); general Y: Banach 1931 with the Borel hierarchy; cf. Kuratowski, Topology I \S31}
Let $Y$ be a separable, arcwise connected metric space. Every Borel-measurable $f:M\to Y$ is of Baire
class $\xi$ for some countable ordinal $\xi$.
\end{fact}

\begin{lemma}[{\cite[Lemma 18]{Diss}}]\label{lem:equiv}
\lean{TJK.CountablyEquiv, TJK.CountablyEquiv.equivalence, TJK.CountablyEquiv.diff,
  TJK.CountablyEquiv.iInter, TJK.CountablyEquiv.iUnion}
Write $A\approx B$ if $A\setminus B$ and $B\setminus A$ are countable. This is an equivalence relation;
$A\approx E$ implies $S\setminus A\approx S\setminus E$; and if $A_n\approx E_n$ for $n$ in a countable
set, then $\bigcap A_n\approx\bigcap E_n$ and $\bigcup A_n\approx\bigcup E_n$.
\end{lemma}
\begin{proof}
$(\bigcap A_n)\setminus(\bigcap E_n)\subseteq\bigcup_n(A_n\setminus E_n)$, and similarly for the rest.
\end{proof}

\begin{lemma}[{\cite[Lemma 19]{Diss}}, {\cite[Lemma 1]{MMJ}}]\label{lem:intervals-open}
\lean{TJK.countable_biUnion_diff_interior_real, TJK.countable_biUnion_diff_interior_circle}
\uses{lem:equiv}
If $\mathcal I$ is a family of nondegenerate intervals of $\R$ (resp.\ nondegenerate closed arcs of $C$),
then $\bigcup_{I\in\mathcal I}I\setminus\bigcup_{I\in\mathcal I}I^*$ is countable. In particular every
union of nondegenerate intervals is $\approx$ an open set.
\end{lemma}
\begin{proof}
$\bigcup I^*$ is open, so it is the disjoint union of its countably many components $J_n$. If $x_0$
lies in $\bigcup I$ but not in $\bigcup I^*$, then $x_0$ is an endpoint of some $I_0$; $I_0^*\subseteq
J_n$ for some $n$, so $x_0\in\overline{J_n}\setminus J_n$ is an endpoint of $J_n$. There are countably
many such endpoints.
\end{proof}

\begin{lemma}[two-sided limit points]\label{lem:two-sided}
\lean{TJK.countable_not_twoSided}
For $F\subseteq\R$, the set of points of $\overline F$ that are not two-sided limit points of $F$ is
countable. (Here $x$ is a \emph{two-sided limit point} of $F$ if every interval $(x-r,x)$ and $(x,x+r)$, $r>0$,
meets $F$.)
\end{lemma}
\begin{proof}
If $x\in\overline F$ and $(x,x+r)\cap F=\emptyset$, then $(x,x+r)$ misses $\overline F$, so $x$ is the
left endpoint of a component of $\R\setminus\overline F$. There are countably many components; argue
symmetrically on the left.
\end{proof}

\begin{lemma}[{\cite[Lemma 20]{Diss}}]\label{lem:monotone-extend}
\lean{TJK.MonotoneOn.exists_monotone_extension_of_dist_le}
Let $h$ be weakly increasing on a nonempty $E\subseteq\R$ with $\abs{x-h(x)}\le1$ on $E$. Then $h$
extends to a weakly increasing $h_1:\R\to\R$.
\end{lemma}
\begin{proof}
Let $e=\inf E$. For $x>e$ put $h_1(x)=\sup h((-\infty,x]\cap E)$, which is finite since
$h(t)\le t+1\le x+1$ there. If $e>-\infty$ put $h_1(x)=\inf h(E)$ for $x\le e$, finite since
$h\ge e-1$ on $E$. Monotonicity and $h_1|_E=h$ are immediate.
\end{proof}

\begin{lemma}[{\cite[Lemma 21]{Diss}}]\label{lem:monotone-compose}
\lean{TJK.BaireClass.comp_monotone}
\uses{lem:baire-precomp}
If $f:\R\to Y$ is of Baire class $\xi$ and $h:\R\to\R$ is weakly increasing, there is a countable
$N\subseteq\R$ such that $(f\circ h)|_{\R\setminus N}$ is of Baire class $\xi$.
\end{lemma}
\begin{proof}
Let $N$ be the (countable) set of discontinuities of $h$ and apply Lemma~\ref{lem:baire-precomp} to
$h|_{\R\setminus N}$.
\end{proof}

\begin{lemma}[{\cite[Lemma 22]{Diss}}]\label{lem:baire-criterion}
\lean{TJK.baireClass_of_sandwich}
\uses{fact:banach, lem:additive-basic}
Let $Y$ be a separable arcwise connected metric space, $E$ a metric space, $\xi\ge2$, and
$\varphi:E\to Y$ such that for every open $U\subseteq Y$ there is $T\in P^{\xi+1}(E)$ with
$\varphi^{-1}(U)\subseteq T\subseteq\varphi^{-1}(\overline U)$. Then $\varphi$ is of Baire class $\xi$.
\end{lemma}
\begin{proof}
Let $\mathcal B$ be a countable base of $Y$. For open $W$ let $\mathcal B(W)=\{B\in\mathcal
B:\overline B\subseteq W\}$; then $W=\bigcup_{\mathcal B(W)}B=\bigcup_{\mathcal B(W)}\overline B$ (for
$p\in W$ with $S(r,p)\subseteq W$ choose $p\in B\subseteq S(r/2,p)$). With $T(B)$ as in the hypothesis,
$\varphi^{-1}(W)=\bigcup_{B\in\mathcal B(W)}T(B)\in P^{\xi+1}(E)$. Conclude by Fact~\ref{fact:banach}.
\end{proof}

% ===========================================================================
\section{Arcs, crosscuts and transfer}\label{sec:plane}
% ===========================================================================

\subsection{Arcs and curvilinear limits}

\begin{definition}[arcs]\label{def:arcs}
\lean{TJK.IsArc, TJK.IsArcAt, TJK.IsCrosscut}
A \emph{path} in $\R^n$ is a continuous $g:[0,1]\to\R^n$; a \emph{simple arc} is the image of an
injective path, and its \emph{endpoints} are $g(0),g(1)$. Let $\Omega\subseteq\R^n$ be open and
$p\in\partial\Omega$. An \emph{arc at $p$} (in $\Omega$) is a simple arc $\gamma$ with one endpoint $p$
and $\gamma\setminus\{p\}\subseteq\Omega$. A \emph{crosscut} of $\D$ is a simple arc whose endpoints
are distinct points of $C$ and whose other points lie in $\D$; a crosscut of $\Hp$ is defined in the
same way with $\R$ in place of $C$.
\end{definition}

\begin{definition}[curvilinear convergence, boundary functions, cluster sets]\label{def:ccset}
\lean{TJK.curvConvSet, TJK.IsBoundaryFunction, TJK.clusterSet}
\uses{def:arcs}
Let $\Omega$ be open, $B\subseteq\partial\Omega$, $Y$ a metric space and $f:\Omega\to Y$. For an arc
$\gamma$ at $p$, $f(z)\to q$ \emph{along $\gamma$} means $f(z)\to q$ as $z\to p$ with
$z\in\gamma\setminus\{p\}$. The \emph{set of curvilinear convergence} of $f$ (relative to $B$) is
\[
\CC(f)=\{p\in B:\ \text{there are an arc $\gamma$ at $p$ and $q\in Y$ with $f\to q$ along $\gamma$}\}.
\]
If $E\subseteq B$, a function $\varphi:E\to Y$ is a \emph{boundary function} for $f$ if for every
$p\in E$ there is an arc $\gamma$ at $p$ with $f\to\varphi(p)$ along $\gamma$. The \emph{cluster set}
of $f$ along $\gamma$ is $C(f,\gamma)=\{q:\exists z_n\in\gamma\setminus\{p\},\ z_n\to p,\ f(z_n)\to q\}$.
The pairs $(\Omega,B)$ used are $(\D,C)$, $(\Hp,\R)$ and $(Q^\circ,X)$, and in Section~\ref{sec:cube}
the open cube with its whole boundary.
\end{definition}

\begin{fact}[Tietze {\cite{Tietze}}]\label{fact:tietze-arc}
\lean{TJK.exists_arc_in_path}
\uses{def:arcs}
\source{Tietze 1919; also: every path-connected Hausdorff space is arc-connected}
If $g$ is a path in $\R^n$ with $g(0)\ne g(1)$, then $g([0,1])$ contains a simple arc with endpoints
$g(0)$ and $g(1)$.
\end{fact}

\begin{corollary}[paths may replace arcs]\label{cor:path-arc}
\lean{TJK.exists_arcAt_of_path, TJK.tendsto_of_nhdsWithin_le_map, TJK.clusterSet_subset_of_nhdsWithin_le_map}
\uses{fact:tietze-arc, def:ccset}
Let $\Omega$ be open, $p\in\partial\Omega$, and $g:[0,1]\to\R^n$ continuous with $g(1)=p$ and
$g([0,1))\subseteq\Omega$. Then $g([0,1])$ contains an arc $\gamma$ at $p$, and for every
$f:\Omega\to Y$: if $f(g(t))\to q$ as $t\to1$ then $f\to q$ along $\gamma$; and $C(f,\gamma)$ is
contained in the set of limits of $f(g(t_n))$ with $t_n\to1$.
\end{corollary}
\begin{proof}
Fact~\ref{fact:tietze-arc} gives a simple arc $\gamma\subseteq g([0,1])$ from $g(0)$ to $p$; since
$g(t)\ne p$ for $t<1$, $\gamma\setminus\{p\}\subseteq\Omega$. If $z_k\in\gamma\setminus\{p\}$ and
$z_k\to p$, write $z_k=g(t_k)$ with $t_k<1$; a subsequence of $t_k$ bounded away from $1$ would have a
limit $t^*<1$ with $g(t^*)=p$, which is impossible. So $t_k\to1$.
\end{proof}

\begin{lemma}[subarcs and exits]\label{lem:subarc}
\lean{TJK.IsArcAt.exists_subarc_subset, TJK.IsArcAt.tendsto_subarc, TJK.IsArcAt.clusterSet_subarc,
  TJK.IsArc.nhdsWithin_subarc, TJK.exists_exit_subpath, TJK.IsArcAt.exists_subarc_firstHit,
  TJK.IsArcAt.exists_subarc_strip}
\uses{def:ccset}
(a) If $\gamma$ is an arc at $p$ and $W$ is a neighbourhood of $p$, there is an arc $\gamma'\subseteq\gamma\cap
W$ at $p$; limits and cluster sets along $\gamma'$ and $\gamma$ coincide. If $f\to q$ along $\gamma$ and
$V\ni q$ is open, some such $\gamma'$ has $f(\gamma'\setminus\{p\})\subseteq V$.
(b) If $g$ is a path, $w=g(t_0)$, and $\abs{g(t_1)-w}\ge r>0$, there is a closed interval $J$ with
endpoints $t_0$ and some $t'$ between $t_0$ and $t_1$ such that $g(J)\subseteq\overline{S(r,w)}$ and
$\abs{g(t')-w}=r$; in particular $\diam g(J)\ge r$.
(c) If $\gamma$ is an arc at $p$ in $\Hp$ and $n$ is large enough that $\gamma$ meets $X_n$, then the subarc
of $\gamma$ from $p$ to its first point on $X_n$ is an arc at $p$ contained in $\overline{\Hp}_n$ with all
points other than its endpoints in $\Hp_n$.
\end{lemma}
\begin{proof}
Parametrize and take initial segments; for (b) let $t'$ be the first parameter between $t_0$ and $t_1$
with $\abs{g(t')-w}\ge r$, and use continuity.
\end{proof}

\begin{lemma}[cluster sets]\label{lem:cluster}
\lean{TJK.clusterSet_eq_iInter_subarc, TJK.clusterSet_nonempty, TJK.isCompact_clusterSet,
  TJK.isConnected_clusterSet, TJK.tendsto_iff_clusterSet_eq_singleton, TJK.clusterSet_comp}
\uses{def:ccset, lem:subarc}
Let $K$ be a compact metric space, $f:\Omega\to K$ continuous and $\gamma$ an arc at $p$. Then
$C(f,\gamma)=\bigcap_{\gamma'}\overline{f(\gamma'\setminus\{p\})}$, the intersection over the subarcs
$\gamma'$ of $\gamma$ at $p$; it is nonempty, compact and connected; and $f\to q$ along $\gamma$ iff
$C(f,\gamma)=\{q\}$. If $\Theta:K\to K'$ is continuous into a metric space, then
$C(\Theta\circ f,\gamma)=\Theta(C(f,\gamma))$.
\end{lemma}
\begin{proof}
The sets $\overline{f(\gamma'\setminus\{p\})}$ are compact, connected (images of half-open arcs) and
decrease as $\gamma'$ shrinks; a decreasing intersection of nonempty compact connected sets is
nonempty, compact and connected. The claim about limits is the usual compactness argument. For the last
claim: $\supseteq$ by continuity of $\Theta$; if $\Theta(f(z_n))\to q'$ with $z_n\to p$ along $\gamma$, a subsequence
of $f(z_n)$ converges in the compact $K$ to some $q\in C(f,\gamma)$, and $\Theta(q)=q'$.
\end{proof}

\begin{fact}[Tietze extension]\label{fact:tietze-ext}
\lean{TJK.tietze_bounded_complex}
\source{Mathlib: ContinuousMap.exists_extension / BoundedContinuousFunction extensions (apply to real and imaginary parts)}
A bounded continuous function from a closed subset of a metric space to $\C$ extends to a bounded continuous function on the
whole space.
\end{fact}

\begin{fact}[hyperspace]\label{fact:hyperspace}
\lean{TJK.compactSpace_nonemptyCompacts}
\source{Hausdorff, Mengenlehre p.~150; Mathlib: EMetric.NonemptyCompacts.compactSpace}
If $K$ is a compact metric space, the space $\mathcal C(K)$ of nonempty closed subsets of $K$ with the Hausdorff metric is
compact.
\end{fact}

\begin{fact}[Baire category theorem]\label{fact:baire-category}
\lean{TJK.baire_category}
\source{Mathlib: BaireSpace instances, dense_iInter_of_isOpen}
A complete metric space is not of the first category in itself; in it, a countable intersection of dense open sets is dense.
\end{fact}

\subsection{Crosscuts}

\begin{fact}[crosscut theorem]\label{fact:crosscut}
\lean{TJK.crosscut_theorem, TJK.closure_inter_circle_of_frontier, TJK.IsCrosscut.exists_angles}
\uses{def:arcs}
\source{Jordan curve theorem; Newman, Elements, Theorem 11.8 (the $\theta$-curve theorem)}
Let $\sigma$ be a crosscut of $\D$ with endpoints $a\ne b$, and let $A_1,A_2$ be the two closed arcs of
$C$ with endpoints $a,b$. Then $\D\setminus\sigma$ has exactly two components $U_1,U_2$, and the frontier
of $U_i$ is $\sigma\cup A_i$ ($i=1,2$). In particular $\overline{U_i}\cap C=A_i$ and
$\sigma\subseteq\overline{U_1}\cap\overline{U_2}$.
\end{fact}

\begin{lemma}[disjoint components have disjoint boundary arcs]\label{lem:disjoint-arcs}
\lean{TJK.disjoint_boundaryArcs}
\uses{fact:crosscut}
Let $\sigma,\tau$ be crosscuts of $\D$, $U$ a component of $\D\setminus\sigma$ and $V$ a component of
$\D\setminus\tau$, with $U\cap V=\emptyset$. Then the arcs $\overline U\cap C$ and $\overline V\cap C$ have
disjoint relative interiors.
\end{lemma}
\begin{proof}
Let $p$ be interior to both arcs. Then $p$ is not an endpoint of $\sigma$ or $\tau$, so $p\notin\sigma\cup\tau$,
and there is $r>0$ with $S(r,p)\cap(\sigma\cup\tau)=\emptyset$. The convex set $S(r,p)\cap\D$ lies in one
component of $\D\setminus\sigma$; as $p\in\overline U$ it meets $U$, so it lies in $U$. Likewise it lies in
$V$, contradicting $U\cap V=\emptyset$.
\end{proof}

\begin{definition}[charts]\label{def:charts}
\lean{TJK.cayley, TJK.chartH, TJK.chartD, TJK.stripExp, TJK.openSquare, TJK.baseX, TJK.gOpen,
  TJK.gClosed}
\begin{enumerate}[label=(\alph*)]
\item The Cayley map $\kappa(z)=(z-i)/(z+i)$.
\item $\Psi_\Hp\langle x,y\rangle=\tan(\pi x/2)+i\,y/(1-y)$ for $\abs x<1$, $0\le y<1$.
\item $\Psi_\D\langle x,y\rangle=1-1/w$ with $w=\tfrac12(1+y)+i\tan(\pi x/2)$, for $\abs x<1$, $0\le y<1$.
\item $\mathcal E_h(z)=e^{\pi z/h}$ for $h>0$.
\end{enumerate}
\end{definition}

\begin{lemma}[charts]\label{lem:charts}
\lean{TJK.cayley_isChart, TJK.tendsto_cayley_cobounded, TJK.tendsto_norm_cayleyInv,
  TJK.chartH_isChart, TJK.chartH_ofReal, TJK.strictMonoOn_tan_pi_mul_div_two, TJK.chartD_isChart,
  TJK.isClosed_gClosed, TJK.exists_ball_inter_disc_subset_gClosed, TJK.ofReal_not_mem_gClosed,
  TJK.stripExp_isChart, TJK.stripExp_ofReal, TJK.stripExp_add_I}
\uses{def:charts}
(a) $\kappa$ is a homeomorphism of $\Hp\cup\R$ onto $(\D\cup C)\setminus\{1\}$ mapping $\Hp$ onto $\D$
and $\R$ onto $C\setminus\{1\}$; $\abs{\kappa(z)-1}\to0$ iff $\abs z\to\infty$.
(b) $\Psi_\Hp$ is a homeomorphism of $Q^\circ\cup X$ onto $\Hp\cup\R$ mapping $X$ onto $\R$
increasingly.
(c) $\Psi_\D$ is a homeomorphism of $Q^\circ\cup X$ onto $G^\circ\cup(C\setminus\{1\})$, where
$G^\circ=\{z\in\D:\operatorname{Re}(1-z)^{-1}<1\}$, mapping $X$ onto $C\setminus\{1\}$. The set
$G''=\{z\in\D:\operatorname{Re}(1-z)^{-1}\le3/4\}$ is closed in $\D$, contains a neighbourhood in $\D$
of each point of $C\setminus\{1\}$, and is disjoint from the radius $[0,1)$.
(d) $\mathcal E_h$ maps $\{0\le\operatorname{Im}z\le h\}$ homeomorphically onto
$(\Hp\cup\R)\setminus\{0\}$, the line $\operatorname{Im}z=0$ increasingly onto $(0,\infty)$ and the line
$\operatorname{Im}z=h$ decreasingly onto $(-\infty,0)$.
\end{lemma}
\begin{proof}
(a) $\abs{\kappa(z)}<1$ iff $\abs{z-i}<\abs{z+i}$ iff $\operatorname{Im}z>0$; the inverse is
$w\mapsto i(1+w)/(1-w)$. (b) is a product of homeomorphisms $(-1,1)\to\R$ and $[0,1)\to[0,\infty)$.
(c) $z\mapsto(1-z)^{-1}$ maps $\D$ onto $\{\operatorname{Re}w>\frac12\}$ and $C\setminus\{1\}$ onto
$\{\operatorname{Re}w=\frac12\}$ (since $\abs{1-1/w}<1\iff\abs{w-1}<\abs w$); $\Psi_\D$ is the inverse
of this map composed with a homeomorphism of $Q^\circ\cup X$ onto $\{\frac12\le\operatorname{Re}w<1\}$.
Near a point $x\in C\setminus\{1\}$, $\operatorname{Re}(1-z)^{-1}$ is close to $\frac12$; on $[0,1)$ it is
$\ge1$. (d) is elementary.
\end{proof}

\begin{lemma}[crosscuts of the half-plane]\label{lem:crosscut-H}
\lean{TJK.crosscutH_theorem}
\uses{fact:crosscut, lem:charts}
Let $\sigma$ be a crosscut of $\Hp$ with endpoints $a<b$. Then $\Hp\setminus\sigma$ has exactly two
components, a bounded one $V$ with $\overline V\cap\R=[a,b]$, and an unbounded one $U$ with
$\overline U\cap\R=\R\setminus(a,b)$; moreover $\sigma\subseteq\overline U\cap\overline V$ and
$U\supseteq\{\operatorname{Im}z>\max_\sigma\operatorname{Im}\}$.
\end{lemma}
\begin{proof}
$\kappa(\sigma)$ is a crosscut of $\D$ with endpoints in $C\setminus\{1\}$
(Lemma~\ref{lem:charts}(a)); apply Fact~\ref{fact:crosscut}. The component whose closed arc avoids $1$
has closure a compact subset of $\overline\D\setminus\{1\}$, so its preimage $V$ is bounded, and $\kappa$
maps $[a,b]$ onto that arc. The half-plane above $\sigma$ is connected, unbounded and disjoint from
$\sigma$, so it lies in $U$.
\end{proof}

\begin{lemma}[order of disjoint arcs in a strip]\label{lem:strip}
\lean{TJK.strip_order}
\uses{lem:crosscut-H, lem:charts}
Let $h>0$, $S=\{0<\operatorname{Im}z<h\}$, and let $\alpha,\alpha'$ be disjoint simple arcs, $\alpha$ from
$x\in\R$ to $p$ with $\operatorname{Im}p=h$, $\alpha'$ from $x'\in\R$ to $p'$ with
$\operatorname{Im}p'=h$, all other points in $S$. If $x<x'$ then $\operatorname{Re}p<\operatorname{Re}p'$.
\end{lemma}
\begin{proof}
By Lemma~\ref{lem:charts}(d), $\mathcal E_h(\alpha)$ is a crosscut of $\Hp$ from $u=e^{\pi x/h}$ to
$v=-e^{\pi\operatorname{Re}p/h}$, and $\mathcal E_h(\alpha')$ one from $u'>u$ to
$v'=-e^{\pi\operatorname{Re}p'/h}\ne v$. By Lemma~\ref{lem:crosscut-H} the bounded component $V$ of
$\Hp\setminus\mathcal E_h(\alpha)$ has $\overline V\cap\R=[v,u]$. The connected set
$\mathcal E_h(\alpha')\setminus\{u',v'\}$ lies in one component; its closure contains $u'\notin[v,u]$, so
it lies in the unbounded component $U$, and then $v'\in\overline U\cap\R=\R\setminus(v,u)$. As $v'<0<u$
and $v'\ne v$, $v'<v$, i.e.\ $\operatorname{Re}p'>\operatorname{Re}p$.
\end{proof}

\begin{lemma}[geometry near a boundary point]\label{lem:disc-geometry}
\lean{TJK.convex_disc_inter_ball, TJK.isConnected_disc_diff_closedBall, TJK.shortArc_subset_ball,
  TJK.exists_shortArc, TJK.norm_expI_sub_le, TJK.sum_le_two_pi_of_pairwiseDisjoint, TJK.card_mul_le_two_pi}
Let $x\in C$ and $0<\delta<2$.
(a) $\D\cap S(\delta,x)$ is convex.
(b) $\D\setminus\overline{S(\delta,x)}$ is connected and has $-x$ in its closure.
(c) If $x'\in C$ and $\abs{x'-x}<\delta$, the shorter closed arc of $C$ from $x$ to $x'$ lies in $S(\delta,x)$.
Any arc of $C$ between $w_1,w_2$ has length at least $\abs{w_1-w_2}$.
(d) If $A_1,\dots,A_m$ are closed arcs of $C$ with pairwise disjoint relative interiors, their lengths add up to at most
$2\pi$. (The length of an arc is the Lebesgue measure of the corresponding set of angles in $[0,2\pi)$.) In particular,
if the endpoints of each $A_i$ are at distance $\ge\eps$, then $m\eps\le2\pi$.
In Lean a closed arc is $\{e^{it}:\theta\le t\le\theta+\ell\}$ with $0\le\ell\le2\pi$ (\texttt{circleArc}), its interior
is the image of $(\theta,\theta+\ell)$, and its length is $\ell$; (d) is proved by pushing the open angle intervals to
$\R/2\pi\Z$ with its Haar measure.
\end{lemma}
\begin{proof}
(a) Intersection of convex sets. (b) For $z\in\D$ the convex function $t\mapsto\abs{(1-t)z-tx-x}^2$ has
derivative $2(1-\abs z^2)>0$ at $t=0$, hence increases on $[0,1]$; so the segment $[z,-x)$ stays in
$\D\setminus\overline{S(\delta,x)}$ and ends in the convex set $\D\cap S(2-\delta,-x)$, which lies in
$\D\setminus\overline{S(\delta,x)}$. (c) On an arc of angle at most $\pi$ the chord from an endpoint
increases with the angle; chords are shorter than arcs. (d) The interiors are disjoint measurable sets of
angles and the endpoints are null.
\end{proof}

\begin{lemma}\label{lem:open-components}
\lean{TJK.countable_connectedComponents_of_isOpen}
An open subset of $\R^n$ has countably many components; they are open and path-connected, and a
connected subset of the open set lies in one of them.
\end{lemma}
\begin{proof}
$\R^n$ is locally path-connected and second countable.
\end{proof}

\subsection{Transfer}

\begin{lemma}[transfer]\label{lem:transfer}
\lean{TJK.IsChart, TJK.transfer_isArcAt, TJK.transfer_tendsto, TJK.transfer_curvConvSet,
  TJK.transfer_isBoundaryFunction, TJK.transfer_clusterSet}
\uses{def:ccset}
Let $\Omega,\Omega'$ be open, $B\subseteq\partial\Omega$, $B'\subseteq\partial\Omega'$, and let
$\Psi:\Omega\cup B\to\Omega'\cup B'$ be a homeomorphism with $\Psi(\Omega)=\Omega'$. For $f:\Omega'\to Y$:
$\gamma$ is an arc at $p\in B$ iff $\Psi(\gamma)$ is an arc at $\Psi(p)$; $f\circ\Psi\to q$ along
$\gamma$ iff $f\to q$ along $\Psi(\gamma)$, and the cluster sets agree; hence
$\CC(f\circ\Psi)=\Psi^{-1}(\CC(f))$, and $\varphi$ is a boundary function for $f$ on $E'$ iff
$\varphi\circ\Psi$ is one for $f\circ\Psi$ on $\Psi^{-1}(E')$.
\end{lemma}
\begin{proof}
Arcs at $p$ lie in $\Omega\cup\{p\}$, where $\Psi$ is a homeomorphism.
\end{proof}

\begin{lemma}[locality]\label{lem:locality}
\lean{TJK.curvConv_local}
\uses{def:ccset, lem:subarc}
Let $p\in\partial\Omega\cap\partial\Omega'$ and suppose $W\cap\Omega=W\cap\Omega'$ for an open $W\ni p$. If
$f=f'$ on $W\cap\Omega$, then $f$ has a limit $q$ along some arc at $p$ in $\Omega$ iff $f'$ has the limit
$q$ along some arc at $p$ in $\Omega'$.
\end{lemma}
\begin{proof}
Lemma~\ref{lem:subarc}(a).
\end{proof}

% ===========================================================================
\section{Sets of curvilinear convergence}\label{sec:cc}
% ===========================================================================
Sources: \cite{MMJ}; \cite[\S4, Lemmas 7--8, Theorem 3]{Diss}. We work in $\D$ and transfer to $\Hp$.

\begin{definition}[accessible points and crosscut arcs]\label{def:acc}
\lean{TJK.McMillan.Gamma, TJK.McMillan.crossArcOpen, TJK.McMillan.I0, TJK.McMillan.acc}
\uses{def:arcs, fact:crosscut}
For a crosscut $\gamma$ of $\D$ with $0\notin\gamma$, let $U$ be the component of $\D\setminus\gamma$ containing
$0$ and $L(\gamma)^*=C\setminus\overline U$; by Theorem~\ref{fact:crosscut} it is the open arc of $C$ between the
endpoints of $\gamma$ on the side of the other component. For a domain $\Omega\subseteq\D\setminus\{0\}$ let
$\Gamma(\Omega)$ be the set of crosscuts $\gamma$ with $\gamma\cap\D\subseteq\Omega$,
\[
I_0(\Omega)=\bigcup_{\gamma\in\Gamma(\Omega)}L(\gamma)^*,
\]
and let $\acc(\Omega)$ be the set of $x\in C$ for which there is an arc at $x$ in $\Omega$. Note that $I_0(\Omega)$ is
open in $C$ and $\acc(\Omega)\subseteq\overline\Omega$.
\end{definition}

\begin{lemma}[{\cite[Lemma 2]{MMJ}}]\label{lem:mm-lemma2}
\lean{TJK.McMillan.countable_acc_diff_I0, TJK.McMillan.exists_arc_join,
  TJK.McMillan.exists_Ioo_subset_crossArcOpen}
\uses{def:acc, lem:intervals-open, fact:tietze-arc, fact:crosscut}
For every domain $\Omega\subseteq\D\setminus\{0\}$, $\acc(\Omega)\setminus I_0(\Omega)$ is countable.
\end{lemma}
\begin{proof}
We may assume $\acc(\Omega)$ has two or more points. Let $a\in\acc(\Omega)$ and $a'\ne a$ in $\acc(\Omega)$,
with arcs $\gamma,\gamma'$ at $a,a'$ in $\Omega$, and join their other endpoints by a path in $\Omega$. By
Fact~\ref{fact:tietze-arc} the resulting path from $a$ to $a'$ contains a simple arc $\delta$ from $a$ to $a'$;
it is a crosscut with $\delta\cap\D\subseteq\Omega$ (so $0\notin\delta$). Write $a=e^{i\theta}$,
$a'=e^{i(\theta+\ell)}$, $0<\ell<2\pi$. By Theorem~\ref{fact:crosscut}, $L(\delta)^*$ is one of the two open arcs
$e^{i(s,t)}$ with $(s,t)=(\theta,\theta+\ell)$ or $(\theta+\ell,\theta+2\pi)$, and $a\in e^{i[s,t]}$. So
$\acc(\Omega)$ lies in a union of closed arcs $e^{i[s,t]}$ whose open arcs $e^{i(s,t)}$ lie in $I_0(\Omega)$, and
Lemma~\ref{lem:intervals-open} shows that $\acc(\Omega)\setminus I_0(\Omega)$ is countable.
\end{proof}

\begin{lemma}[{\cite[Lemma 3]{MMJ}}]\label{lem:mm-lemma3}
\lean{TJK.McMillan.inter_nonempty, TJK.McMillan.inter_nonempty_of_split, TJK.McMillan.exists_split}
\uses{def:acc, fact:crosscut}
Let $\Omega_1,\Omega_2\subseteq\D\setminus\{0\}$ be connected. If the sets
$I_0(\Omega_1)\cap\overline{\acc(\Omega_1)}$ and $I_0(\Omega_2)\cap\overline{\acc(\Omega_2)}$ have a common
point $a$, then $\Omega_1\cap\Omega_2\ne\emptyset$.
\end{lemma}
\begin{proof}
Suppose $\Omega_1\cap\Omega_2=\emptyset$. Choose $\gamma_i\in\Gamma(\Omega_i)$ with $a\in L(\gamma_i)^*$, and let
$U_i,V_i$ be the components of $\D\setminus\gamma_i$, with $0\in U_i$; so $a\notin\overline{U_i}$, while
$a\in\overline{\acc(\Omega_i)}\subseteq\overline{\Omega_i}$. The sets $\tau_i=\gamma_i\cap\D$ are nonempty, connected,
contained in $\Omega_i$ and in $\overline{U_i}$ (Theorem~\ref{fact:crosscut}), and $\tau_1\cap\gamma_2=\emptyset=
\tau_2\cap\gamma_1$. Hence each $\tau_i$ lies in one component of the other complement.

If $\tau_2\subseteq U_1$: the connected set $\Omega_2$ misses $\gamma_1$ (as $\gamma_1\cap\D\subseteq\Omega_1$), so
it lies in $U_1$ or $V_1$; it meets $U_1$ (it contains $\tau_2$), so $\Omega_2\subseteq U_1$ and
$a\in\overline{\Omega_2}\subseteq\overline{U_1}$, a contradiction. Symmetrically $\tau_1\not\subseteq U_2$.

So $\tau_1\subseteq V_2$ and $\tau_2\subseteq V_1$. As $\tau_1\subseteq\overline{U_1}$ and $V_2$ is open, $U_1$ meets
$V_2$; also $0\in U_1\cap U_2$. The connected set $U_1$ meets both components of $\D\setminus\gamma_2$, so it meets
$\gamma_2\cap\D=\tau_2\subseteq V_1$, which is absurd.
\end{proof}

\begin{theorem}[McMillan {\cite[Theorem 5]{McM}}; {\cite[Theorem 1]{MMJ}}]\label{thm:mcmillan}
\lean{TJK.McMillan.isFσδ_curvConvSet, TJK.McMillan.isFσδ_curvConvSet',
  TJK.McMillan.isFσδ_of_skeleton, TJK.McMillan.exists_path_of_seq,
  TJK.McMillan.exists_arcAt_of_chain}
\uses{lem:mm-lemma2, lem:mm-lemma3, cor:path-arc, lem:subarc, lem:open-components, lem:borel-basic}
Let $K$ be a complete separable metric space and $f:\D\to K$ continuous. Then $\CC(f)$ is an $\Fsd$
subset of $C$.
\end{theorem}
\begin{proof}
Let $\{p_k\}$ be dense in $K$ and $\{c_m\}$ dense in $\C$. Put
$Q(n,m)=S(1/(n+1),c_m)\cap\D\setminus\{0\}$, and let $U(n,m,k,l)$, $l\in\N$, be the component of
$O(n,m,k)=f^{-1}(S(2^{-n},p_k))\cap Q(n,m)$ containing $c_l$ (empty if $c_l\notin O(n,m,k)$); every component of the
open set $O(n,m,k)$ is open and nonempty, so it is one of the $U(n,m,k,l)$ (Lemma~\ref{lem:open-components}). Put
$A(n,m,k,l)=\acc U(n,m,k,l)$ and
\[
Y=\bigcap_n\bigcup_{m,k,l}\Bigl(I_0(U(n,m,k,l))\cap\overline{A(n,m,k,l)}\Bigr).
\]
$I_0(\cdot)$ is open in $C$, so $Y$ is $\Fsd$.

\emph{$Y\subseteq\CC(f)$.} Let $y\in Y$ and choose $U_n=U(n,m[n],k[n],l[n])$ with
$y\in I_0(U_n)\cap\overline{\acc U_n}$. By Lemma~\ref{lem:mm-lemma3}, $U_n\cap U_{n+1}$ contains a
point $z_n$; join $z_n$ to $z_{n+1}$ by a path $\gamma_n\subseteq U_{n+1}$. Since
$y\in\overline{U_{n+1}}\subseteq\overline{S(1/(n+2),c_{m[n+1]})}$, $\gamma_n$ lies within $2/(n+2)$ of $y$. So
the paths $\gamma_n$, traversed on $[n,n+1]$ and reparametrized by $t\mapsto(1-t)^{-1}-1$, followed by $y$ at $t=1$,
form a path $g$ with $g(1)=y$, $g([0,1))\subseteq\D$, and $g(t)\in\bigcup_{n\ge N}\gamma_n$ for $t$ near $1$. As
$S(2^{-n},p_{k[n]})$ and $S(2^{-n-1},p_{k[n+1]})$ meet (at $f(z_n)$), $\rho(p_{k[n]},p_{k[n+1]})<2^{-n+1}$, so
$p_{k[n]}\to p$ for some $p\in K$; and $f(\gamma_n)\subseteq S(2^{-n-1},p_{k[n+1]})$, so $f(g(t))\to p$ as $t\to1$.
By Corollary~\ref{cor:path-arc}, $y\in\CC(f)$.

\emph{$\CC(f)\subseteq\bigcap_n\bigcup_{m,k,l}A(n,m,k,l)$.} Let $f\to p'$ along an arc $\gamma_0$ at $x$.
Given $n$, pick $k$ with $\rho(p',p_k)<2^{-n}$ and $m$ with $\abs{x-c_m}<1/(2n+2)$, so that
$S(r,x)\cap\D\subseteq Q(n,m)$ for $r=\min(1/(2n+2),1/2)$. By Lemma~\ref{lem:subarc} a subarc $\gamma_0'$ at $x$ has
$\gamma_0'\setminus\{x\}\subseteq O(n,m,k)$; being connected, it lies in one component, so $x\in A(n,m,k,l)$
for some $l$.

\emph{$\CC(f)\setminus Y$ is countable.} If $x\in\CC(f)\setminus Y$, there is $n$ with
$x\in A(n,m,k,l)$ for some $m,k,l$ but $x\notin I_0(U(n,m,k,l))\cap\overline{A(n,m,k,l)}$, so
$x\in A(n,m,k,l)\setminus I_0(U(n,m,k,l))$, which is countable by Lemma~\ref{lem:mm-lemma2}.

So $\CC(f)$ is the union of an $\Fsd$ set and a countable set, hence $\Fsd$. (In Lean the three steps are
packaged once, for any countable family of open sets $O(n,i)\subseteq\D\setminus\{0\}$, and reused for
Theorem~\ref{thm:mm-thm2}.)
\end{proof}

\begin{theorem}[{\cite[Theorem 2]{MMJ}}]\label{thm:mm-thm2}
\lean{TJK.McMillan.isFσδ_clusterSet_subset, TJK.McMillan.isFσδ_clusterSet_subset',
  TJK.McMillan.mem_map_of_forall_mapClusterPt}
\uses{lem:mm-lemma2, lem:mm-lemma3, cor:path-arc, lem:subarc, lem:cluster, lem:open-components,
  lem:borel-basic, fact:hyperspace}
Let $K$ be a compact metric space, $\mathcal C(K)$ the space of nonempty closed subsets of $K$ with the
Hausdorff metric $\bar\rho$, $\mathcal E\subseteq\mathcal C(K)$ closed, and $f:\D\to K$ continuous. Then
\[
X_{\mathcal E}=\{x\in C:\text{there are an arc $\gamma$ at $x$ and $E\in\mathcal E$ with }C(f,\gamma)\subseteq E\}
\]
is an $\Fsd$ set.
\end{theorem}
\begin{proof}
If $\mathcal E=\emptyset$ there is nothing to prove. $\mathcal E$ is compact (Fact~\ref{fact:hyperspace}), hence has a dense sequence
$P(k)$. For $\eps>0$ let $\mathcal S(\eps,E)=\{a:\dist(a,E)<\eps\}$, which is open, and note that
$\dist(a,F)\le\dist(a,E)+\bar\rho(E,F)$. Define $Q(n,m)$ as before, $U(n,m,k,l)$ as the
components of $f^{-1}(\mathcal S(1/(n+1),P(k)))\cap Q(n,m)$, $A(n,m,k,l)=\acc U(n,m,k,l)$, and $Y$ by the same
formula; $Y$ is $\Fsd$.

\emph{$Y\subseteq X_{\mathcal E}$.} With $U_n$ as before, compactness gives $P\in\mathcal E$ and
$n_1<n_2<\cdots$ with $P(k[n_j])\to P$. The sets $U_{n_j}$ satisfy the hypotheses used above, so as in
Theorem~\ref{thm:mcmillan} we get a path $g$ ending at $y$, with $g(t)\in\bigcup_{j\ge J}U_{n_j}$ for $t$ near $1$,
and (Corollary~\ref{cor:path-arc}) an arc $\gamma$ at $y$ whose cluster set is contained in that of $f\circ g$ at $1$.
Let $p$ be such a cluster point. Given $\eps>0$, for $j$ large $2/(n_j+1)<\eps/2$ and
$\bar\rho(P(k[n_j]),P)<\eps/2$, so $\dist(f(g(t)),P)\le\eps$ for $t$ near $1$, and hence $\dist(p,P)\le\eps$. As $P$
is closed, $p\in P$. So $C(f,\gamma)\subseteq P$.

\emph{$X_{\mathcal E}\subseteq\bigcap_n\bigcup_{m,k,l}A(n,m,k,l)$.} Let $C(f,\gamma_0)\subseteq P_0\in\mathcal E$.
Given $n$, choose $k$ with $\bar\rho(P_0,P(k))<1/(n+1)$, so $C(f,\gamma_0)\subseteq\mathcal S(1/(n+1),P(k))=:O$. Some
subarc of $\gamma_0$ at $x$ is mapped into $O$: otherwise the filter of $f$ along $\gamma_0$ meets the compact
set $K\setminus O$ and has a cluster point there, which would lie in $C(f,\gamma_0)\setminus O$. Conclude as in
Theorem~\ref{thm:mcmillan}; the countability of $X_{\mathcal E}\setminus Y$ is proved in the same way.
\end{proof}

\begin{corollary}[{\cite[Corollaries 1--4]{MMJ}}]\label{cor:mm-corollaries}
\lean{TJK.McMillan.cor1, TJK.McMillan.cor2, TJK.McMillan.cor3, TJK.McMillan.cor4,
  TJK.McMillan.isClosed_subset_greatCircle}
\uses{thm:mm-thm2, lem:cluster}
Let $f:\D\to\Sph$ be continuous. Each of the following sets is $\Fsd$ ((1)--(3) hold for any compact metric
target):
\begin{enumerate}[label=(\arabic*)]
\item (McMillan) $\{x:\exists\gamma\text{ at }x,\ \exists p\in E,\ f\to p\text{ along }\gamma\}$, for closed
  $E\subseteq\Sph$;
\item $\{x:\exists\gamma\text{ at }x,\ \diam C(f,\gamma)\le d\}$, for $d\ge0$;
\item $\{x:\exists\gamma\text{ at }x,\ C(f,\gamma)\subseteq E\}$, for closed $E\subseteq\Sph$;
\item $\{x:\exists\gamma\text{ at }x,\ C(f,\gamma)\text{ is an arc of a great circle}\}$, where (since
  $C(f,\gamma)$ is closed and connected) an ``arc'' may be a point, a closed arc, or a whole great circle; that is,
  $C(f,\gamma)$ is contained in a great circle.
\end{enumerate}
\end{corollary}
\begin{proof}
Apply Theorem~\ref{thm:mm-thm2} with $\mathcal E$ equal to: (1) $\{\{p\}:p\in E\}$, an isometric copy of
the compact set $E$, using that $C(f,\gamma)\subseteq\{p\}$ iff $f\to p$ along $\gamma$ (Lemma~\ref{lem:cluster});
(2) $\{F:\diam F\le d\}$, closed because if $\bar\rho(F,F')<\eps$ then every two points of $F$ are within
$2\eps$ of two points of $F'$, and using $E=C(f,\gamma)$ for the converse; (3) $\{E\}$ (or note the set is empty if
$E=\emptyset$); (4) the sets contained in some great circle $\{v:\langle v,u\rangle=0\}$: if $F_n\to F$ with
$F_n\subseteq\{\langle\cdot,u_n\rangle=0\}$, pass to a subsequence $u_n\to u$; each $v\in F$ is a limit of
$v_n\in F_n$, so $\langle v,u\rangle=\lim\langle v_n,u_n\rangle=0$. With $E=C(f,\gamma)$ for the converse; the
connectedness of $C(f,\gamma)$ (Lemma~\ref{lem:cluster}) justifies the reading of ``arc''.
\end{proof}

\begin{theorem}[McMillan in the half-plane; {\cite[Theorem 3]{Diss}}]\label{thm:mcmillan-H}
\lean{TJK.McMillan.isFσδ_curvConvSet_upperHalfPlane}
\uses{thm:mcmillan, lem:transfer, lem:charts, lem:borel-basic}
If $K$ is a complete separable metric space and $f:\Hp\to K$ is continuous, then $\CC(f)\subseteq\R$ is
$\Fsd$.
\end{theorem}
\begin{proof}
By Lemma~\ref{lem:transfer} with $\Psi=\kappa$ (Lemma~\ref{lem:charts}(a)),
$\CC(f)=\R\cap\kappa^{-1}(\CC(f\circ\kappa^{-1}))$, and the right side is $\Fsd$ by
Theorem~\ref{thm:mcmillan} and Lemma~\ref{lem:borel-basic} (preimage under the continuous map $\kappa|_\R$, then
an $\Fsd$ subset of the closed set $\R$).
\end{proof}

\subsection*{The half-plane lemmas of the dissertation}
For a domain $\Omega\subseteq\Hp$ and a crosscut $\gamma$ of $\Hp$ with endpoints $c<d$ let
$L(\gamma)=[c,d]$, the trace on $\R$ of the closure of the bounded component of $\Hp\setminus\gamma$; let
$I(\Omega)=\bigcup L(\gamma)^*$ over crosscuts with $\gamma\cap\Hp\subseteq\Omega$, and $\acc(\Omega)$ the
points of $\R$ accessible by arcs in $\Omega$. These lemmas are the half-plane forms of
Lemmas~\ref{lem:mm-lemma2} and~\ref{lem:mm-lemma3}; they are not used elsewhere.

\begin{lemma}[{\cite[Lemma 7]{Diss}}]\label{lem:th-lemma7}
\lean{TJK.McMillan.halfPlane_I_eq, TJK.McMillan.IH, TJK.McMillan.accH}
\uses{lem:crosscut-H, fact:tietze-arc}
If $\acc(\Omega)\ne\emptyset$, $a=\inf\acc(\Omega)$ and $b=\sup\acc(\Omega)$, then $I(\Omega)=(a,b)$; that is,
$x\in I(\Omega)$ iff there are $c,d\in\acc(\Omega)$ with $c<x<d$.
\end{lemma}
\begin{proof}
If $x\in L(\gamma)^*=(c,d)$ then $c,d\in\acc(\Omega)$ (initial halves of $\gamma$ from either end are arcs at $c$,
$d$ in $\Omega$), so $a\le c<x<d\le b$. Conversely, for $a<x<b$ pick $c',d'\in\acc(\Omega)$ with $c'<x<d'$, and join
them by a crosscut through $\Omega$ as in Lemma~\ref{lem:mm-lemma2}.
\end{proof}

\begin{lemma}[{\cite[Lemma 8]{Diss}}]\label{lem:th-lemma8}
\lean{TJK.McMillan.halfPlane_inter_nonempty}
\uses{lem:crosscut-H, lem:mm-lemma3}
If $\Omega_1,\Omega_2\subseteq\Hp$ are connected and $I(\Omega_1)\cap\overline{\acc(\Omega_1)}$ and
$I(\Omega_2)\cap\overline{\acc(\Omega_2)}$ meet, then $\Omega_1\cap\Omega_2\ne\emptyset$.
\end{lemma}
\begin{proof}
As in Lemma~\ref{lem:mm-lemma3}, with $U_i$ the unbounded component of $\Hp\setminus\gamma_i$ (any two
unbounded components meet, since each contains a half-plane $\{\operatorname{Im}z>t\}$) and
Lemma~\ref{lem:crosscut-H} in place of Theorem~\ref{fact:crosscut}.
\end{proof}

% ===========================================================================
\section{Boundary functions of continuous functions and homeomorphisms}\label{sec:bf}
% ===========================================================================
Sources: \cite[\S5, Lemma 15, Theorem 5]{Diss}; \cite[Theorems 1--3]{JMM}; \cite[remark after Corollary 4]{MMJ}.

\begin{lemma}[{\cite[Lemma 15]{Diss}}]\label{lem:th-lemma15}
\lean{TJK.isFσMeasurable_of_sandwich}
\uses{def:fn-classes, lem:borel-basic, lem:baire-criterion}
Let $E$ be a metric space, $Y$ a separable metric space and $\varphi:E\to Y$. Suppose that for every open
$U\subseteq Y$ there are $L\in\Fs(E)$ and a countable $N\subseteq E$ with
$\varphi^{-1}(U)\subseteq L\subseteq\varphi^{-1}(\overline U)\cup N$. Then there is a countable $M\subseteq E$
such that $\varphi|_{E\setminus M}$ is of class $\Fs$.
\end{lemma}
\begin{proof}
Let $\mathcal B$ be a countable base, choose $L(B),N(B)$ for $B\in\mathcal B$, put $M=\bigcup_BN(B)$,
$E_0=E\setminus M$, $\varphi_0=\varphi|_{E_0}$. For open $W$ let $\mathcal B(W)=\{B:\overline B\subseteq W\}$,
so $W=\bigcup_{\mathcal B(W)}B$ (proof of Lemma~\ref{lem:baire-criterion}). Then
\[
\varphi_0^{-1}(W)\subseteq E_0\cap\bigcup_{\mathcal B(W)}L(B)\subseteq
E_0\cap\bigcup_{\mathcal B(W)}\bigl(\varphi^{-1}(\overline B)\cup N(B)\bigr)=E_0\cap\bigcup_{\mathcal B(W)}\varphi^{-1}(\overline B)
\subseteq\varphi_0^{-1}(W),
\]
so $\varphi_0^{-1}(W)=E_0\cap\bigcup_{\mathcal B(W)}L(B)\in\Fs(E_0)$.
\end{proof}

\begin{theorem}[{\cite[Theorem 5]{Diss}}]\label{thm:bf-fsigma}
\lean{TJK.BoundaryFn.exists_countable_isFσMeasurable}
\uses{lem:th-lemma15, lem:two-sided, lem:crosscut-H, fact:tietze-arc, lem:subarc, lem:open-components,
  lem:borel-basic, def:ccset}
Let $Y$ be a separable metric space, $f:\Hp\to Y$ continuous, $E\subseteq\R$, and $\varphi:E\to Y$ a
boundary function for $f$. There is a countable $M\subseteq E$ such that $\varphi|_{E\setminus M}$ is of
class $\Fs$.
\end{theorem}
\begin{proof}
By Lemma~\ref{lem:th-lemma15} it suffices, given open $U\subseteq Y$, to find $L\in\Fs(E)$ and countable
$N$ with $\varphi^{-1}(U)\subseteq L\subseteq\varphi^{-1}(\overline U)\cup N$. Let $W=Y\setminus\overline U$ and
\begin{align*}
E_n&=\{x\in\R:\text{some arc $\gamma$ at $x$ has its other endpoint on $X_n$ and }\gamma\setminus\{x\}\subseteq f^{-1}(U)\},\\
K&=\{x\in\R:\text{some arc $\gamma$ at $x$ has }\gamma\setminus\{x\}\subseteq f^{-1}(W)\}.
\end{align*}
If $x\in\varphi^{-1}(U)$, a subarc of the arc at $x$ along which $f\to\varphi(x)$ is mapped into $U$
(Lemma~\ref{lem:subarc}); cutting it at its first point on a line $X_n$ below its other endpoint shows
$x\in E_n$. So $\varphi^{-1}(U)\subseteq\bigcup_nE_n$, and likewise $\varphi^{-1}(W)\subseteq K$.

\emph{$K\cap\overline{E_n}$ is countable.} Fix $n$. For $x\in K$ choose (Lemma~\ref{lem:subarc}) an arc
$\gamma_x$ at $x$ with $\gamma_x\setminus\{x\}\subseteq\Hp_n\cap f^{-1}(W)$, and let $U_x$ be the component of
the open set $\Hp_n\cap f^{-1}(W)$ containing $\gamma_x\setminus\{x\}$. Let $T$ be the set of points of
$K\cap\overline{E_n}$ that are two-sided limit points of $E_n$. We claim $x\mapsto U_x$ is injective on $T$,
so $T$ is countable (Lemma~\ref{lem:open-components}). Let $x<y$ in $T$ with $U_x=U_y$. Joining
$\gamma_x$ and $\gamma_y$ by a path in $U_x$ and applying Fact~\ref{fact:tietze-arc}, we get a crosscut $\alpha'$
of $\Hp$ from $x$ to $y$ with $\alpha'\setminus\{x,y\}\subseteq U_x\subseteq\Hp_n\cap f^{-1}(W)$. Let $A$ be the
unbounded component of $\Hp\setminus\alpha'$ (Lemma~\ref{lem:crosscut-H}); since $\max_{\alpha'}\operatorname{Im}<1/n$,
$X_n\subseteq A$. Choose $w\in(x,y)\cap E_n$ and an arc $\beta$ at $w$ ending on $X_n$ with
$\beta\setminus\{w\}\subseteq f^{-1}(U)$. Then $\beta\setminus\{w\}$ is connected, misses $\alpha'$ (as
$f^{-1}(U)\cap f^{-1}(W)=\emptyset$) and meets $X_n\subseteq A$, so it lies in $A$; hence
$w\in\overline A\cap\R=\R\setminus(x,y)$, a contradiction. The points of $\overline{E_n}$ that are not two-sided
limit points of $E_n$ form a countable set (Lemma~\ref{lem:two-sided}). Hence $K\cap\overline{E_n}$ is countable.

Put $N=E\cap K\cap\bigcup_n\overline{E_n}$ (countable) and $L=E\cap\bigcup_n\overline{E_n}\in\Fs(E)$. Then
$\varphi^{-1}(U)\subseteq L\subseteq N\cup(E\setminus K)\subseteq N\cup(E\setminus\varphi^{-1}(W))=N\cup\varphi^{-1}(\overline U)$.
\end{proof}

\begin{corollary}[{\cite[Theorem 5, Corollary]{Diss}}; {\cite[Theorems 2, 3]{JMM}}]\label{cor:bf-honorary}
\lean{TJK.BoundaryFn.isHonoraryBaireTwo_upperHalfPlane, TJK.BoundaryFn.isHonoraryBaireTwo_disc_real, TJK.BoundaryFn.isHonoraryBaireTwo_disc_sphere}
\uses{thm:bf-fsigma, thm:honorary-char, lem:transfer, lem:charts, lem:borel-basic}
Let $Y$ be $\R^m$ or $\Sph$. If $f:\Hp\to Y$ (resp.\ $f:\D\to Y$) is continuous, $E\subseteq\R$ (resp.\
$E\subseteq C$) and $\varphi:E\to Y$ is a boundary function for $f$, then $\varphi$ is of honorary Baire class
$\le2$. In particular (\cite[Theorems 2, 3]{JMM}) a finite-valued boundary function on $C$ of a continuous
real function in $\D$ is of honorary Baire class $2(C,\R)$, and a boundary function on $C$ of a
continuous map $\D\to\Sph$ is of honorary Baire class $2(C,\Sph)$.
\end{corollary}
\begin{proof}
For $\Hp$ combine Theorem~\ref{thm:bf-fsigma} with Theorem~\ref{thm:honorary-char}(2)$\Rightarrow$(1). For
$\D$, Lemma~\ref{lem:transfer} with $\kappa$ shows that $\varphi\circ\kappa$ is a boundary function for
$f\circ\kappa$ on $\kappa^{-1}(E\setminus\{1\})$; so $\varphi|_{E\setminus(\kappa(M)\cup\{1\})}$ is of class $\Fs$ for
a countable $M$, and Theorem~\ref{thm:honorary-char} applies.
\end{proof}

\begin{fact}[Bagemihl's ambiguous point theorem {\cite{Bag}}]\label{fact:bagemihl}
\lean{TJK.bagemihl_ambiguous_countable}
\uses{def:ccset}
\source{Bagemihl, Proc. Nat. Acad. Sci. 41 (1955)}
For any function $f:\D\to\Sph$, the set of $x\in C$ at which there are two arcs along which $f$ has
different limits is countable.
\end{fact}

\begin{proposition}[{\cite[remark after Corollary 4]{MMJ}}]\label{prop:gdelta-remark}
\lean{TJK.McMillan.isFσδ_limit_mem_of_isGδ}
\uses{thm:mcmillan, cor:bf-honorary, fact:bagemihl, lem:borel-basic, lem:equiv, lem:baire1-fsigma, def:fn-classes}
In Corollary~\ref{cor:mm-corollaries}(1) it suffices that $E$ be a $\Gd$ subset of $\Sph$.
\end{proposition}
\begin{proof}
$A=\CC(f)$ is $\Fsd$ (Theorem~\ref{thm:mcmillan}). Choose any boundary function $\varphi:A\to\Sph$. By
Corollary~\ref{cor:bf-honorary}, $\varphi=\psi$ off a countable $N_0$, with $\psi\in\mathcal B_1(A,\Sph)$. Off the
countable set $Z$ of Fact~\ref{fact:bagemihl}, every arc limit of $f$ at $x\in A$ equals $\varphi(x)$; so the set
in question differs from $\psi^{-1}(E)$ by a countable set. Writing $E=\bigcap_jO_j$ with $O_j$ open,
$\psi^{-1}(E)=\bigcap_j\psi^{-1}(O_j)$ is $\Fsd$ in $A$, hence in $C$; adding or removing countably many points
keeps a set $\Fsd$.
\end{proof}

\subsection{Homeomorphisms}
The following is Theorem~1 of \cite{JMM}, of which we have only the statement.

\begin{lemma}\label{lem:homeo-values}
\lean{TJK.Homeo.boundaryFn_mem_circle}
\uses{def:ccset}
Let $f$ be a homeomorphism of $\D$ onto $\D$, regarded as a map $\D\to\C$. If $f\to q\in\C$ along an arc at
$x\in C$, then $q\in C$.
\end{lemma}
\begin{proof}
If $q\in\D$ then $z=f^{-1}(f(z))\to f^{-1}(q)\in\D$ as $z\to x$ along the arc, which is absurd.
\end{proof}

\begin{theorem}[{\cite[Theorem 1]{JMM}}]\label{thm:homeo}
\lean{TJK.Homeo.exists_countable_continuousOn, TJK.Homeo.exists_countable_continuousOn', TJK.jumpSet,
  TJK.Homeo.card_mul_le_of_subset_jumpSet, TJK.Homeo.finite_jumpSet, TJK.Homeo.ncard_jumpSet_mul_le,
  TJK.exists_crosscut_near, TJK.exists_component_subset_closedBall, TJK.Homeo.exists_image_crosscut,
  TJK.exists_boundaryArc_of_crosscut, TJK.Homeo.exists_boundaryArc}
\uses{lem:homeo-values, fact:crosscut, lem:disjoint-arcs, lem:disc-geometry, fact:tietze-arc, lem:subarc, def:ccset}
Let $f$ be a homeomorphism of $\D$ onto $\D$, regarded as a map $\D\to\C$, let $E\subseteq C$, and let $\varphi:E\to\C$ be a
boundary function for $f$ (so $\varphi(E)\subseteq C$ by Lemma~\ref{lem:homeo-values}). Then there is a countable $N\subseteq E$ such that $\varphi|_{E\setminus N}$ is continuous. More precisely,
for every $\eps>0$ the set
\[
S_\eps=\{x\in E:\ \text{for every }\delta>0\text{ there is }x'\in E\text{ with }0<\abs{x'-x}<\delta,\
\abs{\varphi(x')-\varphi(x)}\ge\eps\}
\]
has at most $2\pi/\eps$ points.
\end{theorem}
\begin{proof}
Since $\varphi$ is continuous (relative to $E$) at every point of $E\setminus\bigcup_kS_{1/k}$, the first claim
follows from the second. Let $x_1,\dots,x_m$ be distinct points of $S_\eps$. Choose $\delta\in(0,1)$ so small
that the balls $S(\delta,x_i)$ are pairwise disjoint, and for each $i$ choose $x_i'\in E$ with
$0<\abs{x_i'-x_i}<\delta/2$ and $\abs{\varphi(x_i')-\varphi(x_i)}\ge\eps$. Fix $i$, and let
$\gamma,\gamma'$ be arcs at $x_i,x_i'$ along which $f$ tends to $\varphi(x_i),\varphi(x_i')$.

\emph{A crosscut near $x_i$.} By Lemma~\ref{lem:subarc} choose subarcs $\beta\subseteq\gamma\cap S(\delta/2,x_i)$ at
$x_i$ and $\beta'\subseteq\gamma'\cap S(\delta/2,x_i')$ at $x_i'$, with other endpoints $u,v\in\D$. The segment
$[u,v]$ lies in the convex set $\D\cap S(\delta,x_i)$ (Lemma~\ref{lem:disc-geometry}(a)). By
Fact~\ref{fact:tietze-arc}, $\beta\cup[u,v]\cup\beta'$ contains a simple arc $\sigma_i$ from $x_i$ to $x_i'$,
which is a crosscut of $\D$ lying in a closed ball $\overline{S(\delta',x_i)}$ with $\delta'<\delta$. Near $x_i$,
$\sigma_i$ consists of points of $\beta$ (the compact sets $[u,v]$ and $\beta'$ miss $x_i$), so $f\to\varphi(x_i)$
along $\sigma_i$ at $x_i$; likewise $f\to\varphi(x_i')$ at $x_i'$. Since $\varphi(x_i)\ne\varphi(x_i')$ lie on $C$
(Lemma~\ref{lem:homeo-values}), $\tau_i=f(\sigma_i\cap\D)\cup\{\varphi(x_i),\varphi(x_i')\}$ is a crosscut of $\D$.

\emph{The small component.} The connected set $\D\setminus\overline{S(\delta',x_i)}$
(Lemma~\ref{lem:disc-geometry}(b)) misses $\sigma_i$, so it lies in one of the two components of $\D\setminus\sigma_i$
(Fact~\ref{fact:crosscut}); let $R_i$ be the other one. Then $R_i\subseteq\overline{S(\delta',x_i)}\subseteq S(\delta,x_i)$,
so $R_1,\dots,R_m$ are pairwise disjoint.

\emph{Images.} $f$ maps $\D\setminus\sigma_i$ homeomorphically onto $\D\setminus\tau_i$, so $f(R_i)$ is a component
of $\D\setminus\tau_i$, and by Fact~\ref{fact:crosscut} $\overline{f(R_i)}\cap C$ is a closed arc $A_i$ with
endpoints $\varphi(x_i),\varphi(x_i')$, of length at least $\abs{\varphi(x_i)-\varphi(x_i')}\ge\eps$
(Lemma~\ref{lem:disc-geometry}(c)). The sets $f(R_i)$ are pairwise disjoint, so the arcs $A_i$ have pairwise
disjoint interiors (Lemma~\ref{lem:disjoint-arcs}). Hence $m\eps\le2\pi$ by Lemma~\ref{lem:disc-geometry}(d).
\end{proof}

% ===========================================================================
\section{Realizing sets of curvilinear convergence and boundary functions}\label{sec:construct}
% ===========================================================================
Sources: \cite[(A)--(F)]{TAMSb}; \cite[\S\S4--5, Lemmas 9--17, Theorems 4, 6, 7]{Diss}. In this
section $X=(-1,1)$ (identified with $(-1,1)\times\{0\}$), $Q=[-1,1]\times(0,1]$, $Q^\circ=(-1,1)\times(0,1)$,
and $\CC$ refers to arcs in $Q^\circ$ unless said otherwise. Points of $\Hp$ are $\langle x,y\rangle$, $y>0$.

\subsection{Regions}

\begin{definition}[Stolz angles, trapezoids, $B(K,\eps,\alpha,\beta)$]\label{def:regions}
\lean{TJK.Construct.stolz, TJK.Construct.trap, TJK.Construct.tri, TJK.Construct.hull, TJK.Construct.region}
For $a\in\R$, $\eps>0$, $0<\theta<\pi/2$:
$s(a,\eps,\theta)=\{\langle x,y\rangle:0<y<\eps,\ \abs{x-a}<y\cot\theta\}$.
For a nonempty interval $J$ with endpoints $a\le b$: $\Trap(J,\eps,\theta)=\{0<y<\eps,\ a+y\cot\theta<x<b-y\cot\theta\}$
and $\Tri(J,\theta)=\{y\ge0,\ a+y\cot\theta\le x\le b-y\cot\theta\}$. For a nonempty compact $K\subseteq\R$ let
$J(K)=[\min K,\max K]$ and let $\mathcal J(K)$ be the set of components of $J(K)\setminus K$; for
$0<\beta<\alpha<\pi/2$ put
\[
B(K,\eps,\alpha,\beta)=\Trap(J(K),\eps,\alpha)\setminus\bigcup_{I\in\mathcal J(K)}\Tri(I,\beta).
\]
\end{definition}

\begin{lemma}[description of $B$]\label{lem:B-char}
\lean{TJK.Construct.mem_region_iff}
\uses{def:regions}
$\langle x,y\rangle\in B(K,\eps,\alpha,\beta)$ iff $0<y<\eps$, $\min K+y\cot\alpha<x<\max K-y\cot\alpha$, and
$\dist(x,K)<y\cot\beta$.
\end{lemma}
\begin{proof}
Let $y>0$ and $x\in J(K)$. If $x\in K$, $x$ lies in no $\Tri(I,\beta)$ (these lie over $\overline I$, and at an
endpoint $a$ of $I$ we would need $a+y\cot\beta\le a$). If $x\in I=(a,b)\in\mathcal J(K)$, then $x\notin\Tri(I,\beta)$
iff $\min(x-a,b-x)<y\cot\beta$, i.e.\ iff $\dist(x,K)<y\cot\beta$, and $x$ lies over no other component.
\end{proof}

\begin{lemma}[{\cite[(VI)--(XVIII)]{TAMSb}}; {\cite[(6)--(17)]{Diss}}]\label{lem:region-geometry}
\lean{TJK.Construct.isOpen_stolz, TJK.Construct.isOpen_region, TJK.Construct.closure_stolz_inter_axis,
  TJK.Construct.closure_region_inter_axis, TJK.Construct.closure_stolz_subset, TJK.Construct.stolz_disjoint,
  TJK.Construct.region_disjoint_of_split, TJK.Construct.region_disjoint_of_disjoint,
  TJK.Construct.closure_region_subset, TJK.Construct.region_stolz_disjoint_of_not_mem_interior,
  TJK.Construct.stolz_region_disjoint_small_stolz, TJK.Construct.stolz_region_disjoint_small_region,
  TJK.Construct.closure_stolz_subset_region}
\uses{lem:B-char, def:special-family}
Let $K,K_1,K_2$ be nonempty compact subsets of $\R$, $0<\beta<\alpha<\pi/2$, $0<\theta<\pi/2$. Write
$c_\theta=\cot\theta$.
\begin{enumerate}[label=(R\arabic*)]
\item $s(a,\eps,\theta)$ and $B(K,\eps,\alpha,\beta)$ are open subsets of $\Hp$.
\item $\cl s(a,\eps,\theta)\cap\R=\{a\}$ and $\cl B(K,\eps,\alpha,\beta)\cap\R\subseteq K$.
\item If $\eps<\eps'$ and $\theta'<\theta$ then $\cl s(a,\eps,\theta)\cap\Hp\subseteq s(a,\eps',\theta')$.
\item If $a\ne b$ and $\eta\le\abs{a-b}/(2c_\theta)$, then $s(a,\eps,\theta)\cap s(b,\eta,\theta)=\emptyset$ for all $\eps$.
\item If $K_1,K_2$ split (Definition~\ref{def:special-family}), $B(K_1,\eps_1,\alpha,\beta)\cap B(K_2,\eps_2,\alpha,\beta)=\emptyset$.
\item If $K_1\cap K_2=\emptyset$ and $\eta\le\dist(K_1,K_2)/(2c_\beta)$, then
  $B(K_1,\eps,\alpha,\beta)\cap B(K_2,\eta,\alpha,\beta)=\emptyset$ for all $\eps$.
\item If $K_1\subseteq K$, $\eps>\eps_1>0$ and $0<\beta<\beta_1<\alpha_1<\alpha<\pi/2$, then
  $\cl B(K_1,\eps_1,\alpha_1,\beta_1)\cap\Hp\subseteq B(K,\eps,\alpha,\beta)$.
\item If $\alpha<\theta$ and $a\notin J(K)^*$, then $B(K,\eps,\alpha,\beta)\cap s(a,\eps_1,\theta)=\emptyset$.
\item If $a\notin K$ and $\eta\le\dist(a,K)/(c_\theta+c_\beta)$, then $s(a,\eta,\theta)\cap B(K,\eps,\alpha,\beta)=\emptyset$
  and $s(a,\eps,\theta)\cap B(K,\eta,\alpha,\beta)=\emptyset$ for all $\eps$.
\item If $a\in K\cap J(K)^*$, $\beta<\theta$, and $0<\eta<\min\{\eps,(a-\min K)/(c_\theta+c_\alpha),(\max K-a)/(c_\theta+c_\alpha)\}$,
  then $\cl s(a,\eta,\theta)\cap\Hp\subseteq B(K,\eps,\alpha,\beta)$.
\end{enumerate}
\end{lemma}
\begin{proof}
All follow from Lemma~\ref{lem:B-char} and the triangle inequality.
(R1) Strict inequalities between continuous functions. (R2) If $\langle x_n,y_n\rangle\to\langle x,0\rangle$ with
$\dist(x_n,K)<y_nc_\beta$ then $\dist(x,K)=0$. (R3) On the closure, $y\le\eps<\eps'$ and $\abs{x-a}\le yc_\theta<yc_{\theta'}$.
(R4) A common point has $\abs{a-b}<2yc_\theta<2\eta c_\theta$. (R5) With splitting point $p$ and $K_1\le p\le K_2$,
points of the first region have $x<\max K_1\le p$ and those of the second $x>\min K_2\ge p$. (R6) A common point
has $y<\eta$ and $\dist(K_1,K_2)\le\dist(x,K_1)+\dist(x,K_2)<2yc_\beta$. (R7) On the closure $y\le\eps_1<\eps$,
$\min K+yc_\alpha<\min K_1+yc_{\alpha_1}\le x$ (similarly on the right), and $\dist(x,K)\le\dist(x,K_1)\le yc_{\beta_1}<yc_\beta$.
(R8) If $a\le\min K$, points of the Stolz angle have $x<a+yc_\theta<\min K+yc_\alpha$; symmetrically if $a\ge\max K$. (R9) A common point has
$\dist(a,K)\le\abs{x-a}+\dist(x,K)<y(c_\theta+c_\beta)$ with $y<\eta$. (R10) On the closure, $0<y\le\eta<\eps$,
$\abs{x-a}\le yc_\theta$, so $\dist(x,K)\le yc_\theta<yc_\beta$ and $x\ge a-yc_\theta>\min K+yc_\alpha$.
\end{proof}

\subsection{Special families and special functions}

\begin{definition}[splitting, special families]\label{def:special-family}
\lean{TJK.Construct.Split, TJK.Construct.IsSpecialFamily, TJK.Construct.endpoints, TJK.Construct.twoPoint}
Two sets $E_1,E_2\subseteq\R$ \emph{split} if there is $p$ with $E_1\le p\le E_2$ or $E_2\le p\le E_1$ (pointwise). A
\emph{special family} is a nonempty countable family $\mathcal F$ of compact subsets of $X$ such that any two members
are equal, disjoint, or split. Let $\mathcal F^2$ be the members with at least two points and $E(\mathcal F)$ the set of
endpoints of the intervals $J(F)$, $F\in\mathcal F$, $F\ne\emptyset$.
\end{definition}
Two distinct members of a special family meet in at most one point, and such a common point is an endpoint of both
hulls. Endpoints of $J(F)$ belong to $F$.

\begin{lemma}[{\cite[(III)]{TAMSb}}; {\cite[Lemma 9]{Diss}}]\label{lem:special-from-fsigma}
\lean{TJK.Construct.exists_specialFamily_of_isFσ}
\uses{def:special-family, lem:borel-basic}
Every $\Fs$ subset $E$ of $X$ is the union of a special family.
\end{lemma}
\begin{proof}
Write $E=\bigcup_nA_n$ with $A_n$ closed in $X$, $A_1=\emptyset$, $A_n\subseteq A_{n+1}$. Every open interval
$I\subseteq X$ is the union of countably many compact intervals any two of which split: take a strictly increasing
bi-infinite sequence $(c_k)_{k\in\Z}$ in $I$ tending to the endpoints and the intervals $[c_k,c_{k+1}]$. As
$X\setminus A_n$ is a disjoint union of countably many open intervals, $X\setminus A_n=\bigcup_jI_{n,j}$ with compact
intervals $I_{n,j}$ any two of which split. Let $\mathcal F=\{A_1\}\cup\{I_{n,j}\cap A_{n+1}\}$. Its members are compact,
and $\bigcup\mathcal F=\bigcup_n(A_{n+1}\setminus A_n)=E$. For distinct members $I_{n,j}\cap A_{n+1}$,
$I_{n',j'}\cap A_{n'+1}$: if $n<n'$ then $n+1\le n'$ and the second lies in $X\setminus A_{n'}\subseteq X\setminus
A_{n+1}$, so they are disjoint; if $n=n'$ then $j\ne j'$ and they split because $I_{n,j},I_{n,j'}$ do.
\end{proof}

\begin{lemma}[{\cite[(IV)]{TAMSb}}]\label{lem:special-sequence}
\lean{TJK.Construct.exists_specialFamily_seq}
\uses{lem:special-from-fsigma, lem:borel-basic}
If $A\subseteq X$ is $\Fsd$, there are special families $\mathcal F_n$ with $A=\bigcap_n\bigcup\mathcal F_n$ such that
every member of $\mathcal F_{n+1}$ is contained in a member of $\mathcal F_n$.
\end{lemma}
\begin{proof}
Write $A=\bigcap_nA_n$ with $A_n$ $\Fs$ and decreasing, choose special $\mathcal S_n$ with $A_n=\bigcup\mathcal S_n$, and
put $\mathcal F_1=\mathcal S_1$, $\mathcal F_{n+1}=\{E\cap F:E\in\mathcal F_n,F\in\mathcal S_{n+1}\}$. For two members
$E\cap F$, $E'\cap F'$: if $E\ne E'$ compare $E,E'$, and if $E=E'$ compare $F,F'$; subsets of disjoint (splitting) sets are
disjoint (split). So $\mathcal F_{n+1}$ is special, and $\bigcup\mathcal F_n=A_1\cap\dots\cap A_n=A_n$.
\end{proof}

\begin{definition}[special functions]\label{def:special-pair}
\lean{TJK.Construct.IsSpecialPair}
\uses{def:special-family, def:regions}
Let $0<\beta<\alpha<\theta<\pi/2$ and $\mathcal F$ special. A \emph{pair of special $\alpha,\beta,\theta$ functions} for
$\mathcal F$ is a pair $(\eps,\delta)$ of positive functions on $E(\mathcal F)$ and $\mathcal F^2$ such that:
(i) for each $\eta>0$ only finitely many $F\in\mathcal F^2$ have $\delta(F)\ge\eta$;
(ii) for each $\eta>0$ only finitely many $x\in E(\mathcal F)$ have $\eps(x)\ge\eta$;
(iii) the sets $s(x,\eps(x),\theta)$, $x\in E(\mathcal F)$, are pairwise disjoint;
(iv) the sets $B(F,\delta(F),\alpha,\beta)$, $F\in\mathcal F^2$, are pairwise disjoint;
(v) $s(x,\eps(x),\theta)\cap B(F,\delta(F),\alpha,\beta)=\emptyset$ for $x\in E(\mathcal F)$, $F\in\mathcal F^2$.
\end{definition}

\begin{lemma}[{\cite[(XX)]{TAMSb}}; {\cite[Lemma 13]{Diss}}]\label{lem:special-pair-exists}
\lean{TJK.Construct.exists_specialPair}
\uses{def:special-pair, lem:region-geometry}
Every special family has a pair of special $\alpha,\beta,\theta$ functions.
\end{lemma}
\begin{proof}
Enumerate $\mathcal F=\{F_1,F_2,\dots\}$; for $F\in\mathcal F$ let $k(F)$ be its least index, and for $e\in E(\mathcal F)$ let
$k(e)$ be the least $k$ such that $e$ is an endpoint of $J(F_k)$. Let $E(k)$ be the set of endpoints of the $J(F_j)$,
$j\le k$, $F_j\ne\emptyset$. Put $c=\cot\theta+\cot\beta$ and define
\begin{align*}
\delta(F)&=\min\Bigl(\Bigl\{\tfrac1{k(F)}\Bigr\}\cup\Bigl\{\tfrac{\dist(F_j,F)}{2\cot\beta}:j<k(F),\ F_j\in\mathcal F^2,\ F_j\cap F=\emptyset\Bigr\}\cup
\Bigl\{\tfrac{\dist(e,F)}c:e\in E(k(F)-1)\cap J(F)^*\setminus F\Bigr\}\Bigr),\\
\eps(e)&=\min\Bigl(\Bigl\{\tfrac1{k(e)}\Bigr\}\cup\Bigl\{\tfrac{\abs{e-e'}}{2\cot\theta}:e'\in E(k(e))\setminus\{e\}\Bigr\}\cup
\Bigl\{\tfrac{\dist(e,F_j)}c:j\le k(e),\ F_j\in\mathcal F^2,\ e\in J(F_j)^*\setminus F_j\Bigr\}\Bigr).
\end{align*}
These are minima of finite sets of positive numbers. Condition (i): $\delta(F)\ge\eta$ forces $k(F)\le1/\eta$; (ii):
$\eps(e)\ge\eta$ forces $k(e)\le1/\eta$, and each index contributes at most two endpoints.
(iii) If $e\ne e'$, say $k(e')\le k(e)$, then $e'\in E(k(e))$ and (R4) applies.
(iv) If $F\ne F'$ in $\mathcal F^2$ split, use (R5); if they are disjoint, say $k(F')<k(F)$, use (R6).
(v) Let $e\in E(\mathcal F)$ and $F\in\mathcal F^2$. If $e\notin J(F)^*$ use (R8) ($\beta<\alpha<\theta$). Otherwise
$e\notin F$: indeed $e$ is an endpoint of some $J(F')$, $e\in F'$, and $F'\ne F$; if $e\in F$ then $F,F'$ would split at
$e$, making $e$ an endpoint of $J(F)$. If $k(F)\le k(e)$ the definition of $\eps(e)$ and (R9) apply; if $k(F)>k(e)$ then
$e\in E(k(F)-1)$ and the definition of $\delta(F)$ and (R9) apply.
\end{proof}

\begin{lemma}[{\cite[(XXI)]{TAMSb}}; {\cite[Lemma 14]{Diss}}]\label{lem:special-pair-shrink}
\lean{TJK.Construct.IsSpecialPair.mono}
\uses{def:special-pair, def:regions}
If $(\eps,\delta)$ is a special pair and $0<\eps_1\le\eps$, $0<\delta_1\le\delta$, then $(\eps_1,\delta_1)$ is a special pair.
\end{lemma}
\begin{proof}
The regions shrink as the parameters $\eps,\delta$ decrease.
\end{proof}

\subsection{The function $g$}

\begin{theorem}[{\cite[(A)]{TAMSb}}; {\cite[Theorem 4, proof]{Diss}}]\label{thm:construct-A}
\lean{TJK.Construct.exists_g}
\uses{lem:special-sequence, lem:special-pair-exists, lem:special-pair-shrink, lem:region-geometry, lem:subarc,
  def:special-family}
Let $A\subseteq X$ be $\Fsd$. There is a continuous $g:\Hp\to[0,1]$ such that
(a) for each $x\in A$, $g(z)\to0$ as $z\to x$ through $s(x,1,3\pi/8)$; and
(b) if $x\in\R$ and $g\to0$ along some arc at $x$ in $\Hp$, then $x\in A$.
\end{theorem}
\begin{proof}
Take $\mathcal F_n$ as in Lemma~\ref{lem:special-sequence} and $E_n=E(\mathcal F_n)$. Fix sequences
$\beta_n\uparrow\pi/4$ in $(0,\pi/4)$, $\alpha_n\downarrow\pi/4$ in $(\pi/4,\pi/3)$, $\theta_n\uparrow3\pi/8$ in
$(\pi/3,3\pi/8)$, all strictly monotone.

\emph{Claim.} If $x\in E_{n+1}\setminus E_n$, there is a unique $F\in\mathcal F_n^2$ with $x\in F\cap J(F)^*$; and no
$x\in E_n$ lies in $F\cap J(F)^*$ for any $F\in\mathcal F_n^2$. Indeed $x\in K\subseteq F$ for some $K\in\mathcal F_{n+1}$,
$F\in\mathcal F_n$; if $x$ were an endpoint of $J(F)$ it would be in $E_n$; uniqueness and the last assertion follow from
the remark after Definition~\ref{def:special-family}. Likewise each $K\in\mathcal F^2_{n+1}$ lies in a unique
$F\in\mathcal F_n$, and $F\in\mathcal F_n^2$.

\emph{Construction.} Let $(\eps(1,\cdot),\delta(1,\cdot))$ be any special $\alpha_1,\beta_1,\theta_1$ pair for $\mathcal F_1$.
Given $(\eps(n,\cdot),\delta(n,\cdot))$, take any special $\alpha_{n+1},\beta_{n+1},\theta_{n+1}$ pair $(\eps,\delta)$ for
$\mathcal F_{n+1}$ and set: for $x\in E_{n+1}\setminus E_n$, with $F$ from the Claim, $\eps(n+1,x)=\min\{\eps(x),\zeta(x)\}$,
where $\zeta(x)>0$ is given by (R10) for $B(F,\delta(n,F),\alpha_n,\beta_n)$ and $\theta_{n+1}>\beta_n$; for
$x\in E_{n+1}\cap E_n$, $\eps(n+1,x)=\min\{\eps(x),\frac12\eps(n,x)\}$; for $K\in\mathcal F^2_{n+1}$, $K\subseteq F$,
$\delta(n+1,K)=\min\{\delta(K),\frac12\delta(n,F)\}$. By Lemma~\ref{lem:special-pair-shrink} this is a special pair, and by
(R3), (R7), (R10):
\begin{enumerate}[label=(\roman*)]
\item if $x\in E_{n+1}$, $F\in\mathcal F_n^2$, $x\in F\cap J(F)^*$, then
  $\cl s(x,\eps(n+1,x),\theta_{n+1})\cap\Hp\subseteq B(F,\delta(n,F),\alpha_n,\beta_n)$;
\item if $x\in E_{n+1}\cap E_n$, then $\cl s(x,\eps(n+1,x),\theta_{n+1})\cap\Hp\subseteq s(x,\eps(n,x),\theta_n)$;
\item if $K\in\mathcal F^2_{n+1}$, $F\in\mathcal F^2_n$, $K\subseteq F$, then
  $\cl B(K,\delta(n+1,K),\alpha_{n+1},\beta_{n+1})\cap\Hp\subseteq B(F,\delta(n,F),\alpha_n,\beta_n)$.
\end{enumerate}

\emph{The sets $U_n$.} Let $U_n$ be the union of the sets $s(x,\eps(n,x),\theta_n)$, $x\in E_n$, and
$B(F,\delta(n,F),\alpha_n,\beta_n)$, $F\in\mathcal F^2_n$. These pieces are open and pairwise disjoint, so each component
$W$ of $U_n$ lies in one piece, and by (R2)
\begin{equation}\label{eq:A1}
\overline W\cap\R\subseteq\textstyle\bigcup\mathcal F_n .
\end{equation}
Each piece lies in $\{y<\eps(n,x)\}$ or $\{y<\delta(n,F)\}$, so by (i), (ii) of Definition~\ref{def:special-pair} the
pieces form a locally finite family in $\Hp$, and $\cl U_n\cap\Hp$ is the union of the closures of the pieces. By
(i)--(iii) and the Claim, $\cl U_{n+1}\cap\Hp\subseteq U_n$; in particular $U_{n+1}\subseteq U_n$.

\emph{The function.} Let $g_n:\Hp\to[0,1]$ be continuous with $g_n=1$ on $\Hp\setminus U_n$ and $g_n=0$ on
$\cl U_{n+1}\cap\Hp$ (Urysohn), and $g=\sum_n2^{-n}g_n$. Then $g$ is continuous with values in $[0,1]$, and
\begin{equation}\label{eq:A23}
g\ge2^{1-n}\text{ on }\Hp\setminus U_n,\qquad g\le2^{-n}\text{ on }U_{n+1}.
\end{equation}

(a) Let $x_0\in A$ and $n\in\N$. Then $x_0\in F$ for some $F\in\mathcal F_{n+1}$. If $x_0\in E_{n+1}$, put
$\eta=\eps(n+1,x_0)$; as $\theta_{n+1}<3\pi/8$, $s(x_0,\eta,3\pi/8)\subseteq s(x_0,\eps(n+1,x_0),\theta_{n+1})\subseteq U_{n+1}$.
Otherwise $x_0\in F\cap J(F)^*$ with $F\in\mathcal F^2_{n+1}$, and (R10) with $\beta_{n+1}<3\pi/8$ gives $\eta>0$ with
$s(x_0,\eta,3\pi/8)\subseteq B(F,\delta(n+1,F),\alpha_{n+1},\beta_{n+1})\subseteq U_{n+1}$. In both cases $g\le2^{-n}$ on
$s(x_0,1,3\pi/8)\cap\{y<\eta\}$.

(b) If $g\to0$ along an arc $\gamma$ at $x$, then for each $n$ a subarc $\gamma'$ at $x$ has $g<2^{1-n}$ on
$\gamma'\setminus\{x\}$, which therefore lies in $U_n$, hence in one component $W$; so $x\in\overline W\cap\R\subseteq\bigcup\mathcal F_n$
by \eqref{eq:A1}. Thus $x\in\bigcap_n\bigcup\mathcal F_n=A$.
\end{proof}

\begin{theorem}[{\cite[Theorem 4]{Diss}}]\label{thm:cc-realize}
\lean{TJK.Construct.exists_curvConvSet_eq_upperHalfPlane}
\uses{thm:construct-A, lem:transfer, lem:charts, lem:borel-basic, lem:subarc, def:ccset}
For every $\Fsd$ set $A\subseteq\R$ there is a bounded continuous $f:\Hp\to\C$ with $\CC(f)=A$.
\end{theorem}
\begin{proof}
Let $A'=\Psi_\Hp^{-1}(A)\subseteq X$, an $\Fsd$ set, $g$ as in Theorem~\ref{thm:construct-A} for $A'$, and
$f_0=g\sin(1/y)+ig$ on $Q^\circ$. If $x\in A'$, $f_0\to0$ along a short vertical segment at $x$ (which lies in $Q^\circ\cap s(x,1,3\pi/8)$ apart from $x$),
since $\abs{f_0}\le2g$.
If $f_0\to c+di$ along an arc $\gamma$ at $x$, then $g\to d$; if $d\ne0$, then since $y\to0$ continuously along
$\gamma$, every tail of $\gamma$ contains points with $\sin(1/y)=1$ and with $\sin(1/y)=-1$, so $\operatorname{Re}f_0$
has no limit. Hence $d=0$ and $x\in A'$ by Theorem~\ref{thm:construct-A}(b). So $\CC(f_0)=A'$, and
$f=f_0\circ\Psi_\Hp^{-1}$ has $\CC(f)=A$ by Lemma~\ref{lem:transfer}.
\end{proof}

\subsection{The function $h$}

\begin{lemma}[{\cite[(B)(I)]{TAMSb}}; {\cite[Lemma 17]{Diss}}]\label{lem:monotone-inverse}
\lean{TJK.Construct.exists_monotone_leftInverse, TJK.Construct.exists_monotone_leftInverse_real}
(a) If $f:(-1,1)\to\R$ is strictly increasing, there is a continuous weakly increasing $f^*:\R\to[-1,1]$ with
$f^*(f(x))=x$ for all $x$. (b) If $f:\R\to\R$ is strictly increasing and $f(\R)$ is unbounded above and below, there is
a continuous weakly increasing $f^*:\R\to\R$ with $f^*\circ f=\mathrm{id}$.
\end{lemma}
\begin{proof}
(a) $f^*(t)=\sup(\{-1\}\cup\{x:f(x)\le t\})$. It is weakly increasing, and $\{x:f(x)\le f(x_0)\}=(-1,x_0]$ gives
$f^*(f(x_0))=x_0$. Its range contains $(-1,1)$, so it is dense in $[-1,1]$, and a monotone map onto a dense subset of an
interval is continuous. (b) The same with $f^*(t)=\sup\{x:f(x)\le t\}$, whose range is $\R$.
\end{proof}

\begin{lemma}[{\cite[(B)(II)]{TAMSb}}; {\cite[Lemma 16]{Diss}}]\label{lem:closed-graph}
\lean{TJK.Construct.continuous_of_closedGraph}
Let $M$ be a metric space and $u:M\to\R$. Suppose that whenever $p_n\to p$ in $M$ and $u(p_n)\to y$ in
$[-\infty,+\infty]$, we have $y\in\R$ and $y=u(p)$. Then $u$ is continuous.
\end{lemma}
\begin{proof}
If $u(p_n)\not\to u(p)$, a subsequence converges in $[-\infty,\infty]$ to some $y\ne u(p)$.
\end{proof}

\begin{theorem}[{\cite[(B)]{TAMSb}}; {\cite[Theorem 6]{Diss}}]\label{thm:construct-B}
\lean{TJK.Construct.exists_h}
\uses{lem:monotone-inverse, lem:closed-graph, thm:baire1-char, def:regions, def:fn-classes}
Let $E\subseteq X$ and let $\varphi:E\to\R^m$ be of honorary Baire class $\le2$. There is a continuous $h:\Hp\to\R^m$ such
that for each $p\in E$ there is a segment $\lambda_p$ at $p$, contained in $s(p,1,3\pi/8)\cup\{p\}$, along which
$h\to\varphi(p)$. If $\varphi$ is bounded, $h$ can be taken bounded.
\end{theorem}
\begin{proof}
Let $\varphi=\psi$ on $E\setminus N_0$ with $N_0$ countable and $\psi\in\mathcal B_1(E,\R^m)$. Let $N\supseteq N_0$ be a countably
infinite subset of $X$, enumerated injectively as $s_1,s_2,\dots$, and let $w(s_n)=(\sqrt2-1)2^{-n-1}$, so
$\sum_sw(s)<\sqrt2-1=\cot(3\pi/8)$. Define on $X$
\[
f(x)=x+\sum_{s\in N,\ s<x}w(s),
\]
which is strictly increasing with $\abs{f(x)-x}<\sqrt2-1$. Let $f^*$ be as in Lemma~\ref{lem:monotone-inverse}(a).

\emph{The function $u$.} For $\langle x,y\rangle$ with $0<y<1$ the function $v\mapsto v-f^*\bigl((x-(1-y)v)/y\bigr)$ is
continuous and strictly increasing, and lies in $[v-1,v+1]$; so it has a unique zero $u(x,y)\in[-1,1]$. If
$\langle x_n,y_n\rangle\to\langle x,y\rangle$ and $u(x_n,y_n)\to u_0$, then $u_0\in[-1,1]$ and continuity of $f^*$ gives
$u_0=f^*((x-(1-y)u_0)/y)$, so $u_0=u(x,y)$; by Lemma~\ref{lem:closed-graph}, $u$ is continuous.

\emph{The function $h_0$.} By Theorem~\ref{thm:baire1-char} there are continuous $g_n:\R\to\R^m$ with $g_n\to\psi$ on $E$.
For $n\ge2$ and $1/(n+1)\le y\le1/n$ put
\[
h_0(x,y)=\bigl(yn(n+1)-n\bigr)g_n(u(x,y))+\bigl((n+1)-yn(n+1)\bigr)g_{n+1}(u(x,y)),
\]
and $h_0(x,y)=h_0(x,\tfrac12)$ for $y\ge\frac12$. The two formulas agree at $y=1/n$, so $h_0$ is continuous on $\Hp$; the
coefficients lie in $[0,1]$ and add up to $1$.

\emph{The correction.} Let $l_n=f(s_n)=\sup_{x\le s_n}f(x)$, $r_n=\inf_{x>s_n}f(x)=l_n+w(s_n)$, $z_n=\varphi(s_n)-\psi(s_n)$ if
$s_n\in E$ and $z_n=0$ otherwise, $t_n(x,y)=(x-(1-y)s_n)/y$, and
\[
\Lambda_n(x,y)=\bigl[(1-ny)\vee0\bigr]\Bigl[\Bigl(1-\frac{\abs{r_n+l_n-2t_n(x,y)}}{r_n-l_n}\Bigr)\vee0\Bigr]z_n .
\]
$\Lambda_n$ is continuous and vanishes for $y\ge1/n$, and its second factor vanishes unless $l_n<t_n<r_n$. So
$h=h_0+\sum_n\Lambda_n$ is continuous on $\Hp$.

\emph{Points of $E$.} Let $p\in E$ and let $\lambda$ be the part with $0<y\le\frac12$ of the line $x=(f(p)-p)y+p$. On $\lambda$,
$t_n=f(p)+(p-s_n)(1/y-1)$, which is $\ge f(p)\ge r_n$ if $p>s_n$, is $\le f(p)<l_n$ if $p<s_n$, and equals $l_n$ if $p=s_n$;
so every $\Lambda_n$ vanishes on $\lambda$. Also $(x-(1-y)p)/y=f(p)$ there, so $u=p$, and $h_0$ lies on the segment from
$g_n(p)$ to $g_{n+1}(p)$ when $1/(n+1)\le y\le1/n$. Hence $h\to\psi(p)$ along $\lambda$, which is $\varphi(p)$ unless
$p\in N_0$.

\emph{Points of $N_0$.} Let $p=s_m\in E$, $\mu_m=(r_m+l_m)/2$, and let $\lambda'$ be the part with $0<y\le\frac12$ of the line
$x=(\mu_m-s_m)y+s_m$. On $\lambda'$, $t_n=\mu_m+(s_m-s_n)(1/y-1)$. If $s_n<s_m$ then $t_n>l_m\ge r_n$; if $s_n>s_m$ then
$t_n<r_m\le l_n$; so $\Lambda_n=0$ for $n\ne m$, while $\Lambda_m=[(1-my)\vee0]z_m\to z_m$. For $a<s_m<b$ we have
$f(a)\le l_m<\mu_m<r_m\le f(b)$, hence $a\le f^*(\mu_m)\le b$; so $f^*(\mu_m)=s_m$ and $u=s_m$ on $\lambda'$. Thus
$h\to\psi(s_m)+z_m=\varphi(s_m)$ along $\lambda'$.

The slopes satisfy $\abs{f(p)-p}<\sqrt2-1$ and $0\le\mu_m-s_m<\sqrt2-1$, so $\lambda,\lambda'\subseteq s(p,1,3\pi/8)$. If
$\varphi$ is bounded, choose $b$ with $\varphi(E)\subseteq(-b,b)^m$ and clamp each coordinate of $h$ to $[-b,b]$; the limits
along $\lambda_p$ are unchanged.
\end{proof}

\subsection{The function $k$}

\begin{theorem}[{\cite[(C)]{TAMSb}}]\label{thm:construct-C}
\lean{TJK.Construct.exists_k}
Let $d:(0,1]\to(0,\infty)$ be weakly increasing. There is a continuous $k:\R\times(0,1]\to\C$ with $\abs k\le\sqrt2$ such
that for every $a\in(0,1]$ and every connected $\gamma\subseteq[-1,1]\times(0,a]$ containing two points at distance
$\ge d(a)$, there are $p,q\in\gamma$ with $\abs{k(p)-k(q)}\ge2$.
\end{theorem}
\begin{proof}
Let $p(x)=\frac12\int_0^xd(t)\,dt$ for $0<x\le1$: positive, continuous, strictly increasing, with $p(x)\le\frac12d(x)$.
For $a\in(0,1]$ the function $1/p$ is uniformly continuous on $[a/4,1]$, so there is $\eps\in(0,a/4]$ such that
($a/2\le x_1\le1$, $x_2\in(0,1]$, $\abs{x_1-x_2}<\eps$) implies $\abs{p(x_1)^{-1}-p(x_2)^{-1}}\le\frac12$. Let $e(a)$ be the
supremum of such $\eps$ in $(0,1]$; the implication holds with $\eps=e(a)$, and $e$ is weakly increasing because an $\eps$
that works for $a$ works for every $a'\ge a$ (the condition $a'/2\le x_1$ is stronger). Let
$q(x)=\int_0^xe(t)\,dt$, so $q$ is positive, continuous, increasing and $q\le e$; let $m=\min\{p,q\}$. Define
\[
k_1(y)=\sin\frac{2\pi}{y\,m(y)},\qquad k_2(x,y)=\sin\frac{4\pi x}{p(y)},\qquad k=k_1+ik_2 .
\]
Let $z_1,z_2\in\gamma$ with $\abs{z_1-z_2}\ge d(a)$, $z_j=\langle x_j,y_j\rangle$, $y_2\le y_1\le a$, and $a'=\min\{a,2y_1\}$, so
$a'/2\le y_1\le a'\le a$. As $m(a')\le\frac12d(a')\le\frac12d(a)$, either
\[
\text{(12)}\ \ \abs{y_1-y_2}\ge m(a'),\qquad\text{or}\qquad\text{(13)}\ \ \abs{y_1-y_2}<m(a')\text{ and }\abs{x_1-x_2}\ge\tfrac12d(a').
\]
In case (12), using $m(y_2)\le m(y_1)\le m(a')$ and $y_1,y_2\le1$,
\[
\frac{2\pi}{y_2m(y_2)}-\frac{2\pi}{y_1m(y_1)}\ge\frac{2\pi(y_1-y_2)}{y_1y_2m(y_1)}\ge\frac{2\pi m(a')}{y_1y_2m(y_1)}\ge2\pi .
\]
The continuous function $2\pi/(y\,m(y))$ on the connected set $\gamma$ therefore takes all values in an interval of length
$2\pi$, so $k_1$ takes the values $1$ and $-1$ on $\gamma$. In case (13), $\abs{y_1-y_2}<m(a')\le e(a')$ gives
$\abs{p(y_1)^{-1}-p(y_2)^{-1}}\le\frac12$, and $p(y_1)\le p(a')\le\frac12d(a')$, so with $\abs{x_2}\le1$
\[
\Bigl|\frac{4\pi x_1}{p(y_1)}-\frac{4\pi x_2}{p(y_2)}\Bigr|\ge4\pi\Bigl[\frac{\abs{x_1-x_2}}{p(y_1)}-\Bigl|\frac1{p(y_1)}-\frac1{p(y_2)}\Bigr|\Bigr]
\ge4\pi\Bigl[1-\frac12\Bigr]=2\pi,
\]
and $k_2$ takes the values $\pm1$ on $\gamma$. In either case some $p,q\in\gamma$ have $\abs{k(p)-k(q)}\ge2$.
\end{proof}

\subsection{Assembly}

\begin{theorem}[{\cite[(D)]{TAMSb}}]\label{thm:construct-D}
\lean{TJK.Construct.exists_f_Q}
\uses{thm:construct-A, thm:construct-B, thm:construct-C, lem:subarc, def:ccset}
Let $A\subseteq X$ be $\Fsd$ and $\varphi:A\to\C$ bounded and of honorary Baire class $\le2$. There is a bounded continuous
$f:Q\to\C$ such that $A$ is the set of curvilinear convergence of $f$ (arcs in $Q^\circ$) and $\varphi$ is a boundary
function for $f$.
\end{theorem}
\begin{proof}
Let $g$ be given by Theorem~\ref{thm:construct-A} and $h$ (bounded, with $\C=\R^2$) by Theorem~\ref{thm:construct-B} for
$E=A$, both restricted to $Q$. For $t\in(0,1]$ let $d_1(t)$ (resp.\ $d_2(t)$) be the supremum of the $\delta\in(0,1]$ such that
$z_1,z_2\in Q$, $\operatorname{Im}z_j\ge t$, $\abs{z_1-z_2}<\delta$ imply $\abs{h(z_1)-h(z_2)}\le t$ (resp.\
$\abs{g(z_1)-g(z_2)}\le t$); these are positive by uniform continuity on $[-1,1]\times[t,1]$, and weakly increasing. Let
$d(t)=\min\{\frac14d_1(t/2),\frac14d_2(t/2),t/4\}$, $k$ as in Theorem~\ref{thm:construct-C}, and $f=h+gk$.

If $x\in A$, $h\to\varphi(x)$ along a segment in $s(x,1,3\pi/8)$, along which $gk\to0$; a short initial piece of this
segment lies in $Q^\circ$ (Lemma~\ref{lem:subarc}). So $\varphi$ is a boundary function and $A\subseteq\CC(f)$. Conversely let $f$ have a limit along an arc $\gamma$ at $x\in X$ in $Q^\circ$; we show $g\to0$ along
$\gamma$, so $x\in A$ by Theorem~\ref{thm:construct-A}(b). If not, there are $\eps\in(0,1]$ and $z_n=\langle x_n,y_n\rangle\in\gamma\setminus\{x\}$,
$z_n\to x$, with $\abs{g(z_n)}\ge\eps$. Fix $n$ large: $y_n<3\eps/16$, and the other endpoint of $\gamma$ is farther than $y_n$
from $z_n$. Let $a=4y_n/3\le1$. By Lemma~\ref{lem:subarc}(b) a connected $\gamma'\subseteq\gamma\cap\overline{S(d(a),z_n)}$
contains $z_n$ and a point at distance $d(a)$ from it. Since $d(a)\le a/4=y_n/3$, $\gamma'\subseteq[-1,1]\times[a/2,a]$, and
Theorem~\ref{thm:construct-C} gives $p,q\in\gamma'$ with $\abs{k(p)-k(q)}\ge2$. As $\abs{p-q}\le2d(a)<d_1(a/2)$ and
$\abs{p-z_n},\abs{q-z_n}\le d(a)<d_2(a/2)$, all with imaginary parts $\ge a/2$, we get $\abs{h(p)-h(q)}\le a/2$ and
$\abs{g(p)-g(z_n)},\abs{g(q)-g(z_n)}\le a/2<\eps/8$. Hence
\[
\abs{f(p)-f(q)}\ge\abs{g(z_n)}\abs{k(p)-k(q)}-\abs{k(p)}\abs{g(p)-g(z_n)}-\abs{k(q)}\abs{g(q)-g(z_n)}-\abs{h(p)-h(q)}
>2\eps-\tfrac{2\sqrt2+1}8\eps>\eps .
\]
Since $\abs{p-z_n},\abs{q-z_n}\le y_n/3\to0$, $f$ has no limit along $\gamma$, a contradiction.
\end{proof}

\begin{theorem}[{\cite[(E)]{TAMSb}}]\label{thm:construct-E}
\lean{TJK.Construct.exists_f_disc}
\uses{thm:construct-D, lem:transfer, lem:locality, lem:charts, lem:borel-basic, fact:tietze-ext, def:ccset,
  lem:honorary-invariance}
Let $A\subseteq C$ be $\Fsd$ and $\varphi:A\to\C$ bounded and of honorary Baire class $\le2$. There is a bounded continuous
$f:\D\to\C$ with $\CC(f)=A$ such that $\varphi$ is a boundary function for $f$.
\end{theorem}
\begin{proof}
If $A=\emptyset$ take $f(z)=\exp\bigl(i/(1-\abs z)\bigr)$: along any arc, $1/(1-\abs z)$ tends to $+\infty$ continuously, so
$f$ has no limit. Otherwise rotate so that $1\in A$. Let $A'=\Psi_\D^{-1}(A\setminus\{1\})\subseteq X$, which is $\Fsd$, and
apply Theorem~\ref{thm:construct-D} to $A'$ and $\varphi\circ\Psi_\D$ (of honorary Baire class $\le2$ by
Lemma~\ref{lem:honorary-invariance}), obtaining $f_1$; put $f_0=f_1\circ\Psi_\D^{-1}$ on
$G^\circ$. On the set $G''$ of Lemma~\ref{lem:charts}(c), which is closed in $\D$, $f_0$ is continuous and bounded; set
$F=f_0$ on $G''$ and $F=\varphi(1)$ on the radius $[0,1)$, which is closed in $\D$ and disjoint from $G''$. By the Tietze
extension theorem (Fact~\ref{fact:tietze-ext}) $F$ extends to a bounded continuous $f$ on $\D$. Along the radius $f\to\varphi(1)$. Each $x\in
C\setminus\{1\}$ has a neighbourhood $W$ with $W\cap\D\subseteq G''$, so by Lemmas~\ref{lem:locality}
and~\ref{lem:transfer} the arc limits of $f$ at $x$ are those of $f_1$ at $\Psi_\D^{-1}(x)$.
\end{proof}

\begin{theorem}[{\cite[(F)]{TAMSb}}; McMillan's problem]\label{thm:construct-F}
\lean{TJK.Construct.exists_f_sphere}
\uses{thm:construct-E, lem:fold, lem:fold-factor, def:fold, lem:baire1-fsigma, thm:honorary-char, lem:cluster}
Let $A\subseteq C$ be $\Fsd$ and $\varphi:A\to\Sph$ of honorary Baire class $\le2$. There is a continuous $f:\D\to\Sph$ with
$\CC(f)=A$ such that $\varphi$ is a boundary function for $f$.
\end{theorem}
\begin{proof}
Let $\varphi=\psi$ off a countable $N$ with $\psi\in\mathcal B_1(A,\Sph)$, and $A_1=A\setminus N$. By
Lemmas~\ref{lem:baire1-fsigma} and~\ref{lem:fold-factor}, applied in the metric space $A_1$, there is $\psi_0:A_1\to Z_U\cup Z_V$ with
$\Phi\circ\psi_0=\varphi|_{A_1}$ and $\sigma\circ\psi_0$ of class $\Fs$. For $x\in N$ choose $\psi_0(x)\in Z_U\cup Z_V$ with
$\Phi(\psi_0(x))=\varphi(x)$ (possible since $\Phi(Z_U\cup Z_V)=\Sph$). Then $\sigma\circ\psi_0:A\to\R^2$ is bounded
(Lemma~\ref{lem:fold}(v)) and of honorary Baire class $\le2$ (Theorem~\ref{thm:honorary-char}). By
Theorem~\ref{thm:construct-E} there is a bounded continuous $f_1:\D\to\C$ with $\CC(f_1)=A$ and boundary function
$\sigma\circ\psi_0$. Put $f=\Phi\circ\sigma^{-1}\circ f_1$. Then $f\to\Phi(\psi_0(x))=\varphi(x)$ along the arcs at $x\in A$. If
$x\notin A$ and $\gamma$ is an arc at $x$, $f_1$ has no limit along $\gamma$, so its cluster set (with values in a closed disk
$K$ containing $f_1(\D)$) is a connected set with more than one point (Lemma~\ref{lem:cluster}), hence infinite; by
Lemma~\ref{lem:cluster} with $\Theta=\Phi\circ\sigma^{-1}|_K$, $C(f,\gamma)$ is its image under $\Theta$, which is at most
three-to-one, so $C(f,\gamma)$ is infinite and $f$ has no limit
along $\gamma$.
\end{proof}

\begin{theorem}[half-plane versions; {\cite[Theorems 4, 6, 7]{Diss}}]\label{thm:construct-H}
\lean{TJK.Construct.exists_f_upperHalfPlane, TJK.Construct.exists_f_upperHalfPlane_sphere,
  TJK.Construct.exists_boundaryFn_euclidean, TJK.Construct.exists_boundaryFn_sphere}
\uses{thm:construct-D, thm:construct-B, thm:construct-F, lem:transfer, lem:charts, lem:fold, lem:fold-factor,
  thm:honorary-char, lem:cluster, lem:subarc, lem:honorary-invariance}
\begin{enumerate}[label=(\alph*)]
\item If $A\subseteq\R$ is $\Fsd$ and $\varphi:A\to\C$ is bounded and of honorary Baire class $\le2$, there is a bounded continuous
  $f:\Hp\to\C$ with $\CC(f)=A$ and boundary function $\varphi$. The same holds with $\Sph$ in place of $\C$ (and $f$
  continuous).
\item (\cite[Theorem 6]{Diss}) If $E\subseteq\R$ is arbitrary and $\varphi:E\to\R^m$ is of honorary Baire class $\le2$, there is a
  continuous $f:\Hp\to\R^m$ having $\varphi$ as a boundary function.
\item (\cite[Theorem 7]{Diss}) The same with $\Sph$ in place of $\R^m$.
\end{enumerate}
\end{theorem}
\begin{proof}
(a) Transfer Theorem~\ref{thm:construct-D} by $\Psi_\Hp$ (Lemma~\ref{lem:transfer}); no exceptional point occurs. For $\Sph$
argue as in Theorem~\ref{thm:construct-F}. (b) Apply Theorem~\ref{thm:construct-B} to $\Psi_\Hp^{-1}(E)$ and
$\varphi\circ\Psi_\Hp$; the segments have tails in $Q^\circ$; transfer. (c) Write $\varphi=\Phi\circ\sigma^{-1}\circ\chi$ with
$\chi:E\to\R^2$ bounded and of honorary Baire class $\le2$, exactly as in Theorem~\ref{thm:construct-F}, apply (b) to $\chi$,
and compose with $\Phi\circ\sigma^{-1}$.
\end{proof}

% ===========================================================================
\section{Boundary functions of discontinuous functions}\label{sec:discont}
% ===========================================================================
Sources: \cite[Chapter II]{Diss}; \cite[Theorems 4--6]{JMM}. We work in $\Hp$ with boundary $\R$.

\subsection{Baire functions}

\begin{definition}\label{def:meets-in}
\lean{TJK.Discont.MeetsIn}
\uses{def:arcs}
For arcs $\gamma_x,\gamma_y$ at $x,y\in\R$ and $c>0$, $\gamma_x$ \emph{meets $\gamma_y$ below height $c$} if there are
subarcs $\gamma_x'\subseteq\gamma_x$ at $x$ and $\gamma_y'\subseteq\gamma_y$ at $y$, contained in $\{\operatorname{Im}z\le c\}$,
with $\gamma_x'\cap\gamma_y'\ne\emptyset$. This relation is symmetric, and if it holds for $c$ it holds for every $c'\ge c$.
\end{definition}

\begin{lemma}[{the ``Figure 5'' argument of \cite[proof of Theorem 8]{Diss}}]\label{lem:figure5}
\lean{TJK.Discont.figure5}
\uses{lem:crosscut-H, fact:tietze-arc}
Let $x<t<y$ be real, let $\alpha,\beta$ be arcs at $x,y$ in $\Hp$, contained in $\{\operatorname{Im}z\le c\}$, with
$\alpha\cap\beta\ne\emptyset$, and let $\tau$ be an arc at $t$ in $\Hp$ whose other endpoint $p$ has $\operatorname{Im}p=c$.
Then $\tau$ meets $\alpha\cup\beta$.
\end{lemma}
\begin{proof}
$\alpha\cup\beta$ is path connected, so by Fact~\ref{fact:tietze-arc} it contains a simple arc $\sigma$ from $x$ to $y$, a
crosscut of $\Hp$ with $\max_\sigma\operatorname{Im}\le c$. Suppose $\tau\cap\sigma=\emptyset$, and let $V,U$ be the bounded and
unbounded components of $\Hp\setminus\sigma$ (Lemma~\ref{lem:crosscut-H}). The connected set $\tau\setminus\{t\}$ lies in
$U\cup V$, hence in $U$ or in $V$; and $t\in\overline{\tau\setminus\{t\}}$. If it lies in $U$, then
$t\in\overline U\cap\R=\R\setminus(x,y)$, which is false. If it lies in $V$, then $p\in V$, so $p+i\eps\in V$ for small $\eps>0$;
but $p+i\eps\in U$ since it lies above $\sigma$. Contradiction.
\end{proof}

\begin{theorem}[{\cite[Theorem 8]{Diss}}; {\cite[Theorem 4]{JMM}}]\label{thm:baire-plus-one}
\lean{TJK.Discont.baireClass_succ_of_boundaryFn}
\uses{def:meets-in, lem:figure5, lem:strip, lem:subarc, lem:intervals-open, lem:equiv, lem:monotone-extend,
  lem:monotone-compose, lem:baire-precomp, lem:baire-imp-class, lem:additive-basic, lem:baire-criterion}
Let $Y$ be a separable arcwise connected metric space, $\xi\ge1$, $f:\Hp\to Y$ of Baire class $\xi$, $E\subseteq\R$, and
$\varphi:E\to Y$ a boundary function for $f$. Then $\varphi$ is of Baire class $\xi+1$.
\end{theorem}
\begin{proof}
Let $U\subseteq Y$ be open, $V=Y\setminus\overline U$, $A=\varphi^{-1}(U)$, $B=\varphi^{-1}(V)$, $P=A\cup B$. By
Lemma~\ref{lem:baire-criterion} (with $\xi+1\ge2$) it suffices to find $T\in P^{\xi+2}(\R)$ with $A\subseteq T$ and
$T\cap B=\emptyset$: then $\varphi^{-1}(U)\subseteq T\cap E\subseteq\varphi^{-1}(\overline U)$ and $T\cap E\in P^{\xi+2}(E)$.
For $x\in P$ choose (Lemma~\ref{lem:subarc}) an arc $\gamma_x$ at $x$ with $f\to\varphi(x)$ along it,
$\gamma_x\subseteq\{\abs{z-x}\le1\}$, and $f(\gamma_x\setminus\{x\})\subseteq U$ if $x\in A$, $\subseteq V$ if $x\in B$. Then
$\gamma_x\cap\gamma_y=\emptyset$ for $x\in A$, $y\in B$. Write $c_n=1/(n+1)$ and ``meets at $n$'' for ``meets below height
$c_n$''. Let
\[
L=\{x\in P:\forall n\ \exists y\in P\setminus\{x\},\ \gamma_x\text{ meets }\gamma_y\text{ at }n\},\qquad M=P\setminus L,
\]
and let $K_n$ be the set of $x\in P$ such that $\gamma_x$ reaches the line $\operatorname{Im}z=c_n$ and, if $x\in M$, $\gamma_x$
meets no $\gamma_y$, $y\in P\setminus\{x\}$, at $n$. Every $x\in P$ is in some $K_n$, and every $x\in M$ is in $K_n$ for all
large $n$ (an arc at $x$ reaches every small height; and ``meets'' is monotone in the height).

\emph{Step 1: $L\cap A\approx G_a\cap P$ and $L\cap B\approx G_b\cap P$ for $\Gd$ sets $G_a,G_b$.} Let $x\in L\cap A$,
$y\in P\setminus\{x\}$ with $\gamma_x$ meeting $\gamma_y$ at $n$. We claim $[x,y]\cap K_n\subseteq L\cap A$ (with $[x,y]$ the
closed interval between $x$ and $y$). First, $\gamma_x\cap\gamma_y\ne\emptyset$, so $y\in A$. Let $t\in[x,y]\cap K_n$, $t\ne x$.
If $t=y$, $\gamma_t$ meets $\gamma_x$ at $n$. Otherwise let $\tau$ be the subarc of $\gamma_t$ from $t$ to its first point on
$\operatorname{Im}z=c_n$; by Lemma~\ref{lem:figure5}, $\tau$ meets one of the two meeting subarcs, so $\gamma_t$ meets $\gamma_x$
or $\gamma_y$ at $n$. Since $t\in K_n$, this gives $t\notin M$, so $t\in L$; and $\gamma_t$ meets an arc of a point of $A$, so
$t\in A$. Let $W_n$ be the union of the intervals $[x,y]$ over all such pairs $(x,y)$. Then
$L\cap A=\bigcap_n(W_n\cap P)$: ``$\subseteq$'' since $x\in[x,y]$; and if $t\in P\cap\bigcap_nW_n$, take $m$ with $t\in K_m$, then
$t\in L\cap A$ by the claim. By Lemma~\ref{lem:intervals-open} each $W_n\approx O_n$ with $O_n$ open; let $G_a=\bigcap_nO_n$. By
Lemma~\ref{lem:equiv}, $L\cap A\approx G_a\cap P$. Exchanging $A$ and $B$ gives $G_b$.

\emph{Step 2: $M\cap A\approx D\cap M$ with $D\in P^{\xi+2}(\R)$.} For $x\in K_n$ let $p_n(x)$ be the first point of
$\operatorname{Im}z=c_n$ reached along $\gamma_x$ from $x$. If $x<x'$ in $M\cap K_n$, the subarcs from $x$ to $p_n(x)$ and from
$x'$ to $p_n(x')$ are disjoint (they would otherwise meet at $n$), so $\operatorname{Re}p_n(x)<\operatorname{Re}p_n(x')$ by
Lemma~\ref{lem:strip}. Also $\abs{x-\operatorname{Re}p_n(x)}\le1$. By Lemma~\ref{lem:monotone-extend} there is a weakly increasing
$h_n:\R\to\R$ equal to $\operatorname{Re}p_n$ on $M\cap K_n$; let $g_n(x)=f(h_n(x)+ic_n)$. By Lemmas~\ref{lem:baire-precomp}
and~\ref{lem:monotone-compose} there is a countable $N$ such that each $g_n|_{M\setminus N}$ is of Baire class $\xi$. For
$x\in M$ and $n$ large, $g_n(x)=f(p_n(x))$ with $p_n(x)\in\gamma_x\setminus\{x\}$ and $\operatorname{Im}p_n(x)\to0$; as $\gamma_x$
is compact and meets $\R$ only at $x$, $p_n(x)\to x$ along $\gamma_x$, so $g_n(x)\to\varphi(x)$. Hence $\varphi|_{M\setminus N}$ is
of Baire class $\xi+1$, hence of class $P^{\xi+2}$ (Lemma~\ref{lem:baire-imp-class}), so $A\cap(M\setminus N)=D\cap(M\setminus N)$
for some $D\in P^{\xi+2}(\R)$ (Lemma~\ref{lem:additive-basic}(f)).

\emph{Step 3.} From Step 1, $L\approx(G_a\cup G_b)\cap P$, so
$M\cap A\approx D\cap M\approx D\cap\bigl(P\setminus(G_a\cup G_b)\bigr)$ and
\[
A=(L\cap A)\cup(M\cap A)\approx T_0\cap P,\qquad T_0=G_a\cup\bigl(D\setminus(G_a\cup G_b)\bigr)\in P^{\xi+2}(\R),
\]
since $G_a\in Q^2\subseteq P^{\xi+2}$ and $\R\setminus(G_a\cup G_b)$ is $\Fs$. With the countable sets
$C_1=A\setminus T_0$ and $C_2=(T_0\cap P)\setminus A$, the set $T=(T_0\setminus C_2)\cup C_1$ is in $P^{\xi+2}(\R)$ (countable sets
are $\Fs$), contains $A$, and misses $B$.
\end{proof}

\begin{corollary}[{\cite[Theorem 8, Corollary]{Diss}}; {\cite[Theorems 4, 5]{JMM}}]\label{cor:borel}
\lean{TJK.Discont.measurable_of_boundaryFn, TJK.Discont.baireClass_succ_of_boundaryFn_real}
\uses{thm:baire-plus-one, fact:borel-baire, lem:baire-imp-class, lem:additive-basic}
Let $Y$ be a separable, arcwise connected metric space, $f:\Hp\to Y$ Borel measurable, and $\varphi:E\to Y$ a boundary function for $f$.
Then $\varphi$ is Borel measurable. In particular (\cite[Theorems 4, 5]{JMM}), if $f$ is a real function of Baire class
$\xi\ge1$ in the upper half-plane, every finite-valued boundary function of $f$ is of Baire class $\xi+1$, and if $f$ is
Borel measurable so is every finite-valued boundary function.
\end{corollary}
\begin{proof}
By Fact~\ref{fact:borel-baire}, $f$ is of some Baire class $\xi<\omega_1$ (if $\xi=0$, also of class $1$); by
Theorem~\ref{thm:baire-plus-one}, $\varphi$ is of Baire class $\xi+1$, hence Borel measurable
(Lemma~\ref{lem:baire-imp-class}). $\R$ is separable and arcwise connected.
\end{proof}

\subsection{Lebesgue measure}
$X_0=\R\times\{0\}$, $X_1=\R\times\{1\}$; $m$ is planar Lebesgue measure, $m^1$ linear measure on a line, and $m_e$, $m^1_e$
the outer measures. A \emph{segment} below means a closed line segment with one endpoint $\langle a_0(s),0\rangle\in X_0$ and the
other $\langle a_1(s),1\rangle\in X_1$; write $\ell_s(y)=(1-y)a_0(s)+ya_1(s)$.

\begin{definition}[trapezoids]\label{def:trapezoid}
\lean{TJK.Discont.closedTrapezoidE, TJK.Discont.closedTrapezoid, TJK.Discont.IsTrapezoid, TJK.Discont.segTrap}
For $a_0\le b_0$, $a_1\le b_1$ in $[-\infty,+\infty]$ with $a_i<+\infty$ and $b_i>-\infty$ (and $0\cdot(\pm\infty)=0$), let
$T(a_0,b_0;a_1,b_1)=\{\langle x,y\rangle:0\le y\le1,\ x\in\R,\ (1-y)a_0+ya_1\le x\le(1-y)b_0+yb_1\}$, a \emph{closed trapezoid}; $\emptyset$
is also one. $S$ is a \emph{trapezoid} if $S=\emptyset$ or $T^\circ\subseteq S\subseteq T$ for a closed trapezoid $T$, where
$T^\circ=\{\langle x,y\rangle:0\le y\le1,\ (1-y)a_0+ya_1<x<(1-y)b_0+yb_1\}$ is the interior of $T$ \emph{relative to the strip}
$\{0\le y\le1\}$. For segments $s,s'$, $T(s,s')$ is the set of $\langle x,y\rangle$ with $0\le y\le1$ and $x$ between $\ell_s(y)$
and $\ell_{s'}(y)$; for disjoint $s,s'$ with $a_i(s)\le a_i(s')$ this is $T(a_0(s),a_0(s');a_1(s),a_1(s'))$, and $T(s,s)=s$.
(In Lean a segment is given by its pair of abscissas $\langle a_0(s),a_1(s)\rangle$.)
\end{definition}
A trapezoid $S$ is Lebesgue measurable and
$m(S)=\frac12\bigl[(b_0-a_0)+(b_1-a_1)\bigr]=\frac12\bigl[m^1(S\cap X_0)+m^1(S\cap X_1)\bigr]$: integrate the lengths
$(1-y)(b_0-a_0)+y(b_1-a_1)$ of the horizontal sections, and note that $S\cap X_i$ lies between $(a_i,b_i)$ and $[a_i,b_i]$.
(The restrictions $a_i<+\infty$, $b_i>-\infty$ exclude the degenerate sides $a_i=b_i=\pm\infty$, for which the formula fails.)
Only the case of finite parameters and $S=T$ is used below (\texttt{volume\_closedTrapezoid}). Two disjoint segments $s\ne s'$
satisfy either $a_i(s)<a_i(s')$ for $i=0,1$ or $a_i(s)>a_i(s')$ for $i=0,1$ (otherwise $\ell_s-\ell_{s'}$, which is affine, changes sign
on $[0,1]$), hence either $\ell_s<\ell_{s'}$ on $[0,1]$ or $\ell_s>\ell_{s'}$ on $[0,1]$. In particular a set $\Sigma$ of pairwise
disjoint segments is totally ordered by $s\preceq s'\iff a_0(s)\le a_0(s')\text{ and }a_1(s)\le a_1(s')$.

\begin{lemma}[{\cite[Lemma 23]{Diss}}]\label{lem:th-lemma23}
\lean{TJK.Discont.segTrap_subset_union}
\uses{def:trapezoid}
Let $s_1,s_2,s_3,s_4$ be segments, any two equal or disjoint. If $T(s_1,s_2)\cap T(s_3,s_4)\ne\emptyset$ then
$T(s_1,s_3)\subseteq T(s_1,s_2)\cup T(s_3,s_4)$.
\end{lemma}
\begin{proof}
The horizontal section of $T(s_i,s_j)$ at height $y$ is the interval between $\ell_{s_i}(y)$ and $\ell_{s_j}(y)$. The order of the
four numbers $\ell_{s_i}(y)$ (with ties exactly between equal segments) does not depend on $y$. Whether two such intervals meet
depends only on this order; so if the two sections meet at one height they meet at every height. The union of two meeting
intervals is an interval, and it contains $\ell_{s_1}(y)$ and $\ell_{s_3}(y)$.
\end{proof}

\begin{lemma}[{\cite[Lemma 24]{Diss}}]\label{lem:th-lemma24}
\lean{TJK.Discont.isTrapezoid_biUnion}
\uses{def:trapezoid}
If $\Sigma$ is a set of pairwise disjoint segments, $\bigcup_{s,s'\in\Sigma}T(s,s')$ is a trapezoid.
\end{lemma}
\begin{proof}
We may assume $\Sigma\ne\emptyset$. Let $A_i=\inf_sa_i(s)$, $B_i=\sup_sa_i(s)$ (in $[-\infty,+\infty]$) and $T=T(A_0,B_0;A_1,B_1)$.
Since $(1-y)A_0+yA_1\le\ell_s(y)\le(1-y)B_0+yB_1$ for all $s$, the union lies in $T$. Let $\langle x,y\rangle\in T^\circ$. From
$(1-y)A_0+yA_1<x$ we get $s_1,s_2\in\Sigma$ with $(1-y)a_0(s_1)+ya_1(s_2)<x$ (if $y=1$ or $y=0$
only one parameter matters; if $A_0=-\infty$ and $y<1$ take any $s_2$ and $a_0(s_1)$ small; symmetrically if $A_1=-\infty$ and $y>0$;
if both are finite approximate each within $\delta=\frac12[x-(1-y)A_0-yA_1]$). The
$\preceq$-smaller $s$ of $s_1,s_2$ has $\ell_s(y)\le(1-y)a_0(s_1)+ya_1(s_2)<x$. Likewise there is $s'\in\Sigma$ with $x<\ell_{s'}(y)$,
so $\langle x,y\rangle\in T(s,s')$.
\end{proof}

\begin{theorem}[{\cite[Theorem 9]{Diss}}]\label{thm:th-thm9}
\lean{TJK.Discont.outerMeasure_union_segments, TJK.Discont.outerMeasure_union_segments_le, TJK.Discont.le_outerMeasure_union_segments}
\uses{def:trapezoid}
Let $\Sigma$ be a set of pairwise disjoint segments and $S=\bigcup\Sigma$. Then
$m_e(S)=\frac12\bigl[m^1_e(S\cap X_0)+m^1_e(S\cap X_1)\bigr]$.
\end{theorem}
\begin{proof}
Let $E_i=S\cap X_i$, identified with $\{a_i(s):s\in\Sigma\}\subseteq\R$. (The dissertation proves both inequalities at once using
Lemmas~\ref{lem:th-lemma23} and~\ref{lem:th-lemma24}; we give a shorter argument for each.)

\emph{Upper bound.} It suffices to show $m_e(S)\le\frac12[m^1(G_0)+m^1(G_1)]$ for all open $G_i\supseteq E_i$ (outer regularity);
we may assume $m^1(G_i)<\infty$, so every component of $G_i$ is a bounded open interval. Group the segments by the pair $(I,J)$ of
components of $G_0,G_1$ containing $a_0(s),a_1(s)$; there are countably many classes $C$. For a class $C$ let
$\alpha_C,\beta_C$ ($\gamma_C,\delta_C$) be the infimum and supremum of $a_0(s)$ ($a_1(s)$) over $s\in C$. Every $s\in C$ lies in the closed
trapezoid $T_C=T(\alpha_C,\beta_C;\gamma_C,\delta_C)$, of area $\frac12[(\beta_C-\alpha_C)+(\delta_C-\gamma_C)]$, and
$(\alpha_C,\beta_C)\subseteq I$, $(\gamma_C,\delta_C)\subseteq J$. The intervals $(\alpha_C,\beta_C)$ are pairwise disjoint: for classes
with different $I$ this is clear; if $C\ne C'$ have the same $I$ but different (disjoint) $J,J'$ and $x$ lies in both intervals,
take $s_1,s_2\in C$, $s_1',s_2'\in C'$ with $a_0(s_1),a_0(s_1')<x<a_0(s_2),a_0(s_2')$; by the order of disjoint segments,
$a_1(s_1)<a_1(s_2')$ and $a_1(s_1')<a_1(s_2)$, so either $a_1(s_1)\in[a_1(s_1'),a_1(s_2')]\subseteq J'$ or $a_1(s_1')\in[a_1(s_1),a_1(s_2)]\subseteq J$,
contradicting $J\cap J'=\emptyset$. The same holds for $(\gamma_C,\delta_C)$. Hence
$m_e(S)\le\sum_Cm(T_C)=\frac12\bigl[\sum_C(\beta_C-\alpha_C)+\sum_C(\delta_C-\gamma_C)\bigr]\le\frac12[m^1(G_0)+m^1(G_1)]$.

\emph{Lower bound.} First a one-dimensional fact: if $Q\subseteq\R^2$ is totally ordered componentwise, then
$m^1_e(\pi_1Q)+m^1_e(\pi_2Q)\le m^1_e(\{q_1+q_2:q\in Q\})$. Indeed, let $G$ be open containing the right-hand set and group
the $q\in Q$ by the component $I$ of $G$ containing $q_1+q_2$. For $q,q',r,r'$ in one class, with $u=q\wedge r$ and
$v=q'\vee r'\vee u$ (again in the class, as $Q$ is a chain), $(q'_1-q_1)^++(r'_2-r_2)^+\le(v_1+v_2)-(u_1+u_2)\le m^1(I)$; taking
suprema and infima, $m^1_e(\pi_1Q_I)+m^1_e(\pi_2Q_I)\le m^1(I)$, and summing over the countably many classes gives the claim.
Now for $0<y<1$ the section $S^y=\{x:\langle x,y\rangle\in S\}=\{(1-y)a_0(s)+ya_1(s)\}$; applying the fact to
$Q=\{\langle(1-y)a_0(s),ya_1(s)\rangle\}$, which is a chain, gives $m^1_e(S^y)\ge(1-y)m^1_e(E_0)+ym^1_e(E_1)$. Let $M\supseteq S$ be a
measurable hull. By Tonelli,
$m_e(S)=m(M)=\int m^1(M^y)\,dy\ge\int_0^1[(1-y)m^1_e(E_0)+ym^1_e(E_1)]\,dy=\frac12[m^1_e(E_0)+m^1_e(E_1)]$.
\end{proof}

\begin{proposition}[{\cite[example after Theorem 9]{Diss}}]\label{prop:th-example}
\lean{TJK.Discont.segments_not_disjoint_example}
\uses{fact:baire-category}
Disjointness cannot be omitted in Theorem~\ref{thm:th-thm9}: there are sets $M_0\subseteq X_0$, $M_1\subseteq X_1$ of measure zero
such that every point of $\{0<y<1\}$ lies on a segment from $M_0$ to $M_1$.
\end{proposition}
\begin{proof}
Let $M_0=M_1=M$ be a dense $\Gd$ subset of $\R$ of measure zero (e.g.\ the Liouville numbers, or
$\bigcap_n\bigcup_k(q_k-2^{-n-k},q_k+2^{-n-k})$ for an enumeration $q_k$ of $\Q$). Given $\langle x_0,y_0\rangle$ with $0<y_0<1$, the
segment from $\langle a,0\rangle$ to $\langle g(a),1\rangle$, $g(a)=(x_0-(1-y_0)a)/y_0$, passes through $\langle x_0,y_0\rangle$. As $g$ is an affine
bijection, $g^{-1}(M)$ is a dense $\Gd$, so $M\cap g^{-1}(M)\ne\emptyset$ by Baire's theorem; any $a$ in it gives the required segment.
\end{proof}

\begin{lemma}[{\cite[Lemma 25]{Diss}}]\label{lem:th-lemma25}
\lean{TJK.Discont.exists_strictMono_null_range}
For every $\eps>0$ there is a strictly increasing $h:\R\to\R$ with $m^1(h(\R))=0$ and $\abs{x-h(x)}\le\eps$ for all $x$.
\end{lemma}
\begin{proof}
Put $g(t)=\sum_{i\ge0}\lfloor2^it\rfloor3^{-i}$ (for $t\in[0,1)$ this is $\tfrac32$ times the point of the Cantor set whose
ternary digits are twice the binary digits of $t$). Each term is nondecreasing in $t$, and if $t<t'$ then
$\lfloor2^it\rfloor<\lfloor2^it'\rfloor$ once $2^i(t'-t)\ge1$, so $g$ is strictly increasing. Since
$2^it-1<\lfloor2^it\rfloor\le2^it$, we get $3t-\tfrac32\le g(t)\le3t$. If $k=\lfloor2^Nt\rfloor$ and $a=k/2^N$, then
$\lfloor2^it\rfloor=\lfloor2^ia\rfloor$ for $i<N$ and $\lfloor2^{N+j}t\rfloor\le k2^j+2^j=\lfloor2^{N+j}a\rfloor+2^j$, so
$g(a)\le g(t)\le g(a)+\sum_j2^j3^{-N-j}=g(a)+3^{1-N}$. Hence $g([m,m+1))$ ($m\in\Z$) is covered by $2^N$ intervals of
length $3^{1-N}$, so it is null, and $g(\R)=\bigcup_mg([m,m+1))$ is null. Finally put $h(x)=\tfrac\eps3g(x/\eps)$: it is
strictly increasing, $h(\R)=\tfrac\eps3g(\R)$ is null, and $0\le x-h(x)\le\eps/2$.
\end{proof}

\begin{theorem}[{\cite[Theorem 10]{Diss}}]\label{thm:th-thm10}
\lean{TJK.Discont.exists_disjoint_null_arcs}
\uses{lem:th-lemma25, thm:th-thm9, def:arcs}
There is a family $(\gamma_x)_{x\in\R}$ of arcs, $\gamma_x$ an arc at $x$, pairwise disjoint, whose union has planar measure zero.
\end{theorem}
\begin{proof}
(Formalized as \texttt{TJK.Discont.exists\_disjoint\_null\_arcs'}.)
Choose by Lemma~\ref{lem:th-lemma25} strictly increasing $h_n$ ($n\ge0$) with $h_n(\R)$ null and $\abs{x-h_n(x)}\le2^{-n}$. Let
$\rho_k(y)=\max(0,\min(1,2^{k+1}y-1))$, which rises from $0$ at $y=2^{-k-1}$ to $1$ at $y=2^{-k}$, and put
$X_x(y)=x+\sum_{k\ge0}(h_k(x)-h_{k+1}(x))\rho_k(y)$. The terms are bounded by $2\cdot2^{-k}$, so $X_x$ is continuous, and $X_x(0)=x$.
For $2^{-k-1}\le y\le2^{-k}$ we have $\rho_n(y)=0$ for $n<k$ and $\rho_n(y)=1$ for $n>k$, and the series telescopes (as
$h_n(x)\to x$) to $X_x(y)=(1-\rho_k(y))h_{k+1}(x)+\rho_k(y)h_k(x)$: on $[2^{-k-1},2^{-k}]$, $X_x$ is affine from $h_{k+1}(x)$ to $h_k(x)$.
So $\gamma_x=\{\langle X_x(y),y\rangle:0\le y\le1\}$ is an arc at $x$ (injective in $y$), lying in $\Hp$ except for $x$.
If $x<x'$ then $X_x<X_{x'}$ on $[0,1]$, so the arcs are disjoint. Let $S_k$ be the union over $x$ of the segments
$\{\langle X_x(y),y\rangle:2^{-k-1}\le y\le2^{-k}\}$; these are pairwise disjoint. Under the affine map
$\langle u,t\rangle\mapsto\langle u,(1+t)2^{-k-1}\rangle$, $S_k$ is the image of the union of the disjoint segments from
$\langle h_{k+1}(x),0\rangle$ to $\langle h_k(x),1\rangle$, whose end sets lie in the null sets $h_{k+1}(\R)$, $h_k(\R)$; by
Theorem~\ref{thm:th-thm9} that union is null, hence so is $S_k$. Thus $\bigcup_x\gamma_x\subseteq X_0\cup\bigcup_kS_k$ is null.
\end{proof}

\begin{corollary}[{\cite[Theorem 10, Corollary 1]{Diss}}; {\cite[Theorem 6]{JMM}}]\label{cor:ae-zero}
\lean{TJK.Discont.exists_ae_zero_boundaryFn}
\uses{thm:th-thm10, def:ccset}
Let $Y$ be a metric space with a point $0$ and $\varphi:\R\to Y$ arbitrary. There is $f:\Hp\to Y$ with $f=0$ outside a set of
measure zero such that $\varphi$ is a boundary function for $f$.
\end{corollary}
\begin{proof}
Put $f=\varphi(x)$ on $\gamma_x\setminus\{x\}$ and $f=0$ elsewhere; $f$ is constant along each $\gamma_x$.
\end{proof}

\begin{lemma}[non-measurable sets; Vitali]\label{fact:vitali}
\lean{TJK.Discont.exists_not_nullMeasurableSet}
There is a subset of $\R$ that is not Lebesgue measurable.
\end{lemma}
\begin{proof}
Let $V\subseteq[0,1]$ contain one point of each coset of $\mathbb Q$ in $\R$. The translates $V+q$, $q\in\mathbb Q\cap[-1,1]$, are pairwise disjoint, cover $[0,1]$ and lie in $[-1,2]$.
If $V$ were measurable they would all have measure $\lambda(V)$, so $1\le\sum_q\lambda(V)\le 3$ by countable additivity, impossible whether $\lambda(V)=0$ or $\lambda(V)>0$.
\end{proof}

\begin{corollary}[{\cite[Theorem 10, Corollary 2]{Diss}}]\label{cor:nonmeasurable}
\lean{TJK.Discont.exists_measurable_nonmeasurable_boundaryFn}
\uses{cor:ae-zero, fact:vitali}
There is a Lebesgue-measurable $f:\Hp\to\R$ with a boundary function on $\R$ that is not Lebesgue measurable.
\end{corollary}
\begin{proof}
Apply Corollary~\ref{cor:ae-zero} to the indicator function of a non-measurable set; $f=0$ almost everywhere.
\end{proof}

% ===========================================================================
\section{Bounded harmonic functions}\label{sec:harmonic}
% ===========================================================================
Source: \cite{TAMSa}. A function $u$ defined in $\D$ has the \emph{asymptotic value} $a$ at $\zeta\in C$ if $u\to a$ along some arc
at $\zeta$, and the \emph{general limit} $a$ at $\zeta$ if $u(z)\to a$ as $z\to\zeta$, $z\in\D$. $P_r(\theta)=(1-r^2)/(1-2r\cos\theta+r^2)$
is the Poisson kernel, $\Delta(S,T)=\inf\{\abs{s-t+2m\pi}:s\in S,t\in T,m\in\Z\}$ for $S,T\subseteq\R$, and $\expi S=\{e^{it}:t\in S\}$.

\begin{fact}[Poisson integrals {\cite[pp.~32--38]{Hoffman}}]\label{fact:poisson}
\lean{TJK.Harmonic.poisson_bounded_harmonic, TJK.Harmonic.poisson_tendsto_of_continuousAt,
  TJK.Harmonic.hp_radial_limit_poisson}
\source{Hoffman, Banach spaces of analytic functions, pp.~32--38; Mathlib has the Poisson kernel (Analysis/Complex/Poisson)}
(i) If $\varphi$ is a bounded measurable real function on $C$, its Poisson integral
$P[\varphi](re^{i\theta})=\frac1{2\pi}\int_{-\pi}^{\pi}\varphi(e^{it})P_r(\theta-t)\,dt$ is a bounded harmonic function in $\D$ with
$\abs{P[\varphi]}\le\sup\abs\varphi$.
(ii) If $\varphi\in L^1(C)$ is defined on a set $S$ of full measure, $\zeta_0\in S$, and $\varphi|_S$ is continuous at $\zeta_0$, then
$P[\varphi]$ has general limit $\varphi(\zeta_0)$ at $\zeta_0$.
(iii) If $u$ is harmonic in $\D$ and $\{\int_{-\pi}^\pi\abs{u(re^{i\theta})}^pd\theta:0\le r<1\}$ is bounded for some $p>1$, then the
radial limit of $u$ exists at almost every point of $C$, the function $\psi$ so defined almost everywhere is in $L^1(C)$, and
$u=P[\psi]$.
\end{fact}

\begin{fact}[Jordan curve theorem and accessibility]\label{fact:jordan-access}
\lean{TJK.Harmonic.jordan_accessible}
\source{Jordan curve theorem; Newman, Elements, corollary on p.~164}
If $\Gamma\subseteq\C$ is a Jordan curve, $\C\setminus\Gamma$ has exactly two components, a bounded one $B$ and an unbounded one,
each with frontier $\Gamma$, and every point $p\in\Gamma$ is accessible from $B$: there is a simple arc $\gamma$ with endpoint $p$ and
$\gamma\setminus\{p\}\subseteq B$.
\end{fact}

\begin{lemma}[{\cite[Lemma]{TAMSa}}]\label{lem:ha-lemma}
\lean{TJK.Harmonic.accessible_avoiding_domains}
\uses{fact:jordan-access}
Let $(I_n)$ be pairwise disjoint nondegenerate closed arcs of $C$ and $(\Omega_n)$ Jordan domains in $\D$ with
$\overline{\Omega_n}\cap C=I_n$, $\overline{\Omega_n}\cap\overline{\Omega_m}=\emptyset$ for $n\ne m$, and $\diam\Omega_n\to0$. Then every point of
$C\setminus\bigcup_nI_n^*$ is accessible by a simple arc in $\D\setminus\bigcup_n\overline{\Omega_n}$.
\end{lemma}
\begin{proof}
$\partial\Omega_n$ is a Jordan curve containing $I_n$; let $\Gamma_n=\partial\Omega_n\setminus I_n^*$, an arc with the endpoints of $I_n$, and
$\phi_n:I_n\to\Gamma_n$ a homeomorphism fixing the endpoints. Let $\Gamma=(C\setminus\bigcup_nI_n^*)\cup\bigcup_n\Gamma_n$ and define $\phi:C\to\Gamma$
by $\phi=\phi_n$ on $I_n$ and $\phi=\mathrm{id}$ elsewhere; let $\phi_p$ use only $\phi_1,\dots,\phi_p$. Each $\phi_p$ is continuous, and
$\abs{\phi_p-\phi}\le\sup_{n>p}\diam\Omega_n\to0$, so $\phi$ is continuous; it is injective and $C$ is compact, so $\Gamma$ is a Jordan
curve. Each $\Omega_n$ misses $\Gamma$; a small disk about an interior point of $I_n$ misses $\Gamma$ and meets both $\Omega_n$ and
$\C\setminus\overline\D$, which lies in the unbounded component; so $\Omega_n$ lies in the unbounded component. The points of
$C\setminus\Gamma=\bigcup I_n^*$ also lie in the unbounded component (they are limits of points outside $\overline\D$ near which
$\Gamma$ is absent), and $\overline{\Omega_n}\subseteq\Omega_n\cup\Gamma_n\cup I_n$. Hence the bounded component $B$
(Fact~\ref{fact:jordan-access}) lies in $\D\setminus\bigcup_n\overline{\Omega_n}$, and each point of $C\setminus\bigcup I_n^*\subseteq\Gamma$
is accessible from $B$.
\end{proof}

\begin{lemma}[general limits along an arc]\label{lem:general-limit-endpoints}
\lean{TJK.Harmonic.asymptotic_at_endpoint}
\uses{def:ccset}
Let $u:\D\to\R$ have general limit $c$ at every point of an open arc $J$ of $C$. Then $u$ has the asymptotic value $c$ at each endpoint
of $J$.
\end{lemma}
\begin{proof}
Let $J=\expi(t_0,t_1)$ and $\tau_k=t_0+(t_1-t_0)2^{-k-1}$. By compactness of $\expi[\tau_k,\tau_0]$ there is $\rho_k\in(0,1/k)$ with
$\abs{u(re^{it})-c}<1/k$ for $t\in[\tau_k,\tau_0]$ and $1-\rho_k<r<1$; we may take $\rho_k$ decreasing. Let $r(t)$ be continuous on
$(t_0,\tau_0]$, affine on each $[\tau_{k+1},\tau_k]$, with $r(\tau_k)=1-\rho_{k+1}/2$. Then $\{r(t)e^{it}\}\cup\{e^{it_0}\}$ is an arc at
$e^{it_0}$ along which $u\to c$.
\end{proof}

\begin{lemma}[the intervals]\label{lem:ha-sequences}
\lean{TJK.Harmonic.exists_intervals}
Let $(p_n)$ be a countable dense subset of $(-\pi,\pi)$. There are natural numbers $n(1)<n(2)<\cdots$ and reals $\eps_k$ such that, with
$\alpha(k)=p_{n(k)}-\pi\eps_k^2/2$ and $\beta(k)=p_{n(k)}+\pi\eps_k^2/2$:
I.~$0<\eps_k\le2^{-k-1}$; II.~$-\pi<\alpha(k)<\beta(k)<\pi$; III.~the intervals $[\alpha(k),\beta(k)]$ are pairwise disjoint;
IV.~$\eps_k\le\bigl(\Delta([\alpha(h),\beta(h)],[\alpha(k),\beta(k)])/4\bigr)^4$ for $h<k$; V.~$p_h\in\bigcup_{i\le k}[\alpha(i),\beta(i)]$
for $h\le n(k)$.
\end{lemma}
\begin{proof}
Let $n(1)=1$ and choose $\eps_1$ satisfying I, II. Given the first $q$ terms, let $n(q+1)$ be the least $n$ with
$p_n\notin\bigcup_{i\le q}[\alpha(i),\beta(i)]$ (it exists because this finite union of closed intervals of total length $<2\pi$ cannot
contain a dense set; and $n(q+1)>n(q)$ by V), and choose $\eps_{q+1}>0$ with I, II and
$\frac\pi2\eps_{q+1}^2+4\eps_{q+1}^{1/4}\le\frac12\min_{k\le q}\Delta(p_{n(q+1)},[\alpha(k),\beta(k)])$. Then
$\Delta([\alpha(q+1),\beta(q+1)],[\alpha(k),\beta(k)])\ge4\eps_{q+1}^{1/4}>0$, which gives III and IV, and V holds by the minimality
of $n(q+1)$.
\end{proof}

\begin{lemma}[estimates]\label{lem:ha-estimates}
\lean{TJK.Harmonic.omega_jordan, TJK.Harmonic.omega_closure_disjoint, TJK.Harmonic.diam_omega_tendsto,
  TJK.Harmonic.fn_le_eps, TJK.Harmonic.fn_le_eps_sqrt, TJK.Harmonic.fn_tendsto_zero}
\uses{lem:ha-sequences, fact:poisson}
With the data of Lemma~\ref{lem:ha-sequences}, let $I_n=\expi[\alpha(n),\beta(n)]$, $\chi_n$ its indicator, $f_n=P[\chi_n]$, and
$\Omega_n=\{re^{i\theta}:1-\eps_n<r<1,\ \Delta(\theta,[\alpha(n),\beta(n)])<(1-r)^{1/4}\}$. Then $\Omega_n$ is a Jordan domain with
$\overline{\Omega_n}\cap C=I_n$; $\overline{\Omega_n}\cap\overline{\Omega_m}=\emptyset$ for $n\ne m$; $\diam\Omega_n\to0$; and
\begin{enumerate}[label=(\arabic*)]
\item $0\le f_n\le\eps_n$ on $\D\setminus\Omega_n$;
\item $f_n(re^{i\theta})\le\eps_n\sqrt{1-r}$ when $r>1-\eps_n$ and $\Delta(\theta,[\alpha(n),\beta(n)])\ge(1-r)^{1/4}$;
\item $f_n\to0$ along every arc at a point of $C$ that does not meet $\Omega_n$.
\end{enumerate}
\end{lemma}
\begin{proof}
$\Omega_n$ is the image of $(0,1)\times(1-\eps_n,1)$ under $(s,r)\mapsto r\exp\bigl(i[\alpha(n)-(1-r)^{1/4}+s(\beta(n)-\alpha(n)+2(1-r)^{1/4})]\bigr)$,
which extends to a homeomorphism of the closed square onto $\overline{\Omega_n}$ taking the edge $r=1$ onto $I_n$. A common point of
$\overline{\Omega_n}$ and $\overline{\Omega_m}$, $n<m$, would give $\Delta([\alpha(n),\beta(n)],[\alpha(m),\beta(m)])\le2\eps_m^{1/4}$, against IV.
$\diam\Omega_n\le\eps_n+\pi\eps_n^2+2\eps_n^{1/4}$.

For $re^{i\theta}\notin\Omega_n$ either $r\le1-\eps_n$ or ($r>1-\eps_n$ and $\Delta(\theta,[\alpha(n),\beta(n)])\ge(1-r)^{1/4}$). In the first case
$P_r\le(1+r)/(1-r)\le2/\eps_n$ and $f_n\le\frac1{2\pi}\cdot\pi\eps_n^2\cdot\frac2{\eps_n}=\eps_n$. In the second case, with $v=(1-r)^{1/4}<1$, for
$t\in[\alpha(n),\beta(n)]$ we have $\cos(\theta-t)\le\cos v\le1-\frac{v^2}2+\frac{v^4}{24}\le1-\frac14\sqrt{1-r}$; as $r>\frac34$,
\[
f_n\le\frac1{2\pi}\cdot\pi\eps_n^2\cdot\frac{1-r^2}{\frac12r\sqrt{1-r}}\le\frac83\eps_n^2\sqrt{1-r}\le\eps_n\sqrt{1-r}\le\eps_n,
\]
since $\eps_n\le\frac14$. This proves (1) and (2); (3) follows from (2) because $r\to1$ along such an arc.
\end{proof}

\begin{theorem}[{\cite{TAMSa}}; answer to {\cite[Problem 5]{BP}}]\label{thm:ha-example}
\lean{TJK.Harmonic.exists_bounded_harmonic_boundaryFn_not_baireOne}
\uses{lem:ha-estimates, lem:ha-sequences, lem:ha-lemma, lem:general-limit-endpoints, fact:poisson,
  cor:baire-continuity-point, lem:baire-precomp, def:ccset}
There are a bounded harmonic function $f$ in $\D$ and a boundary function $\varphi:C\to\R$ for $f$ that is not of Baire class $\le1$.
\end{theorem}
\begin{proof}
With the notation of Lemma~\ref{lem:ha-estimates}, let $\chi=\sum_n\chi_n$ (the indicator of $\bigcup_nI_n$) and $f=P[\chi]=\sum_nf_n$
(monotone convergence); $f$ is bounded harmonic (Fact~\ref{fact:poisson}(i)).

$\chi$ is continuous with value $1$ on each open arc $I_n^*$, so $f$ has general limit $1$ there (Fact~\ref{fact:poisson}(ii)) and, by
Lemma~\ref{lem:general-limit-endpoints}, asymptotic value $1$ at $e^{i\alpha(n)}$ and $e^{i\beta(n)}$.

Let $\zeta\in C\setminus\bigcup_nI_n^*$. By Lemma~\ref{lem:ha-lemma} there is an arc $\gamma$ at $\zeta$ meeting no $\Omega_n$. For $m\in\N$ and
$z\in\gamma$, by (1) and I, $\abs{f(z)}\le\sum_{n\le m}f_n(z)+\sum_{n>m}\eps_n\le\sum_{n\le m}f_n(z)+2^{-m-1}$, and by (3)
$\limsup\abs{f}\le2^{-m-1}$ along $\gamma$. So $f$ has asymptotic value $0$ at $\zeta$; in particular the endpoints of the $I_n$ are
ambiguous points. Let $A=\{e^{i\alpha(n)}\}$, $B=\{e^{i\beta(n)}\}$, and define $\varphi=1$ on $\bigcup_nI_n^*\cup B$ and $\varphi=0$ on the rest of
$C$. Then $\varphi$ is a boundary function for $f$.

By V, every $p_h$ lies in $\bigcup_nI_n$, so $\bigcup_nI_n$ is dense in $C$. Given $e^{i\alpha(n)}$ and a short open arc $O$ ending at it on
the side away from $I_n$, some $I_m$ meets $O$; since $I_m\cap I_n=\emptyset$, $e^{i\beta(m)}\in O$; and some $I_{m'}$ meets the open arc between
$e^{i\beta(m)}$ and $e^{i\alpha(n)}$, so lies inside it and $e^{i\alpha(m')}\in O$. Arguing likewise at points of $B$, each point of $A\cup B$
is a limit of other points of $A$ and of points of $B$. Hence $P=\cl(A\cup B)$ is a nonempty closed (so complete) set in which both $A$
and $B$ are dense. As $\varphi=0$ on $A$ and $\varphi=1$ on $B$, $\varphi|_P$ is continuous at no point. If $\varphi$ were of Baire class $\le1$, so
would be $\varphi|_P$, contradicting Corollary~\ref{cor:baire-continuity-point}.
\end{proof}

\begin{theorem}[{\cite[Theorem]{TAMSa}}]\label{thm:ha-theorem}
\lean{TJK.Harmonic.isMeagre_asymptotic_diff_generalLimit}
\uses{fact:poisson, cor:bf-honorary, thm:kuratowski, lem:meager-subspace}
Let $u$ be harmonic in $\D$ with $\{\int\abs{u(re^{i\theta})}^pd\theta:0\le r<1\}$ bounded for some $p>1$. Let $A$ be the set of points of $C$
at which $u$ has an asymptotic value and $G$ the set at which $u$ has a general limit. Then $A\setminus G$ is of the first category.
\end{theorem}
\begin{proof}
Let $S$ be the set where the radial limit $\psi$ exists; $S\subseteq A$, $C\setminus S$ is null, and $u=P[\psi]$ (Fact~\ref{fact:poisson}(iii)).
Extend $\psi|_S$ to a boundary function $\varphi_1$ on $A$. By Corollary~\ref{cor:bf-honorary} there is $\varphi\in\mathcal B_1(A,\R)$ with
$\varphi=\varphi_1$ off a countable set; then $\varphi=\psi$ almost everywhere and $u=P[\varphi]$. By Theorem~\ref{thm:kuratowski} there is $K\subseteq A$ with
$A\setminus K$ of the first category in $A$, hence in $C$ (Lemma~\ref{lem:meager-subspace}), such that $\varphi|_A$ is continuous at each point of
$K$. By Fact~\ref{fact:poisson}(ii), $K\subseteq G$; so $A\setminus G\subseteq A\setminus K$.
\end{proof}

\begin{corollary}[{\cite[Corollary]{TAMSa}}]\label{cor:ha-corollary}
\lean{TJK.Harmonic.isMeagre_discontinuity_boundaryFn}
\uses{thm:ha-theorem, def:ccset}
Under the hypothesis of Theorem~\ref{thm:ha-theorem}, if $\varphi$ is a boundary function for $u$ defined on $E\subseteq C$, then the set of
points of $E$ at which $\varphi$ is discontinuous is of the first category in $C$.
\end{corollary}
\begin{proof}
At a point $\zeta\in E\cap G$ with general limit $L$: if $\abs{u-L}<\eps$ on $\D\cap S(\delta,\zeta)$, then every $x'\in E\cap S(\delta,\zeta)$ has
$\abs{\varphi(x')-L}\le\eps$, since $\varphi(x')$ is a limit along an arc whose tail lies in $S(\delta,\zeta)$. So $\varphi$ is continuous on $E\cap G$, and
its discontinuities lie in $A\setminus G$.
\end{proof}

\begin{corollary}[analytic functions; {\cite{TAMSa}}]\label{cor:ha-analytic}
\lean{TJK.Harmonic.isMeagre_asymptotic_diff_generalLimit_holomorphic}
\uses{thm:ha-theorem, cor:ha-corollary}
Theorem~\ref{thm:ha-theorem} and Corollary~\ref{cor:ha-corollary} hold for holomorphic $F$ in $\D$ with bounded $p$-means, $p>1$.
\end{corollary}
\begin{proof}
Apply them to $\operatorname{Re}F$ and $\operatorname{Im}F$, whose $p$-means are bounded by those of $F$: if $F$ has an asymptotic value at $x$ so
do both parts, and if both have general limits so does $F$.
\end{proof}

\begin{remark}
The example of \cite[Theorem 6]{BP} shows that the hypothesis on the integral means cannot be omitted.
\end{remark}

% ===========================================================================
\section{A cube with a non-Borel set of curvilinear convergence}\label{sec:cube}
% ===========================================================================
Source: \cite{PAMS}. For an open connected $\Omega\subseteq\R^n$ and $f:\Omega\to\R^m$, the set of curvilinear convergence $\CC(f)$ is the set of
$p\in\partial\Omega$ for which there is an arc $\gamma$ at $p$ with $f$ converging to a finite limit along $\gamma$.

\begin{definition}[asymptotic values]\label{def:asymptotic}
\lean{TJK.Cube.IsAsymptoticValue, TJK.Cube.IsPointAsymptoticValue}
\uses{def:ccset, cor:path-arc}
$a\in\R^m$ is an \emph{asymptotic value} of $f:\Omega\to\R^m$ if there is a continuous $v:[0,1)\to\Omega$ with $\dist(v(t),\R^n\setminus\Omega)\to0$
and $f(v(t))\to a$ as $t\to1$; it is a \emph{point asymptotic value at $p$} if moreover $v(t)\to p\in\partial\Omega$. By
Corollary~\ref{cor:path-arc}, $p\in\CC(f)$ iff $f$ has a (finite) point asymptotic value at $p$.
\end{definition}

\begin{fact}[analytic non-Borel sets]\label{fact:analytic-nonborel}
\lean{TJK.Cube.exists_analytic_not_borel}
\source{Souslin--Lusin; Sierpi\'nski, General topology p.~254. Mathlib has MeasureTheory.AnalyticSet}
There is an analytic subset of $\R$ that is not a Borel set.
\end{fact}

\begin{fact}[Kierst {\cite{Kierst}}]\label{fact:kierst}
\lean{TJK.Cube.kierst}
\source{Kierst 1936}
For every analytic set $A\subseteq\R$ there is a holomorphic function $h$ in $\D$ omitting the values $i$ and $-i$ whose set of finite asymptotic
values is $A\cup\{-i,i\}$.
\end{fact}

\begin{fact}[Lehto--Virtanen {\cite[p.~53]{LV}}; Bagemihl--Seidel {\cite[Theorem 1]{BS}}]\label{fact:normal}
\lean{TJK.Cube.asymptotic_is_point_asymptotic}
\source{Lehto--Virtanen 1957 (functions omitting two values are normal); Bagemihl--Seidel 1961 (Koebe arcs), as used by McMillan}
If $h$ is holomorphic in $\D$ and omits two finite values, then every finite asymptotic value of $h$ is a point asymptotic value of $h$.
\end{fact}

\begin{lemma}[{\cite[Lemma]{PAMS}}]\label{lem:cube-lemma}
\lean{TJK.Cube.exists_s}
\uses{fact:analytic-nonborel, fact:kierst, fact:normal, def:asymptotic, cor:path-arc}
There is a continuous $s:\Hp\to\C$ with $\abs s<1$ such that the set $E$ of real asymptotic values of $s$ lying in $(-1,1)$ equals the set
of real point asymptotic values of $s$ lying in $(-1,1)$, and $E$ is not a Borel set.
\end{lemma}
\begin{proof}
Let $A$ be analytic and not Borel, $h$ as in Fact~\ref{fact:kierst}, and $\rho(w)=w/(1+\abs w)$, a homeomorphism of $\C$ onto $\D$ mapping
$\R$ onto $(-1,1)$. Put $s(x,y)=\rho\bigl(h((1-y)e^{ix})\bigr)$ for $0<y<1$ and $s(x,y)=\rho(h(0))$ for $y\ge1$. Let $b\in(-1,1)$.
If $b$ is an asymptotic value of $s$ along $v=\langle x(t),y(t)\rangle$ (so $y(t)\to0$), then $\rho^{-1}(b)$ is an asymptotic value of $h$ along
$(1-y(t))e^{ix(t)}$, hence lies in $A$. If $a\in A$, by Fact~\ref{fact:normal} $a$ is a point asymptotic value of $h$ at some $e^{ip}$, along a
path $w(t)\to e^{ip}$; writing $w(t)=r(t)e^{i\vartheta(t)}$ with $\vartheta$ continuous (for $t$ near $1$), $\vartheta(t)\to p+2\pi k$, and $\rho(a)$
is a point asymptotic value of $s$ at $p+2\pi k$ along $\langle\vartheta(t),1-r(t)\rangle$. Thus both sets equal $\rho(A)$, which is not Borel.
\end{proof}

\begin{remark}
\cite{PAMS} notes that the lemma can also be proved without meromorphic functions, from \cite[Theorem 113]{SierpGT}, a theorem of
Bagemihl--McMillan, and the methods of \cite{TAMSb}, but omits the ``messy construction''.
\end{remark}

\begin{theorem}[{\cite[Theorem]{PAMS}}]\label{thm:cube}
\lean{TJK.Cube.exists_not_measurableSet_curvConvSet}
\uses{lem:cube-lemma, thm:construct-C, lem:subarc, cor:path-arc, def:asymptotic}
There is a bounded continuous $f:K^\circ\to\C$ on the open cube $K^\circ=(-1,1)\times(0,1)\times(-1,1)$ whose set of curvilinear convergence
is not a Borel set. (So McMillan's theorem does not extend to three dimensions.)
\end{theorem}
\begin{proof}
Let $s,E$ be as in Lemma~\ref{lem:cube-lemma} and $g(x,y)=s(x/y,y)$ on $Q$.

\emph{$E$ is the set of real point asymptotic values of $g$ at $\langle0,0\rangle$ in $(-1,1)$.} If $g\to b$ along a path
$\langle x(t),y(t)\rangle\to\langle0,0\rangle$ in $Q^\circ$, then $s\to b$ along $\langle x/y,y\rangle$, whose second coordinate tends to $0$; so $b\in E$. If
$s\to b$ along $\langle u(t),v(t)\rangle\to\langle p,0\rangle$, then $g\to b$ along $\langle u(t)v(t),v(t)\rangle\to\langle0,0\rangle$, which lies in $Q^\circ$ for $t$ near $1$.

\emph{Construction.} Let $d_0(t)$ be the supremum of the $\delta\in(0,1]$ such that $z,z'\in Q$, $\operatorname{Im}z,\operatorname{Im}z'\ge t$,
$\abs{z-z'}<\delta$ imply $\abs{g(z)-g(z')}\le t$; let $d(t)=\min\{\frac14d_0(t/2),\frac14t\}$, and $k$ as in Theorem~\ref{thm:construct-C}
(so $\abs k\le\sqrt2\le2$). Here $d_0(t)>0$ because $g$ is uniformly continuous on the compact set $[-1,1]\times[t,1]$,
and $d_0$ is weakly increasing (raising $t$ shrinks the region and loosens the tolerance); so $d$ is positive and weakly
increasing, as Theorem~\ref{thm:construct-C} requires. Put $f(x,y,z)=(g(x,y)-z)\,k(x,y)$. For any three points $P=\langle x,y,z\rangle$, $P'$, $P''$ of $K^\circ$, writing
$g'=g(x',y')$ etc.,
\begin{equation}\label{eq:cube1}
\abs{f(P')-f(P'')}\ge\abs{g-z}\abs{k'-k''}-2\abs{g'-g}-2\abs{g-g''}-2\abs{z-z'}-2\abs{z''-z},
\end{equation}
because $f(P')-f(P'')=(g-z)(k'-k'')+\bigl((g'-z')-(g-z)\bigr)k'-\bigl((g''-z'')-(g-z)\bigr)k''$.

Let $L=\{\langle0,0,z\rangle:-1<z<1\}$ and $F=\CC(f)$. We show $F\cap L=\{\langle0,0,b\rangle:b\in E\}$; then $F$ is not Borel, since
otherwise $F\cap L$ would be.

If $b\in E$, take an arc $\gamma$ at $\langle0,0\rangle$ in $Q^\circ$ along which $g\to b$; along $\gamma'=\{\langle x,y,b\rangle:\langle x,y\rangle\in\gamma\}$,
$g-z\to0$, so $f\to0$, and $\langle0,0,b\rangle\in F$.

Conversely let $f$ converge along an arc $\gamma'$ at $\langle0,0,b\rangle$ in $K^\circ$, and let $\gamma$ be its projection to the $\langle x,y\rangle$-plane,
a path ending at $\langle0,0\rangle$ and otherwise in $Q^\circ$. We claim $g(x,y)-z\to0$ along $\gamma'$; then $g\to b$ along $\gamma$, so $b$ is a
point asymptotic value of $g$ at $\langle0,0\rangle$ and $b\in E$. If not, there are $\eps>0$ and points $\langle x_n,y_n,z_n\rangle\in\gamma'$
converging to $\langle0,0,b\rangle$ with $\abs{g(x_n,y_n)-z_n}\ge\eps$. Since the only point of $\gamma'$ with $y=0$ is its endpoint, there is
$\delta>0$ such that points of $\gamma'$ with $y\le\delta$ have $z$-coordinates within $\eps/8$ of each other. Fix $n$ large, with
$y_n<\min\{3\eps/32,3\delta/4,3/4\}$ and with the projection of the other endpoint of $\gamma'$ farther than $y_n$ from
$\langle x_n,y_n\rangle$, and put $a=4y_n/3$. By Lemma~\ref{lem:subarc}(b) applied to (a parametrization of) $\gamma$, there is a
connected $\gamma^*\subseteq\gamma\cap\overline{S(d(a),\langle x_n,y_n\rangle)}$ containing $\langle x_n,y_n\rangle$ and a point at distance $d(a)$ from it.
As $d(a)\le a/4$, $\gamma^*\subseteq[-1,1]\times[a/2,a]$, and Theorem~\ref{thm:construct-C} gives points $\langle x_n',y_n'\rangle$,
$\langle x_n'',y_n''\rangle$ of $\gamma^*$ with $\abs{k'-k''}\ge2$; lift them to points of $\gamma'$ with third coordinates $z_n',z_n''$. These have
$y\le a<\delta$, so $\abs{z_n'-z_n},\abs{z_n''-z_n}<\eps/8$; and $\abs{g(x_n,y_n)-g'},\abs{g(x_n,y_n)-g''}\le a/2<\eps/16$ by the choice of
$d$. By \eqref{eq:cube1}, $\abs{f(P_n')-f(P_n'')}\ge2\eps-\frac\eps4-\frac\eps2>\eps$. Since $P_n',P_n''\to\langle0,0,b\rangle$ along $\gamma'$, $f$ does
not converge along $\gamma'$: a contradiction.
\end{proof}

\begin{remark}
\cite{PAMS} leaves open whether a continuous \emph{real}-valued function in a cube can have a non-Borel set of curvilinear convergence.
\end{remark}

% ===========================================================================
\appendix
\section{The dissertation's special families}\label{app:variants}
% ===========================================================================
The dissertation uses a more restrictive notion of special family than \cite{TAMSb}. Its results are proved here for completeness;
they are not used elsewhere. A \emph{complementary interval} of a closed $F\subseteq\R$ is a component of $\R\setminus F$.

\begin{definition}[{\cite[\S4]{Diss}}]\label{def:th-special}
\lean{TJK.Variant.IsThesisSpecial}
A \emph{thesis-special family} is a nonempty family $\mathcal F$ of bounded closed subsets of $\R$ with an enumeration $F_1,F_2,\dots$ (every
member occurring) such that for $m>n$, either $F_m$ lies in a complementary interval of $F_n$ or $F_m$ splits with $F_n$.
\end{definition}

\begin{lemma}\label{lem:th-special-implies}
\lean{TJK.Variant.IsThesisSpecial.pairwise}
\uses{def:th-special, def:special-family, lem:special-pair-exists, lem:special-pair-shrink}
Any two members of a thesis-special family are equal, disjoint, or split. Hence Lemmas~\ref{lem:special-pair-exists}
and~\ref{lem:special-pair-shrink} apply to thesis-special families of subsets of $X$ (\cite[Lemmas 13, 14]{Diss}).
\end{lemma}
\begin{proof}
A set lying in a complementary interval of $F_n$ is disjoint from $F_n$.
\end{proof}

\begin{lemma}[{\cite[Lemma 9]{Diss}}]\label{lem:th-lemma9}
\lean{TJK.Variant.exists_seq_split_or_disjoint}
\uses{lem:special-from-fsigma}
Every $\Fs$ set $E\subseteq\R$ is $\bigcup_nE_n$ for bounded closed $E_n$ any two of which are disjoint or split.
\end{lemma}
\begin{proof}
As in Lemma~\ref{lem:special-from-fsigma}, with $\R$ in place of $X$; enumerate the resulting countable family, padding
with $\emptyset$ if it is finite (repeating a member with two points would violate the condition).
\end{proof}

\begin{lemma}[{\cite[Lemma 10]{Diss}}]\label{lem:th-lemma10}
\lean{TJK.Variant.exists_thesisSpecial_of_isFσ}
\uses{lem:th-lemma9, def:th-special}
Every $\Fs$ set $E\subseteq\R$ is the union of a thesis-special family.
\end{lemma}
\begin{proof}
Take $E_n$ from Lemma~\ref{lem:th-lemma9}. We construct $1=n_1<n_2<\cdots$ and bounded closed sets $F_1,F_2,\dots$ such that for every $s$:
(a) if $n_s\ge r>t\ge1$, then $F_r$ lies in a complementary interval of $F_t$ or splits with $F_t$; (b) each $F_i$, $i\le n_s$, lies in some
$E_j$, $j\le s$; (c) $\bigcup_{i\le n_s}F_i=\bigcup_{j\le s}E_j$. Start with $F_1=E_1$. Given stage $s$, let $\mathcal J$ be the set of
complementary intervals of the bounded closed set $\bigcup_{j\le s}E_j$.

\emph{$E_{s+1}$ meets only finitely many $I\in\mathcal J$.} Otherwise pick $x_m\in I_m\cap E_{s+1}$ for infinitely many distinct $I_m$; a
subsequence is strictly monotone, say increasing, with limit $x$. Between consecutive terms lies an endpoint $a_k$ of some $I_m$, which
belongs to $\bigcup_{j\le s}E_j$; infinitely many $a_k$ lie in one $E_u$, so $x\in E_u\cap E_{s+1}$. Then $E_u$ and $E_{s+1}$ split at $x$; but
both contain points less than $x$, which is impossible.

Let the new sets $F_{n_s+1},\dots,F_{n_{s+1}}$ be $\emptyset$ and the sets $\overline I\cap E_{s+1}$ for the finitely many $I\in\mathcal J$ meeting
$E_{s+1}$. (We use the closure $\overline I$, so that the new sets are closed; with $I$ itself, as written in \cite{Diss}, they need not be.)
(b) and (c) are clear. For (a) with $r$ new: if $t$ is old, $F_t\subseteq E_v$ with $v\le s$. If $E_v\cap E_{s+1}=\emptyset$, let $J$ be the
smallest closed interval containing $F_r=\overline I\cap E_{s+1}$; $J^*\subseteq I$ misses $E_v$, and the endpoints of $J$ lie in $E_{s+1}$, so
$J\cap E_v=\emptyset$ and $F_r$ lies in a complementary interval of $F_t$. If $E_v$ and $E_{s+1}$ split, so do $F_t$ and $F_r$. If $t$ is new too,
$F_r$ and $F_t$ come from distinct $I_1,I_2\in\mathcal J$, and $\overline{I_1},\overline{I_2}$ split. Finally
$\mathcal F=\{F_k\}$ is thesis-special with $\bigcup\mathcal F=E$.
\end{proof}

\begin{lemma}[{\cite[Lemmas 11, 12]{Diss}}]\label{lem:th-lemma11}
\lean{TJK.Variant.IsThesisSpecial.inf, TJK.Variant.sUnion_inf}
\uses{def:th-special}
If $\mathcal F$ and $\mathcal G$ are thesis-special, so is $\mathcal F\wedge\mathcal G=\{F\cap G:F\in\mathcal F,G\in\mathcal G\}$, and
$\bigcup(\mathcal F\wedge\mathcal G)=\bigcup\mathcal F\cap\bigcup\mathcal G$.
\end{lemma}
\begin{proof}
Enumerate $\N\times\N$ as $(a(k),b(k))$ by increasing $a+b$ and, for equal sums, increasing $b$; then $k<l$ implies $a(k)<a(l)$ or
$b(k)<b(l)$. Put $H_k=F_{a(k)}\cap G_{b(k)}$. If $l>k$ and $a(k)<a(l)$, then $F_{a(l)}$ lies in a complementary interval of $F_{a(k)}$, which
lies in a complementary interval of $H_k\subseteq F_{a(k)}$, or $F_{a(l)}$ splits with $F_{a(k)}$; in either case $H_l$ relates to $H_k$ in the
same way. The case $b(k)<b(l)$ is symmetric.
\end{proof}

% ===========================================================================
\section{Concordance with the sources}\label{app:concordance}
% ===========================================================================
\begin{longtable}{@{}p{0.3\textwidth}p{0.65\textwidth}@{}}
\toprule
Source item & Here \\ \midrule
\multicolumn{2}{@{}l}{\cite{MathMag}, with \cite{Sut}}\\
Sutcliffe Thm 2 & Prop.~\ref{prop:sut-two-composite}, Prop.~\ref{prop:sut-explicit}, Thm.~\ref{thm:sut-sutcliffe2}\\
Sutcliffe Thm 3 & Thm.~\ref{thm:sut-sutcliffe3}\\
Lemma 1, Lemma 2 & Lemma~\ref{lem:sut-lemma1}, Lemma~\ref{lem:sut-lemma2}\\
Theorem, Corollary & Thm.~\ref{thm:sut-main}, Cor.~\ref{cor:sut-main}\\ \midrule
\multicolumn{2}{@{}l}{\cite{JMM} (from the review)}\\
Theorem 1 (homeomorphisms) & Thm.~\ref{thm:homeo} (new proof)\\
Theorems 2, 3 & Cor.~\ref{cor:bf-honorary}\\
Theorems 4, 5 & Thm.~\ref{thm:baire-plus-one}, Cor.~\ref{cor:borel}\\
$f=0$ a.e.\ with given boundary function & Cor.~\ref{cor:ae-zero}\\
remarks on three dimensions & not a theorem; see Section~\ref{sec:cube}\\ \midrule
\multicolumn{2}{@{}l}{\cite{MMJ}}\\
Lemma 1 & Lemma~\ref{lem:intervals-open}\\
Lemmas 2, 3 & Lemmas~\ref{lem:mm-lemma2}, \ref{lem:mm-lemma3}\\
Theorem 1 (McMillan) & Thm.~\ref{thm:mcmillan}\\
Theorem 2, Corollaries 1--4 & Thm.~\ref{thm:mm-thm2}, Cor.~\ref{cor:mm-corollaries}\\
remark ($E$ of type $\Gd$) & Prop.~\ref{prop:gdelta-remark}\\ \midrule
\multicolumn{2}{@{}l}{\cite{Diss}}\\
Lemmas 1--6 & Lemma~\ref{lem:reduction}, \ref{lem:approx-char}, \ref{lem:approx-finite}, \ref{lem:uniform-approx}, \ref{lem:approx-real}, Thm.~\ref{thm:baire1-char}\\
Theorems 1, 2 & Thm.~\ref{thm:sphere-baire1}, Thm.~\ref{thm:honorary-char}\\
Lemmas 7, 8; Theorem 3 & Lemmas~\ref{lem:th-lemma7}, \ref{lem:th-lemma8}; Thm.~\ref{thm:mcmillan-H}\\
Lemmas 9--12 & Lemmas~\ref{lem:th-lemma9}, \ref{lem:th-lemma10}, \ref{lem:th-lemma11}\\
facts (6)--(17); Lemmas 13, 14 & Lemma~\ref{lem:region-geometry}; Lemmas~\ref{lem:special-pair-exists}, \ref{lem:special-pair-shrink}, \ref{lem:th-special-implies}\\
Theorem 4 & Thm.~\ref{thm:cc-realize} (and Thm.~\ref{thm:construct-A})\\
Lemma 15; Theorem 5 and Corollary & Lemma~\ref{lem:th-lemma15}; Thm.~\ref{thm:bf-fsigma}, Cor.~\ref{cor:bf-honorary}\\
Lemmas 16, 17; Theorems 6, 7 & Lemma~\ref{lem:closed-graph}, \ref{lem:monotone-inverse}; Thm.~\ref{thm:construct-B}, \ref{thm:construct-H}\\
Lemmas 18--22 & Lemmas~\ref{lem:equiv}, \ref{lem:intervals-open}, \ref{lem:monotone-extend}, \ref{lem:monotone-compose}, \ref{lem:baire-criterion}\\
Theorem 8 and Corollary & Thm.~\ref{thm:baire-plus-one}, Cor.~\ref{cor:borel}\\
Lemmas 23, 24; Theorem 9; example & Lemmas~\ref{lem:th-lemma23}, \ref{lem:th-lemma24}; Thm.~\ref{thm:th-thm9}; Prop.~\ref{prop:th-example}\\
Lemma 25; Theorem 10, Corollaries & Lemma~\ref{lem:th-lemma25}; Thm.~\ref{thm:th-thm10}; Cor.~\ref{cor:ae-zero}, \ref{cor:nonmeasurable}\\
unsolved problems 1--5 & open questions; not formalized\\ \midrule
\multicolumn{2}{@{}l}{\cite{TAMSa}}\\
Lemma & Lemma~\ref{lem:ha-lemma}\\
example (answer to \cite[Problem 5]{BP}) & Lemmas~\ref{lem:ha-sequences}, \ref{lem:ha-estimates}, Thm.~\ref{thm:ha-example}\\
Theorem, Corollary, analytic case & Thm.~\ref{thm:ha-theorem}, Cor.~\ref{cor:ha-corollary}, Cor.~\ref{cor:ha-analytic}\\ \midrule
\multicolumn{2}{@{}l}{\cite{TAMSb}}\\
(A) with (I)--(XXII) & Thm.~\ref{thm:construct-A}; Def.~\ref{def:special-family}, Lemmas~\ref{lem:special-from-fsigma}, \ref{lem:special-sequence}; Def.~\ref{def:regions}, Lemma~\ref{lem:region-geometry}; Def.~\ref{def:special-pair}, Lemmas~\ref{lem:special-pair-exists}, \ref{lem:special-pair-shrink}\\
(B) with (I)--(IV) & Thm.~\ref{thm:construct-B}; Lemmas~\ref{lem:monotone-inverse}, \ref{lem:closed-graph}; Thm.~\ref{thm:baire1-char}\\
(C), (D), (E), (F) & Thms.~\ref{thm:construct-C}, \ref{thm:construct-D}, \ref{thm:construct-E}, \ref{thm:construct-F}; Def.~\ref{def:fold}, Lemma~\ref{lem:fold}\\
Appendix (a), (b), (c) & Thms.~\ref{thm:sphere-baire1}, \ref{thm:honorary-char}, \ref{thm:sphere-extension}\\
Problems 1, 2 & open questions; not formalized\\ \midrule
\multicolumn{2}{@{}l}{\cite{PAMS}}\\
Lemma, Theorem & Lemma~\ref{lem:cube-lemma}, Thm.~\ref{thm:cube}\\
Problem & open question; not formalized\\
\bottomrule
\end{longtable}


\begin{thebibliography}{99}
\bibitem{JMM} T.~J. Kaczynski, \emph{Boundary functions for functions defined in a disk},
  J. Math. Mech. \textbf{14} (1965), 589--612. (Known here through MR0176080.)
\bibitem{MMJ} T.~J. Kaczynski, \emph{On a boundary property of continuous functions},
  Michigan Math. J. \textbf{13} (1966), 313--320.
\bibitem{Diss} T.~J. Kaczynski, \emph{Boundary functions}, Doctoral dissertation, University of
  Michigan, 1967.
\bibitem{MathMag} T.~J. Kaczynski, \emph{Note on a problem of Alan Sutcliffe}, Math. Mag.
  \textbf{41} (1968), 84--86.
\bibitem{TAMSa} T.~J. Kaczynski, \emph{Boundary functions for bounded harmonic functions},
  Trans. Amer. Math. Soc. \textbf{137} (1969), 203--209.
\bibitem{TAMSb} T.~J. Kaczynski, \emph{Boundary functions and sets of curvilinear convergence for
  continuous functions}, Trans. Amer. Math. Soc. \textbf{141} (1969), 107--125.
\bibitem{PAMS} T.~J. Kaczynski, \emph{The set of curvilinear convergence of a continuous function
  defined in the interior of a cube}, Proc. Amer. Math. Soc. \textbf{23} (1969), 323--327.
\bibitem{Sut} A.~Sutcliffe, \emph{Integers that are multiplied when their digits are reversed},
  Math. Mag. \textbf{39} (1966), 282--287.
\bibitem{BP} F.~Bagemihl and G.~Piranian, \emph{Boundary functions for functions defined in a disk},
  Michigan Math. J. \textbf{8} (1961), 201--207.
\bibitem{Bag} F.~Bagemihl, \emph{Curvilinear cluster sets of arbitrary functions}, Proc. Nat. Acad.
  Sci. U.S.A. \textbf{41} (1955), 379--382.
\bibitem{BS} F.~Bagemihl and W.~Seidel, \emph{Koebe arcs and Fatou points of normal functions},
  Comment. Math. Helv. \textbf{36} (1961), 9--18.
\bibitem{Banach} S.~Banach, \emph{\"Uber analytisch darstellbare Operationen in abstrakten
  R\"aumen}, Fund. Math. \textbf{17} (1931), 283--295.
\bibitem{Hausdorff} F.~Hausdorff, \emph{\"Uber halbstetige Funktionen und deren
  Verallgemeinerung}, Math. Z. \textbf{5} (1919), 292--309.
\bibitem{Hoffman} K.~Hoffman, \emph{Banach spaces of analytic functions}, Prentice-Hall, 1962.
\bibitem{Kierst} S.~Kierst, \emph{Sur l'ensemble des valeurs asymptotiques d'une fonction
  m\'eromorphe dans le cercle-unit\'e}, Fund. Math. \textbf{27} (1936), 226--233.
\bibitem{Kuratowski} K.~Kuratowski, \emph{Topologie I}, 3rd ed., Warsaw, 1952.
\bibitem{LV} O.~Lehto and K.~I. Virtanen, \emph{Boundary behaviour and normal meromorphic
  functions}, Acta Math. \textbf{97} (1957), 47--65.
\bibitem{McM} J.~E. McMillan, \emph{Boundary properties of functions continuous in a disc},
  Michigan Math. J. \textbf{13} (1966), 299--312.
\bibitem{Newman} M.~H.~A. Newman, \emph{Elements of the topology of plane sets of points}, 2nd ed.,
  Cambridge, 1964.
\bibitem{Sierpinski} W.~Sierpi\'nski, \emph{Sur une propri\'et\'e des ensembles ambigus}, Fund.
  Math. \textbf{6} (1924), 1--5.
\bibitem{SierpGT} W.~Sierpi\'nski, \emph{General topology}, Univ. of Toronto Press, 1952.
\bibitem{Tietze} H.~Tietze, \emph{\"Uber stetige Kurven, Jordansche Kurvenb\"ogen und geschlossene
  Jordansche Kurven}, Math. Z. \textbf{5} (1919), 284--291.
\end{thebibliography}

\end{document}
